\documentclass[12pt,titlepage,twoside,openright]{book}
\usepackage{color}
\usepackage{pstricks, pst-node}
\usepackage{graphics,graphicx,graphpap}
\usepackage{mathtools}
\usepackage{amsfonts}
\usepackage{amsmath}
\usepackage{amssymb}
\usepackage{amsthm}
\usepackage[ansinew]{inputenc}
\usepackage[spanish]{babel}
\usepackage{fancyhdr}
\pagestyle{empty}
\usepackage{tikz}
\usepackage{float}
\usetikzlibrary{calc,shapes,backgrounds}
\usetikzlibrary{positioning, fit}
\usepackage{pgfplots}
\usepackage{caption}
\usepackage{color}
\usepackage{algorithm}
\usepackage{algpseudocode}
\usepackage{subcaption}
\usepackage{comment}

\def\proof{{\noindent\bf Demostraci\'on:}\hskip 0.2truecm}
\def\BBox{\kern  -0.2cm\hbox{\vrule width 0.2cm height 0.2cm}}
\def\endproof{\hspace{.20in} \BBox\vspace{.20in}}


%%%%%%%%%%%%%%%%%%%%%%%%%%%%%%%%%%%%%%%%%%%%%%%%%%%%%%%%
%%% Plantilla elaborada por Alma Roc\'io Sagaceta Mej\'ia (2010) y modificada por Juli\'an Alberto Fres\'an Figueroa (2011-2020) %%%
%%%%%%%%%%%%%%%%%%%%%%%%%%%%%%%%%%%%%%%%%%%%%%%%%%%%%%%


%%%%%%%%%%%%%%%%%%%%%%%%%%%%%%%%%%%%%%%%%%%%%%%%%%%%%%%
%Aqui se pueden agregar palabras que Latex separa de forma incorrecta al haber saltos de linea. Al escribir fa-mi-lia se le esta indicando que puede separar la palabra solo en los guiones
%%%%%%%%%%%%%%%%%%%%%%%%%%%%%%%%%%%%%%%%%%%%%%%%%%%%%%%

\hyphenation{Colores fa-mi-lia caracte-rizarlos}



%%%%%%%%%%%%%%%%%%%%%%%%%%%%%%%%%%%%%%%%%%%%%%%%%%%%
%%%              UAM-Logo                   	 %%%
%%%         Por Ismael Vel\'azquez Ram\'irez         %%%
%%%%%%%%%%%%%%%%%%%%%%%%%%%%%%%%%%%%%%%%%%%%%%%%%%%%

\newcommand{\uamlogo}[3][2pt]{
    \psset{unit=#2,linewidth=#1 }
    \psline*[linearc=.25,linecolor=#3](2.8,2)(2,2)(1.8,0)(2.8,2)(3.8,0)(3.6,2)(2,2)(1.8,0)
    \psline*[linecolor=#3](0,0)(.8,0)(1.8,2)(1,2)(0,0)
    \psline*[linecolor=#3](4.8,0)(3.8,2)(4.6,2)(5.6,0)(4.8,0)
    \psline*[linearc=.25,linecolor=#3](3.8,0)(2.8,2)(3.6,2)(4.6,0)(3.8,0)
    \psline*[linearc=.25,linecolor=#3](4.6,0)(3.8,0)(2.8,2)(3.6,2)(4.6,0)
    \rput{180}(5.6,2){%
    \psline*[linearc=.25,linecolor=white](2.8,2)(2,2)(1.8,0)(2.8,2)(3.8,0)(3.6,2)(2,2)(1.8,0)
    \psline*[linearc=.25,linecolor=white](1,0)(1.8,0)(2.8,2)(2,2)(1,0)
    \psline*[linearc=.25,linecolor=white](1.8,0)(1,0)(2,2)(2.8,2)(1.8,0)
    \psline*[linearc=.25,linecolor=white](3.8,0)(2.8,2)(3.6,2)(4.6,0)(3.8,0)
    \psline*[linearc=.25,linecolor=white](4.6,0)(3.8,0)(2.8,2)(3.6,2)(4.6,0)
    \psline[linearc=.25,linecolor=#3](1,0)(2,2)(3.6,2)(4.6,0)
    \psline[linecolor=#3](1,0)(1.8,0)
    \psline[linearc=.25,linecolor=#3](4.6,0)(3.8,0)
    \psline[linearc=.25,linecolor=#3](1.8,0)(2.8,2)(3.8,0)}
    \psline*[linearc=.25,linecolor=#3](1,0)(1.8,0)(2.8,2)(2,2)(1,0)
    \psline*[linearc=.25,linecolor=#3](1.8,0)(1,0)(2,2)(2.8,2)(1.8,0)}
    
    
%%%%%%%%%%%%%%%%%%%%%%%%%%%%%%%%%%%%%%%%%%%%%%%%%%%%%%%
% Aqui se definen los entornos en español.
%%%%%%%%%%%%%%%%%%%%%%%%%%%%%%%%%%%%%%%%%%%%%%%%%%%%%%%

\newtheorem{conjecture}{Conjetura}[section]
\newtheorem{lemma}{Lema}[section]
\newtheorem{theorem}{Teorema}[section]
\newtheorem{definition}{Definici\'on}[section]
\newtheorem{corollary}{Corolario}[section]
\newtheorem{proposition}{Proposici\'on}[section]
\newtheorem{remark}{Observaci\'on}[section]
\newtheorem{claim}{Afirmar}[section]
\newtheorem{nota}{Nota}[section]
\newtheorem{teo}{Teorema}
\newtheorem{defi}{Definici\'on}
\newtheorem{coro}{Corolario}

%Ejemplo de uso:
%\begin{teo}[Teorema de Completitud] 
%Sea...
%\end{teo}



\newcommand{\titulo}[1]{\def\eltitulo{#1}}
\newcommand{\carrera}[1]{\def\lacarrera{#1}}
\newcommand{\nombre}[1]{\def\elnombre{#1}}    %* Del alumno
\newcommand{\director}[1]{\def\eldirector{#1}}  %* De tesis
\newcommand{\fecha}[1]{\def\lafecha{#1}}


%%%%%%%%%%%%%%%%%%%%%%%%%%%%%%%%%%%%%%%%%%%%%%%%%%%%%%%
% Se deben llenar los siguientes datos
%Estos datos apareceran en la portada y los encabezados
%%%%%%%%%%%%%%%%%%%%%%%%%%%%%%%%%%%%%%%%%%%%%%%%%%%%%%%

\titulo{Din\'amica de la infecci\'on por VIH en los ganglios linf\'aticos: Un enfoque de teor\'ia de gr\'aficas}
\nombre{\uppercase{Lizzeth Ariadna S\'anchez Solis}}
\carrera{Licenciatura en Matem\'aticas Aplicadas}
\director{\uppercase{Nombre del asesor(es)}}
\fecha{Diciembre 2025}


\begin{document}

%%%%%%%%%%%%%%%%%%%%%%%%%%%%%%%%%%%%%%%%%%%%%%%%%%%%%%%
% Portada
%%%%%%%%%%%%%%%%%%%%%%%%%%%%%%%%%%%%%%%%%%%%%%%%%%%%%%%


\thispagestyle{empty}

\hskip-1.5cm
\begin{minipage}[c][5cm][s]{4cm}
  \begin{center}
  \rput(-2,.5){\uamlogo{.75}{black}}

  
    \hskip2pt \vrule width2pt height16cm\hskip1mm
    \vrule width1pt height16cm\\[10pt]

  \end{center}
\end{minipage}\quad
%% Barra derecha - Tí­tulos
\begin{minipage}[c][9.5cm][s]{10cm}
  \begin{center}
    % Barra superior
    {\Large \scshape Universidad Aut\'onoma Metropolitana}
    \vspace{.3cm}
    \hrule height2pt
    \vspace{.1cm}
    \hrule height1pt
    \vspace{.3cm}
    % \scshape
    { UNIDAD CUAJIMALPA}


%%%%%%%%%%%%%%%%%%%%%%%%%%%%%%%%%%%%%%%%%%%%%%%%%%%%%%%
% Poner aqu\'i el t\'itulo del trabajo
%%%%%%%%%%%%%%%%%%%%%%%%%%%%%%%%%%%%%%%%%%%%%%%%%%%%%%%
    \vspace{2cm}
    
    {\Large \textsc{Din\'amica de la infecci\'on por VIH en los ganglios linf\'aticos: Un enfoque de teor\'ia de gr\'aficas}} 

    \vspace{2cm}







    \makebox[8cm][c]{\Large Proyecto Terminal}\\[8pt]
    
   
    \makebox[6cm]{QUE PRESENTA:}\\[3pt]
    Lizzeth Ariadna S\'anchez Solis\\
    \vspace{.5cm}
  {\textsc {\large Licenciatura en Matem\'aticas Aplicadas}}\\[3pt]
  \makebox[10cm]{Departamento de Matem\'aticas Aplicadas e Ingenier\'{i}a}\\[13pt]
\makebox[10cm]{Divisi\'on de Ciencias Naturales e Ingenier\'ia}\\[13pt]


    \vspace{2cm}

    { Asesor y Responsable de la tesis:\\ Juli\'an Alberto Fres\'an Figueroa\\}

    \vspace{2cm}
\begin{flushright}
\lafecha 
\end{flushright}

  \end{center}
\end{minipage}


\frontmatter{}



\pagestyle{plain}
	

\tableofcontents
\newpage 
\listoffigures	




\mainmatter{}
\pagestyle{fancy}
\renewcommand{\chaptermark}[1]{%
\markboth{\chaptername
\ \thechapter.\ #1}{}}
\fancyhead[RO,LE]{\bfseries \thepage}
\fancyfoot{}
\setlength{\parindent}{0pt}
\setlength{\parskip}{1.5ex}


%%%%%%%%%%%%%%%%%%%%%%%%%%%%%%%%%%%%%%%%%%%%%%%%%%%%%%%
% Aqu\'i se empieza a colocar la informaci\'on del trabajo
% El documento deber\'a contener al menos los  cap\'itulos que aparecen a continuaci\'on.
% Cada capitulo puede tener secciones y subsecciones
% Las im\'agenes deben ser incluidas en formato .eps
%%%%%%%%%%%%%%%%%%%%%%%%%%%%%%%%%%%%%%%%%%%%%%%%%%%%%%%

\chapter{Resumen}

%Uno o dos p\'arrafos de que se hizo en el proyecto terminal

En este trabajo se estudia la din\'amica de propagaci\'on del Virus de Inmunodeficiencia Humana (VIH) en los ganglios linf\'aticos mediante un enfoque basado en teor\'ia de gr\'aficas. A partir del modelo propuesto por Simon Mukwembi \cite{Mukwembi}, se representa cada sitio susceptible a infecci\'on como un v\'ertice de una gr\'afica y se describe la evoluci\'on temporal de la infecci\'on mediante reglas dependientes de un par\'ametro de reemplazo $R$. El objetivo central consiste en analizar, para una gr\'afica $G$, el valor m\'inimo del par\'ametro $R$ que garantiza la eliminaci\'on total del virus para cualquier configuraci\'on inicia, dicho valor se denota como $VIH(G)$. Este enfoque permite relacionar la estructura de la gr\'afica con la persistencia o desaparici\'on de la infecci\'on, mostrando que la geometr\'ia de la red celular influye de manera decisiva en los procesos de re-infecci\'on y curaci\'on.

Se desarroll\'o un estudio te\'orico y computacional del comportamiento del n\'umero VIH en diversas familias de gr\'aficas. Para cada una de estas familias se obtuvieron tablas experimentales, se identificaron patrones estructurales y se demostraron teoremas que determinan el valor exacto o acotado de $VIH(G)$. Los resultados muestran que la conectividad, el grado de los v\'ertices y la presencia de subgrafos latentes influyen directamente en la persistencia de la infecci\'on. Este estudio integra modelaci\'on matem\'atica, teor\'ia de gr\'aficas y experimentaci\'on computacional, ofreciendo una perspectiva novedosa sobre la propagaci\'on del VIH en redes biol\'ogicas.


\chapter{Introducci\'on}
%una o dos p\'aginas, escrita para ser comprensible por cualquier persona

El Virus de Inmunodeficiencia Humana (VIH) es uno de los agentes infecciosos m\'as estudiados en las \'ultimas d\'ecadas debido a su profundo impacto en la salud p\'ublica. Este virus ataca directamente al sistema inmunitario humano, particularmente a las c\'elulas CD4, encargadas de coordinar la respuesta del organismo frente a infecciones. Aunque los tratamientos actuales permiten controlar el virus y mejorar la calidad de vida de las personas infectadas, el VIH contin\'ua siendo un desaf\'io cient\'ifico y m\'edico, especialmente en lo que respecta a comprender a detalle su propagaci\'on dentro del cuerpo humano.

Diversos estudios han demostrado que la mayor parte del proceso de infecci\'on no ocurre en la sangre, sino en los ganglios linf\'aticos. Su estructura interna puede interpretarse como una red compleja de sitios estrechamente conectados entre s\'i, donde las c\'elulas sanas, infectadas y muertas interact\'uan de manera din\'amica. Debido a esta complejidad geom\'etrica y biol\'ogica, resulta natural recurrir a herramientas matem\'aticas que permitan modelar y analizar el comportamiento del virus en un entorno idealizado pero representativo. 

Dentro de este contexto, Simon Mukwembi propuso un modelo basado en teor\'ia de gr\'aficas para describir la propagaci\'on del VIH en los ganglios linf\'aticos \cite{Mukwembi}. En dicho modelo, cada sitio del ganglio se representa como un v\'ertice, y dos v\'ertices est\'an conectados si sus sitios correspondientes son vecinos en el ganglio. Cada v\'ertice puede encontrarse en tres estados: saludable, infectado o muerto, y la evoluci\'on del sistema se describe paso a paso mediante reglas que imitan el comportamiento biol\'ogico del virus. Una c\'elula sana puede infectarse si tiene un vecino infectado; una c\'elula infectada muere en el siguiente instante; y una c\'elula muerta puede ser reemplazada por una c\'elula saludable o infectada dependiendo del n\'umero de vecinos infectados. 

El par\'ametro central del modelo es $R$, el n\'umero m\'inimo de vecinos infectados necesarios para que una c\'elula muerta se regenere como infectada. Cuando $R$ es peque\~no, la infecci\'on puede mantenerse activa durante largo tiempo; cuando $R$ es grande, el sistema tiende a recuperarse. Un aporte fundamental de Mukwembi fue demostrar que, si $R$ es igual al n\'umero total de v\'ertices menos uno, el sistema siempre termina cur\'andose sin importar c\'omo haya comenzado la infecci\'on. Esto sugiere que ciertos patrones de reemplazo celular podr\'ian detener completamente la propagaci\'on viral. A partir de este resultado surge una pregunta natural: ?`cu\'al es el valor m\'as peque\~no de $R$ que garantiza la curaci\'on del sistema para cualquier configuraci\'on inicial? A este valor se le denomina el n\'umero de VIH de la gr\'afica y se denota por $VIH(G)$.

El estudio del n\'umero VIH permite explorar c\'omo la estructura de una red influye en la persistencia o eliminaci\'on del virus. Algunas redes, presentan patrones de propagaci\'on simples, mientras que otras poseen configuraciones m\'as complejas que pueden facilitar la presencia de ciclos de re-infecci\'on. En este proyecto se analiza detalladamente este comportamiento para distintas familias de gr\'aficas, combinando herramientas te\'oricas con simulaciones computacionales implementadas en SageMath. El objetivo es identificar valores exactos o acotados para $VIH(G)$, descubrir patrones generales y formular conjeturas que describan el comportamiento del modelo en estructuras m\'as amplias.

Este trabajo busca contribuir al entendimiento de c\'omo la geometr\'ia de los ganglios linf\'aticos puede influir en la din\'amica del VIH. Aunque el modelo es una simplificaci\'on de la compleja realidad biol\'ogica, ofrece una perspectiva matem\'atica clara que permite aislar ciertos mecanismos fundamentales de propagaci\'on. Los resultados obtenidos muestran que la topolog\'ia de la red juega un papel crucial en la persistencia de la infecci\'on, y proporcionan una base para estudios futuros que integren modelos m\'as detallados o incorporen informaci\'on adicional sobre el comportamiento del virus.


\chapter{Conocimientos preliminares}
%En este cap\'itulo deber\'an incluir todos los conocimientos necesarios para poder entender los temas que se trataron en el proyecto. 
Antes de abordar el desarrollo del modelo propuesto y los resultados experimentales, es necesario establecer ciertos conceptos te\'oricos y t\'ecnicos que sirven como base para la comprensi\'on integral del trabajo. En primer lugar, se presenta una descripci\'on general del Virus de Inmunodeficiencia Humana VIH. Posteriormente, se introducen los conceptos fundamentales de la teor\'ia de gr\'aficas, incluyendo las principales familias analizadas, as\'i como las nociones de grado, distancia, excentricidad y conectividad, que resultan esenciales para interpretar los resultados obtenidos. Finalmente, se exponen los conocimientos b\'asicos de programaci\'on utilizados en la implementaci\'on computacional del modelo, con el prop\'osito de explicar las herramientas empleadas y la l\'ogica detr\'as de los algoritmos desarrollados.

\textbf{?`Qu\'e es el VIH?}

El Virus de Inmunodeficiencia Humana (VIH) es un retrovirus que ataca al sistema inmunitario humano, espec\'ificamente a las c\'elulas CD4 o linfocitos T cooperadores, que desempe\~nan un papel esencial en la coordinaci\'on de la respuesta inmunol\'ogica del organismo. Una vez que el virus invade estas c\'elulas, utiliza su maquinaria interna para replicarse y producir copias de s\'i mismo.

El ciclo de infecci\'on del VIH comienza cuando el virus se adhiere a la superficie de una c\'elula CD4 mediante receptores espec\'ificos. Posteriormente, el ARN viral se convierte en ADN mediante una enzima denominada transcriptasa inversa, la cual permite integrar el material gen\'etico del virus en el ADN de la c\'elula hu\'esped. A partir de este momento, la c\'elula infectada comienza a producir nuevas part\'iculas virales, que eventualmente se liberan y pueden infectar otras c\'elulas.

El proceso continuo de infecci\'on y destrucci\'on de c\'elulas CD4 debilita progresivamente el sistema inmunitario. Si no se administra tratamiento, esta disminuci\'on en la cantidad de c\'elulas CD4 compromete la capacidad del organismo para combatir infecciones y enfermedades, lo que conduce al S\'indrome de Inmunodeficiencia Adquirida (SIDA).

Aunque actualmente no existe una cura definitiva, los tratamientos antirretrovirales permiten controlar la replicaci\'on del virus, mantener la funci\'on inmunitaria y reducir significativamente el riesgo de transmisi\'on. De esta manera, una persona que vive con VIH y sigue adecuadamente su tratamiento puede alcanzar una carga viral indetectable y llevar una vida saludable.

La mayor parte del proceso de infecci\'on no ocurre en la sangre, sino en los ganglios linf\'aticos, donde las c\'elulas CD4 se encuentran densamente agrupadas. Esta caracter\'istica ha motivado el desarrollo de diversos enfoques te\'oricos y experimentales para estudiar la din\'amica del virus en dichos entornos biol\'ogicos.

\textbf{Conceptos de teor\'ia de gr\'aficas}

%Aqui va definicion de grafica, grado, distancia, arbol, las familias 
En esta secci\'on se presentan las nociones fundamentales de la teor\'ia de gr\'aficas \cite{Harary} que ser\'an empleadas a lo largo del trabajo. Todas las gr\'aficas consideradas son finitas, simples y no dirigidas, salvo que se indique lo contrario.

Una \textit{gr\'afica} $G$ se define como un par ordenado $G = (V(G), E(G))$, donde $V(G)=\{v_1, v_2, \dots, v_n\}$ con $n \in \mathbb{N}$ es un conjunto finito y no vac\'io cuyos elementos se denominan v\'ertices, y $E(G) \subseteq \{ \{v_i,v_j\} : v_i,v_j \in V(G),\, v_i \neq v_j \}$ es el conjunto de aristas, las cuales representan las conexiones entre los v\'ertices. Dos v\'ertices $v_i$ y $v_j$ son adyacentes si $\{v_i,v_j\} \in E(G)$. En la Figura \ref{EjemG}, tenemos una gr\'afica $G$ conformada por 8 v\'ertices denotados como $V(G)=\{v_1,v_2,v_3,v_4,v_5,v_6,v_7,v_8\}$ y 7 aristas: $$E(G)=\{\{v_1,v_2\}, \{v_2,v_3\}, \{v_3,v_4\}, \{v_4,v_5\},\{v_4,v_6\},\{v_4,v_7\},\{v_7,v_8\}\}.$$ 
\begin{figure}[h]
\centering
\begin{tikzpicture}[every node/.style={circle, draw, fill=blue}]
% Nodos
\node[label=below:{$v_1$}] (v1) at (0,0) {};
\node[label=below:{$v_2$}] (v2) at (1.5,0) {};
\node[label=below:{$v_3$}] (v3) at (3,0) {};
\node[label=below:{$v_4$}] (v4) at (4.5,0) {};
\node[label=below:{$v_5$}] (v5) at (6,0) {};
\node[label=above:{$v_6$}] (v6) at (3,1.5) {};
\node[label=above:{$v_7$}] (v7) at (4.5,1.5) {};
\node[label=above:{$v_8$}] (v8) at (6,1.5) {};

\draw (v4)--(v5);
\draw (v4)--(v6);
\draw (v4)--(v7);
\draw (v8)--(v7);

% Aristas del camino rojo
\draw[thick,red] (v1)--(v2)--(v3)--(v4)--(v6);
\end{tikzpicture}
\caption{Ejemplo de una gr\'afica $G$.}
\label{EjemG}
\end{figure}
El \textit{grado} de un v\'ertice $v$, denotado por $\deg_G(v)$, es el n\'umero de aristas incidentes en \'el. El \textit{grado m\'inimo} de $G$, denotado por $\delta(G)$, se define como $\delta(G) = \min\{\deg_G(v) : v \in V(G)\}$, mientras que el \textit{grado m\'aximo}, denotado por $\Delta(G)$, se define como $\Delta(G) = \max\{\deg_G(v) : v \in V(G)\}$, por ejemplo: en la Figura \ref{EjemG}, tenemos que $\delta(G)=1$ y $\Delta(G)=4$. Si todos los v\'ertices de $G$ tienen el mismo grado $k$, se dice que la gr\'afica es \textit{k-regular}. Por ejemplo en la Figura \ref{fig:Completas}, la gr\'afica completa $K_4$ es $3$-regular, pues cada v\'ertice se conecta con los otros tres.

Un \textit{camino} en $G$ es una sucesi\'on finita de v\'ertices $v_1, v_2, \ldots, v_k$ tal que $\{v_{i-1}, v_i\} \in E(G)$ para todo $i = 1, \ldots, k$. La \textit{longitud} de dicho camino es el n\'umero de aristas que lo conforman, es decir, $k$. La \textit{distancia} entre dos v\'ertices $u,v \in V(G)$, denotada por $d(u,v)$, se define como la longitud del camino m\'as corto que los une, por ejemplo en la Figura \ref{EjemG} el camino rojo une $v_1$ y $v_6$ con longitud 4, que coincide con la distancia $d(v_1,v_6)=4$.

A partir de esta noci\'on se define la \textit{excentricidad} de un v\'ertice $v \in V(G)$ como $\varepsilon_G(v) = \max \{ d(v,u) \mid u,v \in V(G) \}$, es decir, la mayor distancia entre $v$ y cualquier otro v\'ertice de $G$. En la Figura \ref{EjemG}, podemos observar que el v\'ertice $v_1$ posee una excentricidad de $\varepsilon_G(v_1) = 5$. 

Una de las familias m\'as elementales es la de las trayectorias. La \textit{trayectoria} de orden $n$, denotada por $P_n$, es una gr\'afica cuyos v\'ertices pueden ordenarse como $v_1, v_2, \dots, v_n$ de modo que dos v\'ertices son adyacentes si y s\'olo si sus \'indices difieren en una unidad. Formalmente, el conjunto de v\'ertices y aristas se define como $$ V(P_n) = \{v_1, v_2, \dots, v_n\}, \quad E(P_n) = \{\{v_i, v_{i+1}\} : 1 \leq i \leq n-1\}.$$ Esta gr\'afica es conexa, contiene $n-1$ aristas y no posee ciclos en la Figura \ref{fig:Trayectorias} se pueden observar algunos ejemplos.

\begin{figure}[h]    
\begin{center}
    \begin{tikzpicture}
        % vertices P_3
        \node[draw, circle, fill=blue, label=left:{$P_3:$}] (1) at (0,0) {};
        \node[draw, circle, fill=blue, label=above:{$ $}] (2) at (1,0) {};
        \node[draw, circle, fill=blue, label=left:{$ $}] (3) at (2,0) {};

        % Aristas P_3
        \draw (1) -- (2);
        \draw (2) -- (3);

        % vertices P_4
        \node[draw, circle, fill=blue, label=left:{$P_4:$}] (4) at (0,-1) {};
        \node[draw, circle, fill=blue, label=above:{$ $}] (5) at (1,-1) {};
        \node[draw, circle, fill=blue, label=below:{$ $}] (6) at (2,-1) {};
        \node[draw, circle, fill=blue, label=below:{$ $}] (7) at (3,-1) {};

        % Aristas P_4
        \draw (4) -- (5);
        \draw (5) -- (6);
        \draw (6) -- (7);

        % vertices P_5
        \node[draw, circle, fill=blue, label=left:{$P_5:$}] (8) at (0,-2) {};
        \node[draw, circle, fill=blue, label=above:{$ $}] (9) at (1,-2) {};
        \node[draw, circle, fill=blue, label=below:{$ $}] (10) at (2,-2) {};
        \node[draw, circle, fill=blue, label=below:{$ $}] (11) at (3,-2) {};
        \node[draw, circle, fill=blue, label=below:{$ $}] (12) at (4,-2) {};

        % Aristas P_4
        \draw (8) -- (9);
        \draw (9) -- (10);
        \draw (10) -- (11);
        \draw (11) -- (12);
    \end{tikzpicture}   
\end{center}
\caption{Ejemplos de las gr\'aficas $P_3$, $P_4$ $P_5$.}
\label{fig:Trayectorias}
\end{figure}
\newpage
Otra familia fundamental es la de los ciclos. El \textit{ciclo} de orden $n$, denotado $C_n$, se obtiene a partir de la trayectoria $P_n$ al a\~nadir la arista que conecta los extremos de la secuencia de v\'ertices. De manera formal,$$ V(C_n) = \{v_1, v_2, \dots, v_n\}, \quad E(C_n) = \{\{v_i, v_{i+1}\} : 1 \leq i \leq n-1\} \cup \{\{v_n, v_1\}\}.$$ Cada v\'ertice en $C_n$ tiene grado $2$, y la gr\'afica es conexa y regular. En la Figura \ref{fig:Ciclos} se muestran algunos ejemplos de est\'a familia de gr\'aficas.

\begin{figure}[h]
    \centering
    \begin{tikzpicture}
        % vertices C_3 
        \node[draw, circle, fill=blue, label=above:{$ $}] (1) at (2,2) {};
        \node[draw, circle, fill=blue, label={[label distance=13pt]355:$C_3$}] (2) at (1,0) {};
        \node[draw, circle, fill=blue, label=left:{$ $}] (3) at (3,0) {};

        % Aristas C_3
        \draw (1) -- (2);
        \draw (2) -- (3);
        \draw (3) -- (1);

        % vertices C_4
        \node[draw, circle, fill=blue, label={[label distance=13pt]355:$C_4$}] (4) at (4,0) {};
        \node[draw, circle, fill=blue, label=above:{$ $}] (5) at (6,0) {};
        \node[draw, circle, fill=blue, label=below:{$ $}] (6) at (6,2) {};
        \node[draw, circle, fill=blue, label=below:{$ $}] (7) at (4,2) {};

        % Aristas P_4
        \draw (4) -- (5);
        \draw (5) -- (6);
        \draw (6) -- (7);
        \draw (4) -- (7);

        % vertices c_5
        \node[draw, circle, fill=blue, label=left:{$ $}] (8) at (7,1.5) {};
        \node[draw, circle, fill=blue, label={[label distance=13pt]355:$C_5$}] (9) at (7.5,0) {};
        \node[draw, circle, fill=blue, label=below:{$ $}] (10) at (9.5,0) {};
        \node[draw, circle, fill=blue, label=below:{$ $}] (11) at (10,1.5) {};
        \node[draw, circle, fill=blue, label=below:{$ $}] (12) at (8.5,2.5) {};

        % Aristas C_4
        \draw (8) -- (9);
        \draw (9) -- (10);
        \draw (10) -- (11);
        \draw (11) -- (12);
        \draw (8) -- (12);
    \end{tikzpicture}   
\caption{Ejemplos de las gr\'aficas $C_3$, $C_4$ $C_5$.}
\label{fig:Ciclos}
\end{figure}
La \textit{gr\'afica completa} de orden $n$, denotada $K_n$, es aquella en la que cada par de v\'ertices distintos est\'a unido por una arista. Esto se expresa formalmente como $$ V(K_n) = \{v_1, v_2, \dots, v_n\}, \quad E(K_n) = \{\{v_i, v_j\} : v_i, v_j \in V(K_n),\, i \neq j\}.$$ En $K_n$ todos los v\'ertices son adyacentes entre s\'i, por lo que el grado de cada uno es $n-1$. Esta gr\'afica es $(n-1)$-regular y contiene $\frac{n(n-1)}{2}$ aristas. A continuaci\'on en la Figura \ref{fig:Completas} se contemplas algunos ejemplos de $K_n$.
\newpage
\begin{figure}[h]  
    \centering
    \begin{tikzpicture}
        % vertices K_3
        \node[draw, circle, fill=blue, label=above:{$ $}] (1) at (2,2) {};
        \node[draw, circle, fill=blue, label={[label distance=13pt]355:$K_3$}] (2) at (1,0) {};
        \node[draw, circle, fill=blue, label=left:{$ $}] (3) at (3,0) {};

        % Aristas K_3
        \draw (1) -- (2);
        \draw (2) -- (3);
        \draw (3) -- (1);

        % vertices K_4
        \node[draw, circle, fill=blue, label={[label distance=13pt]355:$K_4$}] (4) at (4,0) {};
        \node[draw, circle, fill=blue, label=above:{$ $}] (5) at (6,0) {};
        \node[draw, circle, fill=blue, label=below:{$ $}] (6) at (6,2) {};
        \node[draw, circle, fill=blue, label=below:{$ $}] (7) at (4,2) {};

        % Aristas K_4
        \draw (4) -- (5);
        \draw (5) -- (6);
        \draw (6) -- (7);
        \draw (4) -- (7);
        \draw (4) -- (6);
        \draw (5) -- (7);

        % vertices K_5
        \node[draw, circle, fill=blue, label=left:{$ $}] (8) at (7,1.5) {};
        \node[draw, circle, fill=blue, label={[label distance=13pt]355:$K_5$}] (9) at (7.5,0) {};
        \node[draw, circle, fill=blue, label=below:{$ $}] (10) at (9.5,0) {};
        \node[draw, circle, fill=blue, label=below:{$ $}] (11) at (10,1.5) {};
        \node[draw, circle, fill=blue, label=below:{$ $}] (12) at (8.5,2.5) {};

        % Aristas K_4
        \draw (8) -- (9);
        \draw (8) -- (10);
        \draw (8) -- (11);
        \draw (8) -- (12);
        \draw (9) -- (10);
        \draw (9) -- (11);
        \draw (9) -- (12);
        \draw (10) -- (11);
        \draw (10) -- (12);
        \draw (11) -- (12);        
    \end{tikzpicture}   
\caption{Ejemplos de las gr\'aficas $K_3$, $K_4$ $K_5$.}
\label{fig:Completas}
\end{figure}
Una gr\'afica de inter\'es relacionada con $K_n$ es la \textit{estrella} de orden $n$, denotada por $S_n$. Se define como la gr\'afica con un v\'ertice central $v_0$ adyacente a los dem\'as v\'ertices, los cuales no est\'an conectados entre s\'i. En forma formal, $$ V(S_n) = \{v_0, v_1, v_2, \dots, v_{n-1}\}, \quad E(S_n) = \{\{v_0, v_i\} : 1 \leq i \leq n-1\}.$$ El v\'ertice $v_0$ tiene grado $n-1$, mientras que los dem\'as v\'ertices tienen grado $1$, esto se observa mejor en la Figura \ref{fig:Estrellas}.

\begin{figure}[h]
    \centering
    \begin{tikzpicture}
        % vertices S_3
        \node[draw, circle, fill=blue, label=right:{$ $}] (1) at (0,0) {};
        \node[draw, circle, fill=blue, label=above:{$ $}] (2) at (0,1) {};
        \node[draw, circle, fill=blue, label={[label distance=13pt]355:$S_3$}] (3) at (-1,-1) {};
        \node[draw, circle, fill=blue, label=below:{$ $}] (4) at (1,-1) {};

        % Aristas s_3
        \draw (1) -- (2);
        \draw (1) -- (3);
        \draw (1) -- (4);

        % vertices S_4
        \node[draw, circle, fill=blue, label=right:{$ $}] (4) at (3,0) {};
        \node[draw, circle, fill=blue, label=above:{$ $}] (5) at (4,1) {};
        \node[draw, circle, fill=blue, label=above:{$ $}] (6) at (2,1) {};
        \node[draw, circle, fill=blue, label={[label distance=13pt]355:$S_4$}] (7) at (2,-1) {};
        \node[draw, circle, fill=blue, label=below:{$ $}] (8) at (4,-1) {};

        % Aristas S_4
        \draw (4) -- (5);
        \draw (4) -- (6);
        \draw (4) -- (7);
        \draw (4) -- (8);

        % vertices S_5
        \node[draw, circle, fill=blue, label=right:{$ $}] (9) at (6,0) {};
        \node[draw, circle, fill=blue, label=above:{$ $}] (10) at (6.7,0.9) {};
        \node[draw, circle, fill=blue, label=above:{$ $}] (11) at (5.5,1) {};
        \node[draw, circle, fill=blue, label=below:{$ $}] (12) at (5,0) {};
        \node[draw, circle, fill=blue, label=below:{$S_5$}] (13) at (6,-1) {};
        \node[draw, circle, fill=blue, label=below:{$ $}] (14) at (7,-0.5) {};
        
        % Aristas S_5
        \draw (9) -- (10);
        \draw (9) -- (11);
        \draw (9) -- (12);
        \draw (9) -- (13);
        \draw (9) -- (14);
    \end{tikzpicture}   
\caption{Ejemplos de las gr\'aficas $S_3$, $S_4$ $S_5$.}
\label{fig:Estrellas}
\end{figure}
Una generalizaci\'on de la estructura de las estrellas es la \textit{rueda} de orden $n$, denotada $W_n$, que se construye a partir de un ciclo $C_{n-1}$ al a\~nadir un v\'ertice central $v_0$ que se conecta con todos los v\'ertices del ciclo. De manera formal, $$V(W_n) = \{v_0, v_1, v_2, \dots, v_{n-1}\}, \quad E(W_n) = E(C_{n-1}) \cup \{\{v_0, v_i\} : 1 \leq i \leq n-1\}.$$ El v\'ertice central tiene grado $n-1$, mientras que los v\'ertices del ciclo tienen grado 3, como se puede observar en la Figura \ref{fig:Ruedas}.
\newpage
\begin{figure}[h]
    \centering
    \begin{tikzpicture}
        % vertices W_3
        \node[draw, circle, fill=blue, label=above:{$ $}] (9) at (-3,1) {};
        \node[draw, circle, fill=blue, label={[label distance=13pt]355:$W_3$}] (10) at (-4,-1) {};
        \node[draw, circle, fill=blue, label=below:{$ $}] (11) at (-2,-1) {};

        % Aristas W_3
        \draw (9) -- (10);
        \draw (9) -- (11);
        \draw (10) -- (11);
        
        % vertices W_4
        \node[draw, circle, fill=blue, label=right:{$ $}] (1) at (0,0) {};
        \node[draw, circle, fill=blue, label=above:{$ $}] (2) at (0,1) {};
        \node[draw, circle, fill=blue, label={[label distance=13pt]355:$W_4$}] (3) at (-1,-1) {};
        \node[draw, circle, fill=blue, label=below:{$ $}] (4) at (1,-1) {};

        % Aristas W_4
        \draw (1) -- (2);
        \draw (1) -- (3);
        \draw (1) -- (4);
        \draw (2) -- (3);
        \draw (2) -- (4);
        \draw (3) -- (4);

        % vertices W_5
        \node[draw, circle, fill=blue, label=right:{$ $}] (4) at (3,0) {};
        \node[draw, circle, fill=blue, label=above:{$ $}] (5) at (4,1) {};
        \node[draw, circle, fill=blue, label=above:{$ $}] (6) at (2,1) {};
        \node[draw, circle, fill=blue, label={[label distance=13pt]355:$W_5$}] (7) at (2,-1) {};
        \node[draw, circle, fill=blue, label=below:{$ $}] (8) at (4,-1) {};

        % Aristas W_5
        \draw (4) -- (5);
        \draw (4) -- (6);
        \draw (4) -- (7);
        \draw (4) -- (8);
        \draw (5) -- (6);
        \draw (6) -- (7);
        \draw (7) -- (8);
        \draw (5) -- (8);
    \end{tikzpicture}   
\caption{Ejemplos de las gr\'aficas $W_3$, $W_4$ $W_5$.}
\label{fig:Ruedas}
\end{figure}
Una \textit{gr\'afica bipartita} es aquella cuyo conjunto de v\'ertices puede dividirse en dos subconjuntos disjuntos $V_1$ y $V_2$ tales que ninguna arista une v\'ertices de un mismo subconjunto. Una de las m\'as estudiadas dentro de esta familia es la \textit{gr\'afica bipartita completa}, denotada $K_{n,m}$, en la cual cada v\'ertice de $V_1$ est\'a conectado con todos los v\'ertices de $V_2$. Su definici\'on formal es$$V(K_{n,m}) = V_1 \cup V_2, \quad E(K_{n,m}) = \{\{u,v\} : u \in V_1,\,v \in V_2\},$$ donde $|V_1|=n$ y $|V_2|=m$. Esta gr\'afica tiene $n+m$ v\'ertices y $nm$ aristas como lo podemos observar en la Figura \ref{fig:Bipartitas}.
 
\begin{figure}[h]
    \centering
    \begin{tikzpicture}[every node/.style={circle, draw, fill=blue}]
        % --- K_{2,3} ---
        \node (a1) at (0,0.7) {};
        \node (a2) at (0,-0.7) {};
        \node (b1) at (2,1.2) {};
        \node (b2) at (2,0) {};
        \node [label={[label distance=12pt]200:$K_{2,3}$}] (b3) at (2,-1.2) {};
        \foreach \i in {1,2}{
            \foreach \j in {1,2,3}{
                \draw (a\i)--(b\j);
            }
        }
        % --- K_{3,2} ---
        \node (c1) at (4.5,1.2) {};
        \node (c2) at (4.5,0) {};
        \node [label={[label distance=12pt]340:$K_{3,2}$}] (c3) at (4.5,-1.2) {};
        \node (d1) at (6.5,0.7) {};
        \node (d2) at (6.5,-0.7) {};
        \foreach \i in {1,2,3}{
            \foreach \j in {1,2}{
                \draw (c\i)--(d\j);
            }
        }
        % --- K_{3,3} ---
        \node (e1) at (8.5,1.2) {};
        \node (e2) at (8.5,0) {};
        \node [label={[label distance=12pt]340:$K_{3,3}$}] (e3) at (8.5,-1.2) {};
        \node (f1) at (10.5,1.2) {};
        \node (f2) at (10.5,0) {};
        \node (f3) at (10.5,-1.2) {};
        \foreach \i in {1,2,3}{
            \foreach \j in {1,2,3}{
                \draw (e\i)--(f\j);
            }
        }
    \end{tikzpicture}
    \caption{Ejemplos de gr\'aficas bipartitas completas: $K_{2,3}$, $K_{3,2}$ y $K_{3,3}$.}
    \label{fig:Bipartitas}
\end{figure}
Otra familia de inter\'es es la de las \textit{rejillas} o \textit{mallas}, las cuales pueden interpretarse como el producto cartesiano de dos trayectorias. La rejilla $R_{n,m}$ se define formalmente como \begin{gather*} V(R_{n,m}) = \{(i,j) : 1 \leq i \leq n,\, 1 \leq j \leq m\}, \\ E(R_{n,m}) = \{\{(i,j),(i',j')\} : |i-i'| + |j-j'| = 1\}.\end{gather*} En esta estructura, cada v\'ertice est\'a conectado con sus vecinos inmediatos en direcci\'on horizontal o vertical, siempre que estos existan. Las rejillas son gr\'aficas planas, conexas y no regulares, cuyo grado m\'aximo es $4$. Para una mejor visualizaci\'on observe los ejemplos en la Figura \ref{fig:Rejillas}.  
\newpage
\begin{figure}[h!]
    \centering
    \begin{tikzpicture}[every node/.style={circle, draw, fill=blue}]
        % --- R_{3,2} ---
        \foreach \i in {1,...,3}{
            \foreach \j in {1,...,2}{
                \node (a\i\j) at (1.0*\j,1.0*\i) {};
            }
        }
        \foreach \i in {1,...,3}{
            \foreach \j [evaluate=\j as \jp using int(\j+1)] in {1}{
                \draw (a\i\j)--(a\i\jp);
            }
        }
        \foreach \j in {1,...,2}{
            \foreach \i [evaluate=\i as \ip using int(\i+1)] in {1,2}{
                \draw (a\i\j)--(a\ip\j);
            }
        }
        \node [draw=white, circle, fill=white, label={[label distance=12pt]200:$R_{3,2}$}] (a) at (2.5,1) {};
        % --- R_{2,3} ---
        \foreach \i in {1,...,2}{
            \foreach \j in {1,...,3}{
                \node (b\i\j) at (3.8+1.0*\j,1.0*\i) {};
            }
        }
        \foreach \i in {1,...,2}{
            \foreach \j [evaluate=\j as \jp using int(\j+1)] in {1,2}{
                \draw (b\i\j)--(b\i\jp);
            }
        }
        \foreach \j in {1,...,3}{
            \foreach \i [evaluate=\i as \ip using int(\i+1)] in {1}{
                \draw (b\i\j)--(b\ip\j);
            }
        }
        \node [draw=white, circle, fill=white, label={[label distance=35pt]350:$R_{2,3}$}] (b) at (4,1) {};
        % --- R_{3,3} ---
        \foreach \i in {1,...,3}{
            \foreach \j in {1,...,3}{
                \node (c\i\j) at (7.4+1.0*\j,1.0*\i) {};
            }
        }
        \foreach \i in {1,...,3}{
            \foreach \j [evaluate=\j as \jp using int(\j+1)] in {1,2}{
                \draw (c\i\j)--(c\i\jp);
            }
        }
        \foreach \j in {1,...,3}{
            \foreach \i [evaluate=\i as \ip using int(\i+1)] in {1,2}{
                \draw (c\i\j)--(c\ip\j);
            }
        }
        \node [draw=white, circle, fill=white, label={[label distance=25pt]350:$R_{3,3}$}] (c) at (7.9,1) {};
    \end{tikzpicture}
    \caption{Ejemplos de rejillas o mallas: $R_{3,2}$, $R_{2,3}$ y $R_{3,3}$.}
    \label{fig:Rejillas}
\end{figure}

\textbf{Aspectos de programaci\'on}
%Que es un for, un while un arreglo, que es sagemath que se usa sobre python

Para realizar las simulaciones y obtener los resultados experimentales del modelo, se emple\'o el software \textbf{SageMath} \cite{Matthes}, un entorno computacional dise\~nado para el estudio de estructuras algebraicas, combinatorias y gr\'aficas, que se basa en el lenguaje de programaci\'on \textbf{Python}. Este entorno permite la creaci\'on y manipulaci\'on eficiente de objetos matem\'aticos, as\'i como la automatizaci\'on de c\'alculos mediante rutinas programadas. SageMath dispone de funciones predefinidas para construir autom\'aticamente diversas familias de gr\'aficas, como las trayectorias $P_n$, los ciclos $C_n$, las estrellas $S_n$, las ruedas $W_n$ y las completas $K_n$ y $K_{n,m}$, lo que facilita su uso en el an\'alisis de propiedades estructurales. Sin embargo, en el caso de las rejillas bidimensionales $R_{n,m}$, cuya estructura no se encuentra implementada por defecto, fue necesario desarrollar una funci\'on espec\'ifica que definiera sus v\'ertices y aristas de manera sistem\'atica.

En el lenguaje Python, un \textbf{bucle for} permite ejecutar un conjunto de instrucciones un n\'umero determinado de veces, mientras que un \textbf{bucle while} repite la ejecuci\'on de un bloque de c\'odigo mientras se cumpla una condici\'on l\'ogica. Estas estructuras de control resultan fundamentales en la simulaci\'on de procesos iterativos, como el seguimiento temporal de la propagaci\'on de la infecci\'on. Asimismo, un \textbf{arreglo} o \textbf{lista} constituye una estructura de datos que almacena una secuencia de elementos accesibles mediante un \'indice, lo cual permite representar, por ejemplo, el estado de cada v\'ertice en la gr\'afica. Por otra parte, una \textbf{funci\'on} agrupa un conjunto de instrucciones que ejecutan una tarea espec\'ifica, favoreciendo la organizaci\'on del c\'odigo y su reutilizaci\'on. En este trabajo, cada funci\'on desarrollada corresponde a una operaci\'on clave del modelo, como la generaci\'on de estados iniciales, la actualizaci\'on temporal del sistema o la construcci\'on de las gr\'aficas.

Adem\'as, se emplearon \textbf{diccionarios}, estructuras que asocian claves con valores, permitiendo vincular cada v\'ertice con un color o estado determinado, lo cual facilita la representaci\'on visual y la interpretaci\'on de los resultados. Tambi\'en se utilizaron \textbf{tuplas}, que son secuencias inmutables de elementos, ideales para generar y almacenar las configuraciones iniciales de los v\'ertices sin que se alteren durante el proceso de simulaci\'on. Finalmente, SageMath integra una amplia variedad de librer\'ias que ampl\'ian las capacidades del lenguaje base, haciendo posible la ejecuci\'on de c\'alculos combinatorios y la visualizaci\'on directa de las gr\'aficas empleadas en el estudio.

En conjunto, los conceptos presentados en esta secci\'on permiten comprender tanto la formulaci\'on matem\'atica del modelo como su implementaci\'on computacional. La integraci\'on de la teor\'ia de gr\'aficas con las herramientas de programaci\'on proporciona un marco metodol\'ogico s\'olido para el an\'alisis de la din\'amica de propagaci\'on y curaci\'on en distintas topolog\'ias. De esta forma, los conocimientos te\'oricos y t\'ecnicos aqu\'i expuestos constituyen la base sobre la cual se desarrollan los algoritmos, las simulaciones y los teoremas presentados en las secciones posteriores.


\chapter{Desarrollo del proyecto}
%En Este cap\'itulo se incluir\'an los resultados obtenidos durante el proyecto

\section{El modelo de Mukwembi}
%Explicas el modelo, explicas los resultados del articulo
Como lo mencionamos antes, la gran mayor\'ia de las infecciones por VIH ocurren en los ganglios linf\'aticos, donde las c\'elulas CD4 susceptibles al VIH est\'an densamente empaquetadas. El modelo propuesto por Simon Mukwembi \cite{Mukwembi} aborda la din\'amica de la infecci\'on por VIH en los ganglios linf\'aticos mediante un enfoque basado en teor\'ia de gr\'aficas. La premisa central del modelo es que la propagaci\'on viral depende no solo de la presencia de c\'elulas infectadas, sino tambi\'en de la geometr\'ia y las interacciones locales de las c\'elulas dentro del tejido linf\'atico. Cada sitio del ganglio linf\'atico se representa como un v\'ertice de un gr\'afica, y las relaciones de vecindad entre sitios se modelan mediante aristas. Los v\'ertices pueden encontrarse en uno de tres estados: saludable, infectado o muerto, representados respectivamente por los colores blanco, rojo y negro en un esquema de coloraci\'on.

La propagaci\'on del virus en pasos de tiempo discretos cumple las siguientes reglas de transici\'on:

R1: Una c\'elula sana se infecta si tiene al menos un vecino infectado, de lo contrario, permanece sana.

R2: Una c\'elula infectada muere en el siguiente paso de tiempo.

R3: Una c\'elula muerta es reemplazada por una c\'elula infectada si tiene al menos $R$ vecinos infectados; de lo contrario, es reemplazada por una c\'elula sana.

Por lo tanto, el estado de las c\'elulas vecinas de una c\'elula muerta influye en el estado de la c\'elula que la reemplaza. Los cambios en $R$ determinan la din\'amica de propagaci\'on del virus del VIH.

La configuraci\'on inicial se compone de c\'elulas sanas con una peque\~na fracci\'on de c\'elulas infectadas. En este proyecto investigamos, utilizando m\'etodos de teor\'ia de gr\'aficas, c\'omo los cambios determinan la din\'amica de propagaci\'on del virus del VIH.

Dado un estado inicial, el sistema evoluciona bajo las reglas de transici\'on hasta alcanzar un estado estable, en el que ya no se producen cambios. Si despu\'es de un n\'umero finito de pasos todas las c\'elulas se encuentran saludables, diremos que el sistema se ha curado, en la Figura \ref{fig:Dinamica} se muestra una gr\'afica $G$ bajo esta din\'amica de propagaci\'on.

\begin{figure}[h]
    \centering
    \begin{subfigure}[h]{0.8\linewidth}
    \begin{tikzpicture}
        % vertices t=0
        \node[draw, circle, fill=white, label=above:{$t=0$}] (1) at (1.5,2.5) {};
        \node[draw, circle, fill=red, label=above:{$ $}] (2) at (0,1.5) {};
        \node[draw, circle, fill=white, label=left:{$ $}] (3) at (0.5,0) {};
        \node[draw, circle, fill=red, label=right:{$ $}] (4) at (2.5,0) {};
        \node[draw, circle, fill=white, label=above:{$ $}] (5) at (3,1.5) {};
        
        % Aristas t=0
        \draw (2) -- (3);
        \draw (3) -- (4);
        \draw (4) -- (5);
        \draw (5) -- (1);
        \draw (2) -- (1);
        \draw (2) -- (5);

        % vertices t=1
        \node[draw, circle, fill=red, label=above:{$t=1$}] (1) at (5.5,2.5) {};
        \node[draw, circle, fill=black, label=above:{$ $}] (2) at (4,1.5) {};
        \node[draw, circle, fill=red, label=left:{$ $}] (3) at (4.5,0) {};
        \node[draw, circle, fill=black, label=right:{$ $}] (4) at (6.5,0) {};
        \node[draw, circle, fill=red, label=above:{$ $}] (5) at (7,1.5) {};
        
        % Aristas t=1
        \draw (2) -- (3);
        \draw (3) -- (4);
        \draw (4) -- (5);
        \draw (5) -- (1);
        \draw (2) -- (1);
        \draw (2) -- (5);

        % vertices t=2
        \node[draw, circle, fill=black, label=above:{$t=2$}] (1) at (9.5,2.5) {};
        \node[draw, circle, fill=red, label=above:{$ $}] (2) at (8,1.5) {};
        \node[draw, circle, fill=black, label=left:{$ $}] (3) at (8.5,0) {};
        \node[draw, circle, fill=white, label=right:{$ $}] (4) at (10.5,0) {};
        \node[draw, circle, fill=black, label=above:{$ $}] (5) at (11,1.5) {};
        
        % Aristas t=2
        \draw (2) -- (3);
        \draw (3) -- (4);
        \draw (4) -- (5);
        \draw (5) -- (1);
        \draw (2) -- (1);
        \draw (2) -- (5);
    \end{tikzpicture} 
    \end{subfigure}
    \begin{subfigure}[h]{0.5\linewidth}
    \begin{tikzpicture}
        % vertices t=3
        \node[draw, circle, fill=white, label=above:{$t=3$}] (1) at (1.5,2.5) {};
        \node[draw, circle, fill=black, label=above:{$ $}] (2) at (0,1.5) {};
        \node[draw, circle, fill=white, label=left:{$ $}] (3) at (0.5,0) {};
        \node[draw, circle, fill=white, label=right:{$ $}] (4) at (2.5,0) {};
        \node[draw, circle, fill=white, label=above:{$ $}] (5) at (3,1.5) {};
        
        % Aristas t=3
        \draw (2) -- (3);
        \draw (3) -- (4);
        \draw (4) -- (5);
        \draw (5) -- (1);
        \draw (2) -- (1);
        \draw (2) -- (5);

        % vertices t=4
        \node[draw, circle, fill=white, label=above:{$t=4$}] (1) at (5.5,2.5) {};
        \node[draw, circle, fill=white, label=above:{$ $}] (2) at (4,1.5) {};
        \node[draw, circle, fill=white, label=left:{$ $}] (3) at (4.5,0) {};
        \node[draw, circle, fill=white, label=right:{$ $}] (4) at (6.5,0) {};
        \node[draw, circle, fill=white, label=above:{$ $}] (5) at (7,1.5) {};
        
        % Aristas t=4
        \draw (2) -- (3);
        \draw (3) -- (4);
        \draw (4) -- (5);
        \draw (5) -- (1);
        \draw (2) -- (1);
        \draw (2) -- (5);
    \end{tikzpicture} 
    \end{subfigure}
\caption{Ejemplo de la din\'amica de propagaci\'on en $G$ con $R=3$.}
\label{fig:Dinamica}
\end{figure}
Entre los hallazgos m\'as relevantes, Mukwembi demuestra que si $R = n-1$ donde $n$ es el n\'umero total de v\'ertices, el sistema converge de manera finita a un estado en el que todas las c\'elulas son saludables, independientemente de la configuraci\'on inicial. Este resultado sugiere que, en escenarios extremos donde las c\'elulas muertas son muy probablemente reemplazadas por c\'elulas saludables, la infecci\'on puede detenerse completamente, lo que tiene implicaciones te\'oricas sobre la eficacia de terapias que favorezcan la regeneraci\'on de c\'elulas sanas. Adem\'as, el modelo permite calcular de manera algebraica el n\'umero de c\'elulas infectadas en cualquier instante y establecer l\'imites superiores e inferiores para las poblaciones de c\'elulas saludables, infectadas y muertas. Esto se logra mediante la introducci\'on de subgrafos latentes, los cuales representan agrupaciones de c\'elulas que mantienen la infecci\'on activa y act\'uan como fuentes de propagaci\'on viral.

Como antes dicho Mukwembi modela la situaci\'on mediante una gr\'afica $G$, donde el conjunto de v\'ertices es el conjunto de sitios del ganglio linf\'atico y dos v\'ertices estan\'an unidos por una arista si y solo si los sitios correspondientes en los ganglios linf\'aticos son vecinos. En cada instante $i$, los v\'ertices de $G$ se colorean con un color $f_{i,R}: V(G) \rightarrow \{0,1,2\}$, tal que $f_{i,R}(v)=0$ si el ganglio linf\'aticos est\'a sano en el instante $i$, $f_{i,R}(v)=1$ si est\'a infectado y $f_{i,R}(v)=2$ si est\'a muerto. Obs\'ervese que en la configuraci\'on inicial, $f_{0,R}$, satisface que $f_{0,R}(v)\neq 2$ para cada $v\in V(G)$. Dado $R$, un par\'ametro de reemplazo, Mukwembi defini\'o la din\'amica del modelo mediante un conjunto de reglas que pueden escribirse de la siguiente manera. Sea $N_{i,j}(v) =\{ x\in N(v) \mid f_{i,R}(v)=j\}$ y sea $d_{i,j}(v) = |N_{i,j}(v)|$, entonces: 

$$f_{i+1,R}(v)= \left\{ \begin{array}{lcl} 
0 & \text{si }  & f_{i,R}(v)=0\text{ y }d_{i,1}(v)=0 \text{, o } \\
 & &  f_{i,R}(v)=2 \text{ y } d_{i,1}(v)<R\\ \\
1 & \text{si } & f_{i,R}(v)=0\text{ y }d_{i,1}(v)\geq 1 \text{, o } \\
 & &  f_{i,R}(v)=2 \text{ y } d_{i,1}(v)\geq R\\ \\
2 & \text{si } &   f_{i,R}(v)=1
\end{array} \right.$$

El modelo de Mukwembi determina una secuencia de funciones que dependen de $R$, digamos $P_{f_0,R}(G)=(f_{0,R}, f_{1,R}, f_{2,R}, \ldots)$.

Mukwembi introduce en su modelo \cite{Mukwembi} el concepto de subgrafo latente como una estructura que permite la persistencia de la infecci\'on dentro del sistema. Formalmente, un subgrafo latente de una gr\'afica $G$, denotado por $K_{a,b}$, es un subgrafo bipartito completo cuyas clases de v\'ertices son $V_1$ y $V_2$, que cumple las siguientes condiciones: $a,b \geq R$, donde $R$ es el par\'ametro de reemplazo del modelo, y existe un instante de tiempo $i$ tal que todos los v\'ertices de $V_1$ se encuentran infectados en el tiempo $i$, y todos los v\'ertices de $V_2$ se encuentran infectados en el tiempo $i+1$, es decir, $f_{i,R}(V_1)=1=f_{i+1,R}(V_2)$. De esta forma, en un subgrafo latente la infecci\'on alterna entre las dos clases de v\'ertices de manera peri\'odica, produciendo un ciclo de reinfecci\'on continua que impide la eliminaci\'on completa del virus, este comportamiento se puede observar en la Figura \ref{fig:Latente}. Dichas estructuras modelan, desde el punto de vista biol\'ogico, los reservorios latentes del VIH, es decir, aquellos conglomerados de c\'elulas en los ganglios linf\'aticos que mantienen activa la infecci\'on a lo largo del tiempo.
\begin{figure}[h]
    \centering
    \begin{tikzpicture}[every node/.style={circle, draw}]
        % t=0
        \node [fill=red, label={[label distance=10pt]3:$t=0$}] (e4) at (0.5,3.6) {};
        \node [fill=red] (e1) at (0.5,2.4) {};
        \node [fill=red] (e2) at (0.5,1.2) {};
        \node [fill=red] (e3) at (0.5,0) {};
        \node [fill=white] (f4) at (2.5,3.6) {};
        \node [fill=white] (f1) at (2.5,2.4) {};
        \node [fill=white] (f2) at (2.5,1.2) {};
        \node [fill=white] (f3) at (2.5,0) {};
        \foreach \i in {1,2,3,4}{
            \foreach \j in {1,2,3,4}{
                \draw (e\i)--(f\j);
            }
        }

        % t=1
        \node [fill=black, label={[label distance=10pt]3:$t=1$}] (e4) at (3.5,3.6) {};
        \node [fill=black] (e1) at (3.5,2.4) {};
        \node [fill=black] (e2) at (3.5,1.2) {};
        \node [fill=black] (e3) at (3.5,0) {};
        \node [fill=red] (f4) at (5.5,3.6) {};
        \node [fill=red] (f1) at (5.5,2.4) {};
        \node [fill=red] (f2) at (5.5,1.2) {};
        \node [fill=red] (f3) at (5.5,0) {};
        \foreach \i in {1,2,3,4}{
            \foreach \j in {1,2,3,4}{
                \draw (e\i)--(f\j);
            }
        }

        % t=2
        \node [fill=red, label={[label distance=10pt]3:$t=2$}] (e4) at (6.5,3.6) {};
        \node [fill=red] (e1) at (6.5,2.4) {};
        \node [fill=red] (e2) at (6.5,1.2) {};
        \node [fill=red] (e3) at (6.5,0) {};
        \node [fill=black] (f4) at (8.5,3.6) {};
        \node [fill=black] (f1) at (8.5,2.4) {};
        \node [fill=black] (f2) at (8.5,1.2) {};
        \node [fill=black] (f3) at (8.5,0) {};
        \foreach \i in {1,2,3,4}{
            \foreach \j in {1,2,3,4}{
                \draw (e\i)--(f\j);
            }
        }
    
         % t=3
        \node [fill=black, label={[label distance=10pt]3:$t=3$}] (e4) at (9.5,3.6) {};
        \node [fill=black] (e1) at (9.5,2.4) {};
        \node [fill=black] (e2) at (9.5,1.2) {};
        \node [fill=black] (e3) at (9.5,0) {};
        \node [fill=red] (f4) at (11.5,3.6) {};
        \node [fill=red] (f1) at (11.5,2.4) {};
        \node [fill=red] (f2) at (11.5,1.2) {};
        \node [fill=red] (f3) at (11.5,0) {};
        \foreach \i in {1,2,3,4}{
            \foreach \j in {1,2,3,4}{
                \draw (e\i)--(f\j);
            }
        }

        % t=4
        \node [fill=red, label={[label distance=10pt]3:$t=4$}] (e4) at (12.5,3.6) {};
        \node [fill=red] (e1) at (12.5,2.4) {};
        \node [fill=red] (e2) at (12.5,1.2) {};
        \node [fill=red] (e3) at (12.5,0) {};
        \node [fill=black] (f4) at (14.5,3.6) {};
        \node [fill=black] (f1) at (14.5,2.4) {};
        \node [fill=black] (f2) at (14.5,1.2) {};
        \node [fill=black] (f3) at (14.5,0) {};
        \foreach \i in {1,2,3,4}{
            \foreach \j in {1,2,3,4}{
                \draw (e\i)--(f\j);
            }
        }
    \end{tikzpicture}
    \caption{Ejemplo de subgr\'afica $4-latente$.}
    \label{fig:Latente}
\end{figure}

\begin{theorem}\label{Teo:Mukwembi}
    Sea $G$ un gr\'afica de orden $n$ y $f_{0,R}$ una configuraci\'on inicial. Si $R= n-1$, entonces existe un entero $t_0$ tal que para cada $t\geq t_0$, $f_{t,R}(v)=0$ para cada $v\in V(G)$.
\end{theorem}
\begin{proof}
    Sea $f_{0,R}$ una configuraci\'on inicial arbitraria en la que los v\'ertices de $G$ pueden encontrarse en estado sano ($0$) o infectado ($1$), pero ninguno en estado muerto ($2$). Denotemos por $S = \{v \in V(G) : f_{0,R}(v) = 1\}$ al conjunto de v\'ertices infectados inicialmente, y por $s = |S|$ su cardinalidad. Para cada instante discreto $t$, definimos $r(t)$, $w(t)$ y $b(t)$ como el n\'umero de v\'ertices infectados, sanos y muertos, respectivamente. Si $S = V(G)$, el proceso es trivial, pues todos los v\'ertices mueren en el siguiente paso y posteriormente se vuelven sanos, cumpli\'endose el resultado. Supongamos entonces que $S \neq V(G)$.

    Sea $e = \max\{d(v,S) : v \in V(G)\}$ la excentricidad del conjunto inicial de v\'ertices infectados, y definamos para cada $k = 0, 1, \ldots, e$ el conjunto $L_k = \{\, v \in V(G) : d(v,S) = k \,\}$, con $|L_k| = \ell_k$. Por construcci\'on, los conjuntos $L_k$ son disjuntos y ning\'un v\'ertice de $L_k$ es adyacente a uno de $L_{k+2}$, pues esto violar\'ia la definici\'on de distancia.

    Bajo la hip\'otesis $R = n-1$, la condici\'on $d_{t,1}(v) \geq R$ significa que el v\'ertice $v$ tiene a todos los dem\'as v\'ertices de la gr\'afica infectados simult\'aneamente. Por lo tanto, un v\'ertice muerto s\'olo puede volverse infectado si todos los dem\'as v\'ertices de la gr\'afica est\'an infectados en ese instante. Este hecho ser\'a clave para analizar la evoluci\'on del sistema.

    Consideremos primero el caso en que no existe ning\'un v\'ertice de grado $n-1$. En este escenario, ning\'un v\'ertice puede cumplir $d_{t,1}(v) \geq R$, de modo que la condici\'on de infecci\'on de un v\'ertice muerto nunca se satisface. As\'i, todo v\'ertice muerto en el tiempo $t$ ser\'a necesariamente sano en el instante $t+1$. Aplicando las reglas de transici\'on de forma sincr\'onica sobre las capas $L_0, L_1, \ldots, L_e$, se obtiene que los v\'ertices de $L_t$ son infectados en el tiempo $t$, mientras que los de $L_{t-1}$ mueren y el resto permanecen sanos. En consecuencia, para $1 \le t \le e$, se cumple $f_{t,R}(v) = 1$ si $v \in L_t$ y $f_{t,R}(v)=2$ si $v \in L_{t-1}$. Una vez que la onda de infecci\'on alcanza la capa m\'as lejana $L_e$, en los dos pasos siguientes todos los v\'ertices muertos son reemplazados por v\'ertices sanos, de manera que para todo $t \ge e+2$ se cumple $f_{t,R}(v)=0$ para todo $v \in V(G)$. As\'i, el sistema alcanza un estado completamente sano en un n\'umero finito de pasos.

    Consideremos ahora el caso en que existe al menos un v\'ertice de grado $n-1$. Denotemos por $v$ a dicho v\'ertice, el cual es adyacente a todos los dem\'as de la gr\'afica. Si $v \in S$, es decir, si $v$ est\'a infectado en el instante inicial $f_{0,R}(v) = 1$, la excentricidad $e(S)$ es igual a $1$. En $t=1$, por la regla de muerte, $f_{1,R}(v) = 2$, mientras que todos los dem\'as v\'ertices, al tener un vecino infectado, se infectan: $f_{1,R}(u) = 1$ para todo $u \neq v$. Por tanto, $r(1) = n-1$ y $b(1)=1$. En el siguiente paso, los v\'ertices infectados mueren $f_{2,R}(u)=2$, y como el v\'ertice $v$ muerto tiene $n-1$ vecinos infectados, se cumple $d_{1,1}(v)=n-1=R$, lo cual implica que $f_{2,R}(v)=1$. En consecuencia, $r(2)=1$ y $b(2)=n-1$. Finalmente, en $t=3$, los v\'ertices muertos, al tener menos de $R$ vecinos infectados, se vuelven sanos, y el \'unico infectado muere, obteniendo $f_{3,R}(v)=2$ y $f_{3,R}(u)=0$ para todo $u\neq v$. Por tanto, para $t \ge 3$, $f_{t,R}(v)=0$ para todo $v \in V(G)$. Esta evoluci\'on produce la secuencia $r(0)=1$, $r(1)=n-1$, $r(2)=1$, $r(t)=0$ para $t \ge 3$, que corresponde al caso excepcional mencionado por Mukwembi.

    Si, en cambio, $v \notin S$, entonces $v$ debe estar a distancia $1$ de alg\'un v\'ertice infectado inicial, es decir, $v \in L_1$. En este caso, la excentricidad satisface $e \le 2$. Si $e=1$, la infecci\'on se propaga \'unicamente entre $L_0$ y $L_1$: los v\'ertices de $L_0$ mueren y los de $L_1$ se infectan en $t=1$; luego, en $t=2$, los de $L_0$ se vuelven sanos y los de $L_1$ mueren, alcanzando la salud total en $t=3$. Si $e=2$, un razonamiento an\'alogo muestra que en $t=1$, $f_{1,R}(L_0)=2$, $f_{1,R}(L_1)=1$ y $f_{1,R}(L_2)=0$; luego, en $t=2$, los v\'ertices de $L_2$ se infectan, los de $L_1$ mueren y los de $L_0$ sanan; finalmente, en $t=3$, mueren los \'ultimos infectados y todos los dem\'as v\'ertices se vuelven sanos. Dado que ning\'un v\'ertice tiene $n-1$ vecinos infectados simult\'aneamente (pues $\ell_2 < n-1$), no puede reactivarse la infecci\'on, y por tanto para $t \ge 4$ todo el sistema permanece sano.

    En conclusi\'on, bajo la hip\'otesis $R = n - 1$, la infecci\'on desaparece completamente en un n\'umero finito de pasos, independientemente de la configuraci\'on inicial, es decir, la din\'amica definida por $f_{0,R}$ siempre converge al estado sano global $f_{t,R} = 0$, independientemente de la configuraci\'on inicial.
\end{proof}
\section{El n\'umero de VIH de una gr\'afica}
%Explicas el invariante, que por PBO esta bien definido, y aqui iran los resultados teoricos como el del Delta+1
\begin{definition}\label{Def:VIH(G)}
    Sea $G$ una gr\'afica simple y conexa. Definimos el par\'ametro $VIH(G)$ como el m\'inimo $R$ tal que la enfermedad se extingue para cualquier condici\'on inicial en $G$. 
\end{definition}
Como vimos antes, Mukwembi demostr\'o que para toda gr\'afica $G$ con $n$ v\'ertices, se cumple que $VIH(G) \leq n - 1$. Este resultado implica que, bajo ciertas condiciones, el sistema puede alcanzar un estado saludable en un n\'umero finito de pasos, y que el valor $n-1$ act\'ua como una cota superior del par\'ametro. Por lo tanto, resulta natural investigar el valor m\'inimo posible de $VIH(G)$, as\'i como las gr\'aficas en las que esta cota se alcanza.
%A continuacion, se presentan los resultados teoricos asociados al modelo propuesto.
\begin{lemma}\label{Lem:BorrarVertices}
    Sea $G$ una gr\'afica simple y sea $H := G - v$ para un v\'ertice $v \in V(G)$ de grado cero. Entonces se cumple que $VIH(H) \leq VIH(G)$.
\end{lemma}
\begin{proof}
    Recordemos que, para un par\'ametro de reemplazo $R \in \mathbb{N}$ y una configuraci\'on inicial $f_{0,R}: V(G) \rightarrow \{0,1,2\}$, la din\'amica del modelo de Mukwembi define en cada instante $i$ una funci\'on $f_{i,R}$ seg\'un las reglas de evoluci\'on que dependen \'unicamente de los estados de los vecinos de cada v\'ertice. Adem\'as, para cualquier v\'ertice $x$ y estado $j \in \{0,1,2\}$, se define el conjunto $N_{i,j}(x)=\{y \in N(x) \mid f_{i,R}(y)=j\}$ y su cardinalidad $d_{i,j}(x)=|N_{i,j}(x)|$. Dichas definiciones determinan completamente la transici\'on de estados en la red linf\'atica modelada por la gr\'afica $G$.

    Sea $R$ un entero tal que es un valor \emph{vanishing} para $G$, es decir, para cualquier configuraci\'on inicial $f_{0,R}$ la sucesi\'on de estados $f_{i,R}$ con $i\geq0$ converge hacia una configuraci\'on completamente sana, en la cual todos los v\'ertices est\'an en estado $0$. Consideremos una configuraci\'on inicial arbitraria $f_0$ sobre $G$ y definamos su restricci\'on a $H$ como $f_0^H(u)=f_0(u)$ para todo $u \in V(H)=V(G)\setminus\{v\}$. Denotemos por $f_{i,R}^G$ y $f_{i,R}^H$ las configuraciones en el tiempo $i$ obtenidas al aplicar las reglas de evoluci\'on sobre $G$ y $H$, respectivamente, partiendo de las configuraciones iniciales $f_0$ y $f_0^H$.

    Como el v\'ertice $v$ tiene grado cero en $G$, se cumple que $N_H(u)=N_G(u)$ para todo $u \in V(H)$. Dado que las reglas del modelo de Mukwembi dependen exclusivamente de los vecinos de cada v\'ertice, y las configuraciones iniciales coinciden sobre $V(H)$, la evoluci\'on temporal del proceso es id\'entica en ambas gr\'aficas para todos los v\'ertices de $H$. Es decir, se cumple que $f_{i,R}^H(u)=f_{i,R}^G(u)$ para todo $u \in V(H)$ y para todo $i \geq 0$. El v\'ertice $v$ no influye ni es influido por los dem\'as, puesto que carece de vecinos, y por tanto su presencia o ausencia no altera la din\'amica del resto del sistema.

    De esta forma, si el par\'ametro $R$ garantiza la curaci\'on completa de $G$, entonces el mismo valor $R$ garantiza tambi\'en la curaci\'on completa de $H$. En consecuencia, el m\'inimo valor de $R$ que provoca la desaparici\'on de la infecci\'on en $H$ no puede ser mayor que el correspondiente en $G$. Por lo tanto, se concluye que $VIH(H) \leq VIH(G)$, lo cual demuestra el resultado.
\end{proof}
\begin{lemma}\label{Lem:BorrarAristas}
    Sea $G$ una gr\'afica simple y sea $H:=G-e$ para una arista $e\in E(G)$. Entonces $VIH(H)\leq VIH(G)$.
\end{lemma}
\begin{proof}
    Sea $R\in\mathbb{N}$ un valor fijo y consid\'erese una configuraci\'on inicial $f_0:V(G)\to\{0,1,2\}$. Para cada una de las gr\'aficas $X\in\{G,H\}$, denotemos por $f^X_{t,R}$ la configuraci\'on en el instante $t$ obtenida al aplicar las reglas del modelo de Mukwembi con par\'ametro $R$ sobre la gr\'afica $X$, comenzando desde $f_0$.
    
    Sea adem\'as $I^X_t=\{v\in V(G): f^X_{t,R}(v)=1\}$ el conjunto de v\'ertices infectados en el instante $t$ en $X$, y definamos $d^X_{t,1}(v)=|\{u\in N_X(v): f^X_{t,R}(u)=1\}|$ como el n\'umero de vecinos infectados de un v\'ertice $v$ en el instante $t$.

    Demostramos por inducci\'on sobre $t$, que para todo $t\geq 0$ se cumple la inclusi\'on $I^H_t\subseteq I^G_t$. La afirmaci\'on es trivial para $t=0$, ya que $f^H_{0,R}=f^G_{0,R}=f_0$.
    
    Supongamos que $I^H_t\subseteq I^G_t$ para alg\'un $t$. Entonces, para cada v\'ertice $v$, como $N_H(v)\subseteq N_G(v)$ al eliminar una arista, y por hip\'otesis inductva $I^H_t\subseteq I^G_t$, se sigue que el n\'umero de vecinos infectados en $H$ nunca excede al de $G$, es decir, $d^H_{t,1}(v)\le d^G_{t,1}(v)$ para todo $v \in V(G)$.
    
    Consideremos ahora un v\'ertice $v\in I^H_{t+1}$. Seg\'un las reglas de evoluci\'on del modelo, esto ocurre en uno de los dos siguientes casos; caso 1 $f^H_{t,R}(v)=0$ y $d^H_{t,1}(v)\ge 1$. Caso 2 $f^H_{t,R}(v)=2$ y $d^H_{t,1}(v)\ge R$.
    
    En ambos casos, dado que $d^G_{t,1}(v)\ge d^H_{t,1}(v)$ y $f^G_{t,R}(v)\in\{0,2\}$, el mismo tipo de transici\'on ocurre en $G$ y, por lo tanto, $v\in I^G_{t+1}$. Esto demuestra $I^H_{t+1}\subseteq I^G_{t+1}$, completando as\'i el paso inducti\'vo. 

    Finalmente, Si $R$ se anula para $G$, entonces existe un entero $t_0$ tal que $I^G_t=\emptyset$ para todo $t\ge t_0$. Por la inclusi\'on anterior, tambi\'en se deduce que $I^H_t=\emptyset$ para todo $t\ge t_0$. Adem\'as, dado que los v\'ertices infectados siempre pasan al estado muerto en el siguiente paso y no pueden aparecer nuevas infecciones en ausencia de vecinos infectados, dos pasos despu\'es todos los v\'ertices estar\'an sanos. Por lo tanto, $R$ tambi\'en se anula para $H$, lo que demuestra que $VIG(H)\le VIH(G)$.
\end{proof}
\begin{theorem}\label{Teo:Delta+1}
    Sea $G$ una gr\'afica. Si $G$ contiene una arista, entonces $\Delta+1\geq VIH(G) \geq 2$ 
\end{theorem}
\begin{proof}
    Comenzamos estableciendo la cota inferior, si $R=1$ basta demostrar que existe una configuraci\'on inicial $f_{0,1}$, tal que en la sucesi\'on $P_{f_{0}, 1} (G)$ siempre existe un v\'ertice $v$ cuya evoluci\'on nunca alcanza el estado sano, es decir, que para todo entero $t$ se cumple $f_{t,1}(v)\neq 0$.

    Sea $v$ un v\'ertice cualquiera con al menos un vecino, es decir, tal que $N(v)\neq \emptyset$. Definimos la configuraci\'on inicial $f_{0,1}$ de la siguiente manera, asignando el estado infectado al v\'ertice $v$, es decir, $f_{0,1}(v) =1$, y el estado sano a todos los dem\'as v\'ertices, esto es $f_{0,1}(x)=0$ para todo $x\in V(G)\setminus \{v\}$. 
    
    Entonces, siguiendo las reglas de evoluci\'on del modelo de Mukwemni con $R=1$, el v\'ertice $v$ alternar\'a c\'iclicamente entre los estados infecto y muerto. Dado que $v$ posee al menos un vecino, se cumple que para todo entero $k$, $f_{2k,1}(v)= 1$ y $f_{2k+1,1}(v)= 2$, porque $f_{2k,1}(u)=2$ y $f_{2k+1,1}(u)=1$ para todo $u \in N(v)$. Por otro lado, sus vecinos alternar\'an de forma complementaria entre los estados infectado y muerto, de manera que nunca se alcanza el estado completamente sano. Por tanto, el virus persiste indefinidamente en la red y el n\'umero VIH de la gr\'afica satisface $VIH(G)\geq2$.

    Para demostrar ahora la cota superior, Consideremos una configuraci\'on inicial $f_{0,\Delta+1}$ y sea $U_0$ el conjunto de v\'ertices que est\'an inicialmente infectados, es decir, $U_0 = \{ x \in V(G) : f_{0,\Delta +1}(x)=1\}$. Para cada $i=1, \ldots \ell$, denotemos por $U_i$ el conjunto de v\'ertices situados a distancia $i$ de $U_0$.

    Observemos que bajo el par\'ametro $R= \Delta +1$, ning\'un v\'ertice puede alcanzar el estado infectado en el paso siguiente si en el instante actual se encuentra en estado muerto. En efecto, dado que ning\'un v\'ertice posee m\'as de $\Delta$ vecinos, la condici\'on de tener al menos $R= \Delta +1$ vecinos infectados nunca se satisface. Por lo tanto, si en alg\'un instante $t$ un v\'ertice $v$ est\'a muerto, esto es, $f_{0,\Delta+1}(v)=2$, necesariamente en el siguiente paso se cumple $f_{t+1, \Delta+1} (v) =0$.

    De esta manera, la din\'amica del sistema progresa en capas a lo largo de la distancia desde los v\'ertices inicialmente infectados. En particular, para cada $i =1 \ldots \ell -2$, se cumple que:$$f_{i,\Delta+1}(v)= \left\{ \begin{array}{ll} 
    2 \text{ si } x\in U_{i-1} \\
    1 \text{ si } x\in U_{i}\\
    0 \text{ e.o.c} 
    \end{array} \right.$$
    Al cabo de un n\'umero finito de pasos, concretamente cuando $t=\ell$, todos los v\'ertices de la gr\'afica se encuentran en estado sano, es decir, $f_{\ell, \Delta +1} (v) = 0$ para cada $v\in V(G)$ y el resultado se cumple. 
    
    Por consiguiente, el par\'ametro $R= \Delta +1$ garantiza la desaparici\'on total de la infecci\'on, y concluimos que $\Delta+1\geq VIH(G)$. Combinando ambas desigualdades, se obtiene finalmente $\Delta+1\geq VIH(G) \geq 2$.
\end{proof}
\section{Exploraci\'on computacional}
%El pseudocodigo, y los resulatdos de las tablas y de las familias
Con el objetivo de determinar de manera experimental el n\'umero $VIH(G)$ asociado a distintas familias de gr\'aficas, se implement\'o un modelo computacional en el entorno de \textit{SageMath}. Dicho modelo se basa en un aut\'omata celular que simula la propagaci\'on y recuperaci\'on de una infecci\'on sobre una gr\'afica simple $G$.

Cada v\'ertice de la gr\'afica puede encontrarse en uno de tres estados posibles: sano (0), infectado (1) o muerto (2). La evoluci\'on del sistema se realiza de manera discreta en el tiempo, y en cada iteraci\'on los v\'ertices actualizan su estado en funci\'on de los estados de sus vecinos, de acuerdo con el conjunto de reglas de Mukwembi. El par\'ametro $R$ desempe\~na un papel fundamental dentro del modelo, pues controla el n\'umero m\'inimo de vecinos infectados necesarios para que un v\'ertice recuperado vuelva a infectarse. El valor de $R$ se ajusta en cada simulaci\'on hasta encontrar el m\'inimo que garantiza la curaci\'on total del sistema, identific\'andose as\'i con el n\'umero $VIH(G)$.

Para la implementaci\'on del modelo, se defini\'o un conjunto de funciones fundamentales que permiten generar los estados iniciales, aplicar las reglas del aut\'omata y construir las distintas topolog\'ias de gr\'aficas. En primer lugar, el Algoritmo \ref{fun:EstadoInicial} de la funci\'on \textit{EstadoInicial(G)} genera una configuraci\'on inicial aleatoria de la gr\'afica $G$, asignando a cada v\'ertice un valor 0 o 1 seg\'un su estado de salud. Esta funci\'on proporciona el punto de partida de la simulaci\'on, pues cada posible configuraci\'on inicial representa un escenario independiente de infecci\'on.

\begin{algorithm}[h]
\caption{EstadoInicial($G$)}
\label{fun:EstadoInicial}
\begin{algorithmic}[1]
\State \textbf{Entrada:} Una gr\'afica $G$.
\State \textbf{Salida:} Un vector \textit{estados} que representa el estado inicial aleatorio de cada v\'ertice.
\State Sea $n \gets |V(G)|$.
\For{$i \gets 1$ \textbf{hasta} $n$}
    \State Asignar aleatoriamente $estados[i] \in \{0,1\}$.
\EndFor
\State \Return $estados$
\end{algorithmic}
\end{algorithm}
\newpage
Posteriormente, el Algoritmo \ref{fun:HacerDicCol} de la funci\'on \textit{HacerDicCol(estados)} se encarga de asociar colores a los v\'ertices de acuerdo con su estado, lo cual permite una visualizaci\'on clara del proceso de propagaci\'on. Los v\'ertices sanos se representan en color blanco, los infectados en rojo y los recuperados en negro. Aunque esta funci\'on no modifica el comportamiento del modelo, resulta \'util para representar gr\'aficamente la din\'amica del sistema y facilitar la interpretaci\'on de los resultados.
\begin{algorithm}[h!]
\caption{HacerDicCol(\textit{estados})}
\label{fun:HacerDicCol}
\begin{algorithmic}[1]
\State \textbf{Entrada:} Lista \textit{estados} que representa el estado de cada v\'ertice.
\State \textbf{Salida:} Diccionario $d$ que asocia colores a conjuntos de v\'ertices.
\State Inicializar listas vac\'ias: $ceros$, $unos$, $dos$.
\For{$i \gets 1$ \textbf{hasta} longitud(\textit{estados})}
    \If{$estados[i] = 0$} agregar $i$ a $ceros$.
    \ElsIf{$estados[i] = 1$} agregar $i$ a $unos$.
    \ElsIf{$estados[i] = 2$} agregar $i$ a $dos$.
    \EndIf
\EndFor
\State Definir $d = \{\text{white}: ceros, \text{red}: unos, \text{black}: dos\}$
\State \Return $d$
\end{algorithmic}
\end{algorithm}

El Algoritmo \ref{fun:Paso} que es sobre la funci\'on \textit{Paso(G, estado, R)} constituye el n\'ucleo del modelo, ya que implementa la regla de evoluci\'on temporal del modelo. Dada una gr\'afica $G$, un vector de estados y un valor del par\'ametro $R$, la funci\'on actualiza el estado de cada v\'ertice siguiendo las reglas del modelo de Mukwembi, las cuales recordemos que son:

R1: Un v\'ertice sano pasa a infectado si al menos uno de sus vecinos est\'a infectado;

R2: Un v\'ertice infectado pasa a recuperado en el siguiente paso; y

R3: Un v\'ertice recuperado puede reinfectarse si cuenta con al menos $R$ vecinos infectados, o de lo contrario regresa al estado sano.

Estas reglas reflejan un equilibrio entre propagaci\'on y curaci\'on, y permiten observar la din\'amica del sistema hacia un estado estacionario o de curaci\'on total.
\newpage
\begin{algorithm}[h]
\caption{Paso($G$, \textit{estado}, $R$)}
\label{fun:Paso}
\begin{algorithmic}[1]
\State \textbf{Entrada:} Gr\'afica $G$, lista de estados \textit{estado}, par\'ametro $R$.
\State \textbf{Salida:} Nuevo vector de estados actualizado.
\State Copiar $nuevoestado \gets estado$.
\For{cada v\'ertice $i$ en $V(G)$}
    \If{$estado[i] = 0$}
        \If{alg\'un vecino de $i$ est\'a en estado $1$}
            \State $nuevoestado[i] \gets 1$
        \EndIf
    \ElsIf{$estado[i] = 1$}
        \State $nuevoestado[i] \gets 2$
    \ElsIf{$estado[i] = 2$}
        \State $vecinos\_infectados \gets$ n\'umero de vecinos de $i$ tales que $estado[vecino] = 1$
        \If{$vecinos\_infectados \ge R$}
            \State $nuevoestado[i] \gets 1$
        \Else
            \State $nuevoestado[i] \gets 0$
        \EndIf
    \EndIf
\EndFor
\State \Return $nuevoestado$
\end{algorithmic}
\end{algorithm}

\begin{algorithm}[H]
\caption{Rejilla($n$, $m$)}
\label{fun:Rejilla}
\begin{algorithmic}[1]
\State \textbf{Entrada:} Enteros $n, m$.
\State \textbf{Salida:} Gr\'afica $G$ con estructura de rejilla $n \times m$.
\State Inicializar una gr\'afica vac\'ia $G$.
\State $contador \gets 0$
\For{$i \gets 1$ \textbf{hasta} $n$}
    \For{$j \gets 1$ \textbf{hasta} $m$}
        \State A\~nadir v\'ertice $contador$ a $G$.
        \State $contador \gets contador + 1$
    \EndFor
\EndFor
\State A\~nadir aristas verticales: conectar $a,b$ si $a-b = n$.
\State A\~nadir aristas horizontales: conectar $a,b$ si $b-a = 1$ y $b \bmod n \neq 0$.
\State \Return $G$
\end{algorithmic}
\end{algorithm}

Finalmente, el Algoritmo \ref{fun:Rejilla} de la funci\'on \textit{Rejilla(n,m)} genera una gr\'afica con estructura de rejilla rectangular de dimensiones $n \times m$. En esta gr\'afica, cada v\'ertice se conecta con sus vecinos adyacentes de manera horizontal y vertical, formando una red bidimensional regular. Esta funci\'on resulta especialmente \'util para analizar el comportamiento del modelo en topolog\'ias espaciales, donde la conectividad local es uniforme y las interacciones entre v\'ertices son homog\'eneas.

Una vez definidas las funciones que gobiernan la din\'amica del modelo, se implement\'o un procedimiento general que permite calcular experimentalmente el n\'umero $VIH(G)$ para cualquier familia de gr\'aficas, el cual se presenta en el Algoritmo \ref{fun:General}. El algoritmo consiste en recorrer todas las configuraciones iniciales posibles de los v\'ertices, y simular la evoluci\'on del sistema durante un n\'umero determinado de pasos de tiempo. Para cada configuraci\'on inicial se eval\'ua si la gr\'afica alcanza la curaci\'on completa, es decir, si al cabo de la simulaci\'on todos los v\'ertices se encuentran en el estado sano.

El proceso se repite para distintos valores del par\'ametro $R$, increment\'andolo gradualmente hasta encontrar el m\'inimo valor que garantiza la curaci\'on en todos los casos iniciales posibles, este valor se interpreta como el n\'umero $VIH(G)$. Este algoritmo fue aplicado a distintas familias de gr\'aficas con el prop\'osito de determinar experimentalmente el valor de $VIH(G)$ en cada caso. .

En el Algoritmo \ref{fun:VIH(Pn)} se presenta el pseudoc\'odigo correspondiente al c\'alculo experimental del n\'umero $VIH(P_n)$, donde se consideran todas las configuraciones iniciales posibles y se analiza la evoluci\'on del sistema hasta alcanzar un estado estacionario. De manera an\'aloga, el Algoritmo \ref{fun:VIH(Cn)} describe el procedimiento utilizado para las gr\'aficas $C_n$, mientras que el Algoritmo \ref{fun:VIH(Sn)} y el Algoritmo \ref{fun:VIH(Wn)} muestran la implementaci\'on aplicada a las gr\'aficas de estrellas y ruedas, respectivamente.

Por otro lado, el Algoritmo \ref{fun:VIH(Kn)} detalla el proceso seguido para las gr\'aficas completas $K_n$, en las cuales la alta conectividad influye de manera significativa en la din\'amica de propagaci\'on del modelo. En el caso de las gr\'aficas bipartitas completas, el procedimiento se describe en el Algoritmo \ref{fun:VIH(Kn,m)}, donde, debido a la simetr\'ia de la estructura, \'unicamente se considera el caso $n \geq m$, sin p\'erdida de generalidad. Finalmente, el Algoritmo \ref{fun:VIH(Rn,m)} corresponde al an\'alisis de las rejillas $R_{n,m}$

\newpage
\begin{algorithm}[h]
\caption{C\'alculo del n\'umero VIH(G) para una familia de gr\'aficas}
\label{fun:General}
\begin{algorithmic}[1]
\State \textbf{Entrada:} Familia de gr\'aficas $\mathcal{F}$, par\'ametro de reemplazo $R$, n\'umero m\'aximo de pasos $T$.
\State \textbf{Salida:} Valor aproximado del n\'umero $VIH(G)$ para cada $G \in \mathcal{F}$.
\For{cada orden $n$ en el rango de inter\'es}
    \State Generar la gr\'afica $G \in \mathcal{F}$ de orden $n$.
    \State Inicializar listas vac\'ias: \textit{estados}, \textit{sumas}.
    \For{cada configuraci\'on inicial $\varphi \in \{0,1\}^{|V(G)|}$}
        \State Sea $local[0] \gets \varphi$
        \For{$t \gets 0$ \textbf{hasta} $T-1$}
            \State $local[t+1] \gets Paso(G, local[t], R)$
        \EndFor
        \State Agregar $\sum local[T]$ a la lista \textit{sumas}.
    \EndFor
    \If{$\sum \textit{sumas} = 0$}
        \State Imprimir: Para $G$ con $n$ v\'ertices y $R$, todos los casos iniciales se curan.
    \Else
        \State Imprimir: Para $G$ con $n$ v\'ertices y $R$, existen configuraciones que no se curan.
    \EndIf
\EndFor
\end{algorithmic}
\end{algorithm}
\begin{algorithm}[H]
\caption{C\'alculo de $VIH(P_n)$}
\label{fun:VIH(Pn)}
\begin{algorithmic}[1]
\For{$n \gets 1$ \textbf{hasta} $10$}
    \State Generar $G \gets$ \texttt{PathGraph}$(n)$.
    \State Fijar $R \gets 2$.
    \State Aplicar el procedimiento general del Algoritmo \ref{fun:General}.
\EndFor
\end{algorithmic}
\end{algorithm}
\begin{algorithm}[H]
\caption{C\'alculo de $VIH(C_n)$}
\label{fun:VIH(Cn)}
\begin{algorithmic}[1]
\For{$n \gets 1$ \textbf{hasta} $10$}
    \State Generar $G \gets$ \texttt{CycleGraph}$(n)$.
    \State Fijar $R \gets 3$.
    \State Aplicar el procedimiento general del Algoritmo \ref{fun:General}.
\EndFor
\end{algorithmic}
\end{algorithm}
\newpage
\begin{algorithm}[H]
\caption{C\'alculo de $VIH(S_n)$}
\label{fun:VIH(Sn)}
\begin{algorithmic}[1]
\For{$n \gets 1$ \textbf{hasta} $10$}
    \State Generar $G \gets$ \texttt{StarGraph}$(n)$.
    \State Fijar $R \gets 6$.
    \State Aplicar el procedimiento general del Algoritmo \ref{fun:General}.
\EndFor
\end{algorithmic}
\end{algorithm}
\begin{algorithm}[H]
\caption{C\'alculo de $VIH(W_n)$}
\label{fun:VIH(Wn)}
\begin{algorithmic}[1]
\For{$n \gets 1$ \textbf{hasta} $10$}
    \State Generar $G \gets$ \texttt{WheelGraph}$(n)$.
    \State Fijar $R \gets 9$.
    \State Aplicar el procedimiento general del Algoritmo \ref{fun:General}.
\EndFor
\end{algorithmic}
\end{algorithm}
\begin{algorithm}[H]
\caption{C\'alculo de $VIH(K_n)$}
\label{fun:VIH(Kn)}
\begin{algorithmic}[1]
\For{$n \gets 1$ \textbf{hasta} $10$}
    \State Generar $G \gets$ \texttt{CompleteGraph}$(n)$.
    \State Fijar $R \gets 6$.
    \State Aplicar el procedimiento general del Algoritmo \ref{fun:General}.
\EndFor
\end{algorithmic}
\end{algorithm}
\begin{algorithm}[H]
\caption{C\'alculo de $VIH(K_{n,m})$}
\label{fun:VIH(Kn,m)}
\begin{algorithmic}[1]
\For{$n \gets 1$ \textbf{hasta} $10$}
    \For{$m \gets 1$ \textbf{hasta} $10$}
        \If{$n \ge m$}
            \State Generar $G \gets$ \texttt{CompleteBipartiteGraph}$(n,m)$.
            \State Fijar $R \gets 6$.
            \State Aplicar el procedimiento general del Algoritmo \ref{fun:General}.
        \EndIf
    \EndFor
\EndFor
\end{algorithmic}
\end{algorithm}
\newpage
\begin{algorithm}[H]
\caption{C\'alculo de $VIH(R_{n,m})$}
\label{fun:VIH(Rn,m)}
\begin{algorithmic}[1]
\For{$n \gets 10$ \textbf{hasta} $10$}
    \For{$m \gets 10$ \textbf{hasta} $10$}
        \If{$n \ge m$}
            \State Generar $G \gets$ \texttt{Rejilla}$(n,m)$.
            \State Fijar $R \gets 4$.
            \State Aplicar el procedimiento general del Algoritmo \ref{fun:General}.
        \EndIf
    \EndFor
\EndFor
\end{algorithmic}
\end{algorithm}
El objetivo de esta fase fue determinar de forma experimental el n\'umero $VIH(G)$, entendido como el menor valor del par\'ametro $R$ que garantiza la curaci\'on del sistema para todas las configuraciones iniciales posibles.

En cada simulaci\'on, se exploraron exhaustivamente todos los estados iniciales de los v\'ertices, representados por los vectores binarios en $\{0,1\}^{|V(G)|}$. A partir de cada configuraci\'on inicial $\varphi$, se aplic\'o la funci\'on \textit{Paso(G, estado, R)} durante 20 iteraciones consecutivas, verificando al final del proceso si la suma de los estados finales era igual a cero, lo que indica que la gr\'afica alcanz\'o un estado completamente sano. Si esta condici\'on se cumpl\'ia para todas las configuraciones posibles, se conclu\'ia que el valor elegido de $R$ era suficiente, y por tanto, que $R = VIH(G)$.

\section{Resultados sobre VIH(G)}

Los resultados obtenidos a partir de la exploraci\'on computacional descrita en la secci\'on anterior permitieron identificar patrones regulares en el comportamiento del n\'umero $VIH(G)$ en funci\'on de la estructura de las distintas familias de gr\'aficas analizadas. La observaci\'on de estos patrones condujo a la formulaci\'on de una serie de teoremas que relacionan el valor del n\'umero $VIH(G)$ con propiedades estructurales espec\'ificas, tales como el grado m\'aximo $\Delta(G)$, la regularidad, la conectividad y la densidad de aristas.

En esta secci\'on se presentan dichos teoremas junto con sus respectivas demostraciones. El prop\'osito es establecer una base te\'orica que respalde las tendencias observadas emp\'iricamente, proporcionando as\'i una interpretaci\'on formal del comportamiento del modelo de infecci\'on bajo diferentes topolog\'ias.

\subsection{VIH en estrellas $S_n$}
Posteriormente, se analiz\'o la familia de gr\'aficas estrella $S_n$, caracterizada por la presencia de un v\'ertice central $C$ conectado a todos los dem\'as $n$. En este tipo de estructuras, el v\'ertice central $C$ act\'ua como un punto de alta concentraci\'on de infecciones, lo que dificulta la curaci\'on completa del sistema. 

\begin{figure}[h!]
    \centering
    \begin{tabular}{|c|cccccccccc|}

        \hline
        \textbf{n} & \textbf{1} & \textbf{2} & \textbf{3} & \textbf{4} & \textbf{5} & \textbf{6} & \textbf{7} & \textbf{8} & \textbf{9} & \textbf{10}\\
        \hline
        \textbf{$VIH(S_n)$} & 2 & 2 & 2 & 3 & 3 & 4 & 4 & 5 & 5 & 6 \\
        \hline
    \end{tabular}
    \caption{Tabla $VIH(S_n)$}
    \label{tabla:VIH(S_n)}
\end{figure}
Los resultados obtenidos en la Figura \ref{tabla:VIH(S_n)}, mostraron que el par\'ametro $R$ necesario para lograr la curaci\'on es considerablemente mayor, lo que sugiere una fuerte dependencia del n\'umero $VIH(S_n)$ respecto al grado del v\'ertice central. 
\begin{theorem}\label{Teo:VIH(Sn)}
    $$
    VIH(S_n) = \left\lbrace
    \begin{array}{ll}
        2 \textup{ si } n = 1\\
        \lceil \frac{n+1}{2} \rceil \textup{ e.o.c}
    \end{array}
    \right.
    $$
\end{theorem}
\begin{proof}
    Comencemos considerando el caso en que $n=1$. En este caso, la gr\'afica $S_1$ est\'a compuesta por el centro $C$ y un \'unico v\'ertice exterior $u$. Si tomamos $R=1$ y la configuraci\'on inicial $f_{0,1}(C)=0$, $f_{0,1}(u)=1$, observamos que en $t=0$ el centro tiene $d_{0,1}(C)=1$, por lo que en el siguiente paso se tiene $f_{1,1}(C)=1$ y $f_{1,1}(u)=2$. En el tiempo $t=1$, el exterior muerto tiene $d_{1,1}(u)=1$, de manera que $f_{2,1}(u)=1$ y $f_{2,1}(C)=2$. 
    
    Si continuamos analizando la din\'amica, el proceso se repite indefinidamente, alternando entre los estados en los que el centro est\'a infectado y el exterior muerto, y viceversa, en la Figura \ref{S1yR=1} se puede observar mejor este comportamiento. Por lo tanto, para $R=1$, la infecci\'on no se extingue. Sin embargo, si $R=2$, un v\'ertice muerto requiere al menos dos vecinos infectados para re-infectarse, lo cual es imposible en $S_1$ ya que cada v\'ertice tiene un solo vecino. En consecuencia, para $R=2$ el sistema se cura completamente en un n\'umero finito de pasos, como se observa en la Figura \ref{S1yR=2}, de modo que $VIH(S_1)=2$.
    
    \begin{figure}[h]
    \centering
    \begin{subfigure}[h]{0.65\linewidth}
    \begin{tikzpicture}
    %tiempo 0
    % vertices
    \node[draw, circle, fill=red, label=below:{$t=0$}] (1) at (0,0) {};
    \node[draw, circle, fill=white, label=above:{$C$}] (2) at (0,1) {};

    % Aristas
    \draw (1) -- (2);
    
    %tiempo 1
    % vertices
    \node[draw, circle, fill=black, label=below:{$t=1$}] (3) at (2,0) {};
    \node[draw, circle, fill=red, label=above:{$C$}] (4) at (2,1) {};

    % Aristas
    \draw (3) -- (4);

    %tiempo 2
    % vertices
    \node[draw, circle, fill=red, label=below:{$t=2$}] (5) at (4,0) {};
    \node[draw, circle, fill=black, label=above:{$C$}] (6) at (4,1) {};

    % Aristas
    \draw (5) -- (6);

    %tiempo 3
    % vertices
    \node[draw, circle, fill=black, label=below:{$t=3$}] (7) at (6,0) {};
    \node[draw, circle, fill=red, label=above:{$C$}] (8) at (6,1) {};

    % Aristas
    \draw (7) -- (8);

    %tiempo 4
    % vertices
    \node[draw, circle, fill=red, label=below:{$t=4$}] (9) at (8,0) {};
    \node[draw, circle, fill=black, label=above:{$C$}] (10) at (8,1) {};

    % Aristas
    \draw (9) -- (10);
    \end{tikzpicture}
    \caption{Comportamiento con $R=1$}
    \label{S1yR=1}
    \end{subfigure}
    \begin{subfigure}[h]{0.5\linewidth}
    \begin{tikzpicture}
    %tiempo 0
    % vertices
    \node[draw, circle, fill=red, label=below:{$t=0$}] (1) at (0,0) {};
    \node[draw, circle, fill=white, label=above:{$C$}] (2) at (0,1) {};

    % Aristas
    \draw (1) -- (2);
    
    %tiempo 1
    % vertices
    \node[draw, circle, fill=black, label=below:{$t=1$}] (3) at (2,0) {};
    \node[draw, circle, fill=red, label=above:{$C$}] (4) at (2,1) {};

    % Aristas
    \draw (3) -- (4);

    %tiempo 2
    % vertices
    \node[draw, circle, fill=white, label=below:{$t=2$}] (5) at (4,0) {};
    \node[draw, circle, fill=black, label=above:{$C$}] (6) at (4,1) {};

    % Aristas
    \draw (5) -- (6);

    %tiempo 3
    % vertices
    \node[draw, circle, fill=white, label=below:{$t=3$}] (7) at (6,0) {};
    \node[draw, circle, fill=white, label=above:{$C$}] (8) at (6,1) {};

    % Aristas
    \draw (7) -- (8);
    \end{tikzpicture}
    \caption{Comportamiento con $R=2$}
    \label{S1yR=2}
    \end{subfigure}
    \caption{Comportamiento del $VIH$ en $S_1$}
    \label{S1}
    \end{figure}
    Consideremos ahora el caso en que $n=2$ y $R<\lceil (n+1)/2\rceil$. Como $\lceil (n+1)/2\rceil=\lceil 3/2\rceil=2$, tenemos que $R=1$. Si en la configuraci\'on inicial $f_{0,1}(C)=0$ y ambos exteriores est\'an infectados, es decir, $f_{0,1}(v)=1$ para todo v\'ertice exterior $v$, entonces $d_{0,1}(C)=2$ y en el siguiente paso el centro se infecta $f_{1,1}(C)=1$, mientras que los exteriores infectados mueren $f_{1,1}(v)=2$. 
    
    En el tiempo $t=1$, los exteriores muertos tienen $d_{1,1}(v)=1$, por lo que en el tiempo $t=2$ se re-infectan $f_{2,1}(v)=1$ y el centro pasa a muerto $f_{2,1}(C)=2$. Este patr\'on se repite peri\'odicamente, alternando entre estados en los que el centro est\'a infectado y los exteriores muertos, y viceversa. Por lo tanto, para $R=1$ la infecci\'on persiste indefinidamente y no se extingue, como se puede observar en la Figura \ref{S2yR=1}. 
    
    En cambio, si $R\ge 2$, los exteriores muertos no pueden re-infectarse, pues cada uno tiene como m\'aximo un vecino infectado. As\'i, tras algunos pasos, todos los v\'ertices pasan al estado sano, garantizando la extinci\'on de la enfermedad. En consecuencia, para $n=2$ se cumple que $VIH(S_2)=2$, en la Figura \ref{S2yR=2} se observan dos posibles estados iniciales para este caso en particular y el como con $R=2$ la enfermedad se cura.

    \begin{figure}[h]
    \centering
    \begin{subfigure}[h]{0.75\linewidth}
    \begin{tikzpicture}
    %tiempo 0
    % vertices
    \node[draw, circle, fill=white, label=above:{$C$}] (1) at (0,0) {};
    \node[draw, circle, fill=red, label=below:{$ $}] (2) at (1,1) {};
    \node[draw, circle, fill=red, label=below:{$t=0$}] (3) at (1,-1) {};

    % Aristas
    \draw (1) -- (2);
    \draw (1) -- (3);

    %tiempo 1
    % vertices
    \node[draw, circle, fill=red, label=above:{$C$}] (4) at (2,0) {};
    \node[draw, circle, fill=black, label=below:{$ $}] (5) at (3,1) {};
    \node[draw, circle, fill=black, label=below:{$t=1$}] (6) at (3,-1) {};

    % Aristas
    \draw (4) -- (5);
    \draw (4) -- (6);

    %tiempo 2
    % vertices
    \node[draw, circle, fill=black, label=above:{$C$}] (7) at (4,0) {};
    \node[draw, circle, fill=red, label=below:{$ $}] (8) at (5,1) {};
    \node[draw, circle, fill=red, label=below:{$t=2$}] (9) at (5,-1) {};

    % Aristas
    \draw (7) -- (8);
    \draw (7) -- (9);

    %tiempo 3
    % vertices
    \node[draw, circle, fill=red, label=above:{$C$}] (10) at (6,0) {};
    \node[draw, circle, fill=black, label=below:{$ $}] (11) at (7,1) {};
    \node[draw, circle, fill=black, label=below:{$t=3$}] (12) at (7,-1) {};

    % Aristas
    \draw (10) -- (11);
    \draw (10) -- (12);
    
    %tiempo 4
    % vertices
    \node[draw, circle, fill=black, label=above:{$C$}] (13) at (8,0) {};
    \node[draw, circle, fill=red, label=below:{$ $}] (14) at (9,1) {};
    \node[draw, circle, fill=red, label=below:{$t=4$}] (15) at (9,-1) {};

    % Aristas
    \draw (13) -- (14);
    \draw (13) -- (15);
    \end{tikzpicture}
    \caption{Comportamiento con $R<\lceil \frac{n+1}{2} \rceil$.}
    \label{S2yR=1}
    \end{subfigure}
    \begin{subfigure}[h]{0.6\linewidth}
    \begin{tikzpicture}
    %tiempo 0
    % vertices
    \node[draw, circle, fill=white, label=above:{$C$}] (1) at (0,0) {};
    \node[draw, circle, fill=red, label=below:{$ $}] (2) at (1,1) {};
    \node[draw, circle, fill=red, label=below:{$t=0$}] (3) at (1,-1) {};

    % Aristas
    \draw (1) -- (2);
    \draw (1) -- (3);

    %tiempo 1
    % vertices
    \node[draw, circle, fill=red, label=above:{$C$}] (4) at (2,0) {};
    \node[draw, circle, fill=black, label=below:{$ $}] (5) at (3,1) {};
    \node[draw, circle, fill=black, label=below:{$t=1$}] (6) at (3,-1) {};

    % Aristas
    \draw (4) -- (5);
    \draw (4) -- (6);

    %tiempo 2
    % vertices
    \node[draw, circle, fill=black, label=above:{$C$}] (7) at (4,0) {};
    \node[draw, circle, fill=white, label=below:{$ $}] (8) at (5,1) {};
    \node[draw, circle, fill=white, label=below:{$t=2$}] (9) at (5,-1) {};

    % Aristas
    \draw (7) -- (8);
    \draw (7) -- (9);

    %tiempo 3
    % vertices
    \node[draw, circle, fill=white, label=above:{$C$}] (10) at (6,0) {};
    \node[draw, circle, fill=white, label=below:{$ $}] (11) at (7,1) {};
    \node[draw, circle, fill=white, label=below:{$t=3$}] (12) at (7,-1) {};

    % Aristas
    \draw (10) -- (11);
    \draw (10) -- (12);
    \end{tikzpicture}
    \end{subfigure}

    \begin{subfigure}[b]{0.6\linewidth}
    \begin{tikzpicture}
    %tiempo 0
    % vertices
    \node[draw, circle, fill=red, label=above:{$C$}] (1) at (0,0) {};
    \node[draw, circle, fill=white, label=below:{$ $}] (2) at (1,1) {};
    \node[draw, circle, fill=red, label=below:{$t=0$}] (3) at (1,-1) {};

    % Aristas
    \draw (1) -- (2);
    \draw (1) -- (3);

    %tiempo 1
    % vertices
    \node[draw, circle, fill=black, label=above:{$C$}] (4) at (2,0) {};
    \node[draw, circle, fill=red, label=below:{$ $}] (5) at (3,1) {};
    \node[draw, circle, fill=black, label=below:{$t=1$}] (6) at (3,-1) {};

    % Aristas
    \draw (4) -- (5);
    \draw (4) -- (6);

    %tiempo 2
    % vertices
    \node[draw, circle, fill=white, label=above:{$C$}] (7) at (4,0) {};
    \node[draw, circle, fill=black, label=below:{$ $}] (8) at (5,1) {};
    \node[draw, circle, fill=white, label=below:{$t=2$}] (9) at (5,-1) {};

    % Aristas
    \draw (7) -- (8);
    \draw (7) -- (9);

    %tiempo 3
    % vertices
    \node[draw, circle, fill=white, label=above:{$C$}] (10) at (6,0) {};
    \node[draw, circle, fill=white, label=below:{$ $}] (11) at (7,1) {};
    \node[draw, ellipse, fill=white, label=below:{$t=3$}] (12) at (7,-1) {};

    % Aristas
    \draw (10) -- (11);
    \draw (10) -- (12);
    \end{tikzpicture}
    %\caption{$f_{0,R}(C)=1$, $f_{0,R}(v_1)=0$ y $f_{0,R}(v_2)=1$}
    \caption{Comportamiento con $R=\lceil \frac{n+1}{2} \rceil$=2.}
    \label{S2yR=2}
    \end{subfigure}
    \caption{Comportamiento del $VIH$ en $S_2$.}
    \label{S2}
    \end{figure}
    Analicemos ahora el caso en que $n\ge 3$ y $R<\lceil (n+1)/2\rceil$. Supongamos una configuraci\'on inicial en la que el centro $C$ est\'a sano $f_{0,R}(C)=0$, y de los $n$ v\'ertices exteriores, supongamos que hay un conjunto de estos v\'ertices que lo denotaremos como $B$, en donde, $|B|=n-\lceil n/2\rceil$ est\'an infectados y el otro conjunto $A$ de v\'ertices exteriores $|A|=\lceil n/2\rceil$ est\'an sanos. En el tiempo inicial $t=0$, el centro tiene $d_{0,1}(C)=|B|\ge 1$, por lo que en el siguiente paso se infecta $f_{1,R}(C)=1$, y los exteriores infectados mueren $f_{1,R}(B)=2$, mientras que los v\'ertices del conjunto $A$ permanecen sanos $f_{1,R}(A)=0$. En el tiempo $t=1$, los exteriores muertos tienen $d_{1,1}(B)=1$ por tener al centro infectado, por lo que si $R=1$, vuelven a infectarse en el tiempo $t=2$ $f_{2,R}(B)=1$, el centro pasa a muerto $f_{2,R}(C)=2$ y como este estaba infectado en el tiempo anterior, entonces $f_{2,R}(A)=1$. En $t=3$, los v\'ertices exteriores mueren, y el centro vuelve a infectarse $f_{3,R}(C)=1$, ya que todos sus vecinos estaban infectados en el paso anterior y este ciclo se repite, manteniendo la infecci\'on activa de manera indefinida como se muestra en la Figura \ref{SnyR<vih}. Si $R$ es mayor pero sigue cumpliendo $R<\lceil (n+1)/2\rceil$, la estructura del proceso sigue siendo oscilatoria: en algunos pasos el centro est\'a infectado y los exteriores muertos, y en otros el centro muerto y los exteriores infectados. Por lo tanto, para cualquier $R<\lceil (n+1)/2\rceil$ existen configuraciones iniciales que provocan un comportamiento peri\'odico en el que la infecci\'on nunca se extingue, lo que implica que dichos valores de $R$ no son suficiente para que la enfermedad se cure.

    \begin{figure}[h]
    \centering
    \begin{subfigure}[h]{0.8\linewidth}
    \begin{tikzpicture}
    %tiempo 0
    % vertices
    \node[draw, circle, fill=white, label=above:{$C$}] (1) at (0,0) {};
    \node[draw, ellipse, minimum height=60pt,minimum width=20pt, fill=white, label=above:{$A$}] (2) at (-1,0) {};
    \node[draw, ellipse, minimum height=60pt,minimum width=20pt, fill=red, label={[label distance=1pt]253:$t=0$}] (3) at (1,0) {};
    \node[draw, ellipse, minimum height=60pt,minimum width=20pt, fill=red, label=above:{$B$}] (a) at (1,0) {};

    % Aristas
    \draw (1) -- (2);
    \draw (1) -- (3);

    %tiempo 1
    % vertices
    \node[draw, circle, fill=red, label=above:{$C$}] (4) at (4,0) {};
    \node[draw, ellipse, minimum height=60pt,minimum width=20pt, fill=white, label=above:{$A$}] (5) at (3,0) {};
    \node[draw, ellipse, minimum height=60pt,minimum width=20pt, fill=black, label={[label distance=1pt]253:$t=1$}] (6) at (5,0) {};
    \node[draw, ellipse, minimum height=60pt,minimum width=20pt, fill=black, label=above:{$B$}] (b) at (5,0) {};

    % Aristas
    \draw (4) -- (5);
    \draw (4) -- (6);
    
    %tiempo 2
    % vertices
    \node[draw, circle, fill=black, label=above:{$C$}] (7) at (8,0) {};
    \node[draw, ellipse, minimum height=60pt,minimum width=20pt, fill=red, label=above:{$A$}] (8) at (7,0) {};
    \node[draw, ellipse, minimum height=60pt,minimum width=20pt, fill=red, label={[label distance=1pt]253:$t=2$}] (9) at (9,0) {};
    \node[draw, ellipse, minimum height=60pt,minimum width=20pt, fill=red, label=above:{$B$}] (c) at (9,0) {};

    % Aristas
    \draw (7) -- (8);
    \draw (7) -- (9);
    \end{tikzpicture}
    \end{subfigure}
    
    \begin{subfigure}[h]{0.5\linewidth}
    \begin{tikzpicture}
    %tiempo 3
    % vertices
    \node[draw, circle, fill=red, label=above:{$C$}] (1) at (0,0) {};
    \node[draw, ellipse, minimum height=60pt,minimum width=20pt, fill=black, label=above:{$A$}] (2) at (-1,0) {};
    \node[draw, ellipse, minimum height=60pt,minimum width=20pt, fill=black, label={[label distance=1pt]253:$t=3$}] (3) at (1,0) {};
    \node[draw, ellipse, minimum height=60pt,minimum width=20pt, fill=black, label=above:{$B$}] (a) at (1,0) {};

    % Aristas
    \draw (1) -- (2);
    \draw (1) -- (3);
    
    %tiempo 4
    % vertices
    \node[draw, circle, fill=black, label=above:{$C$}] (4) at (4,0) {};
    \node[draw, ellipse, minimum height=60pt,minimum width=20pt, fill=red, label=above:{$A$}] (5) at (3,0) {};
    \node[draw, ellipse, minimum height=60pt,minimum width=20pt, fill=red, label={[label distance=1pt]253:$t=4$}] (6) at (5,0) {};
    \node[draw, ellipse, minimum height=60pt,minimum width=20pt, fill=red, label=above:{$B$}] (b) at (5,0) {};

    % Aristas
    \draw (4) -- (5);
    \draw (4) -- (6);
    \end{tikzpicture}
    \end{subfigure}
    \caption{Comportamiento del $VIH$ en $S_n$ con $n\geq3$ y $R<\lceil \frac{n+1}{2} \rceil$.}
    \label{SnyR<vih}
    \end{figure}
    Finalmente, consideremos el caso en que $R= \lceil (n+1)/2\rceil$. En este escenario se verificar\'a que, sin importar la configuraci\'on inicial, la infecci\'on siempre se extingue en un n\'umero finito de pasos. Sea $f_{0,R}$ cualquier configuraci\'on inicial con $f_{0,R}(v)\neq 2$ para todo $v$. Si el centro $C$ est\'a sano $f_{0,R}(C)=0$ y todos los exteriores tambi\'en lo est\'an $f_{0,R}(v)=0$, el sistema est\'a desde el inicio curado. Si el centro est\'a infectado $f_{0,R}(C)=1$ y todos los exteriores tambi\'en lo est\'an $f_{0,R}(v)=1$, entonces en el siguiente tiempo $t=1$ todos los v\'ertices mueren $f_{1,R}(C)=2$ y $f_{1,R}(v)=2$ para todo $v$. En el tiempo $t=2$, cada v\'ertice muerto se re-infectar\'a solo si tiene al menos $R$ vecinos infectados en $t=1$. Sin embargo, como cada v\'ertice en tiempo $t=1$ se encontraba muerto, entonces esto nunca sucede, por lo que todos los v\'ertices pasan al estado sano en $t=2$, es decir, $f_{2,R}(C)=0$ y $f_{2,R}(v)=0$. 
    
    Si el centro est\'a sano $f_{0,R}(C)=0$ y todos los exteriores infectados $f_{0,R}(v)=1$, en $t=1$ el centro se infecta $f_{1,R}(C)=1$ y los exteriores infectados mueren $f_{1,R}(v)=2$. En $t=2$, los exteriores muertos tienen $d_{1,1}(v)=1<R$, por lo que pasan a sanos $f_{2,R}(v)=0$, y el centro pasa a muerto $f_{2,R}(C)=2$. En el siguiente paso, al no tener vecinos infectados, el centro se vuelve sano $f_{3,R}(C)=0$, y la infecci\'on se extingue completamente en un tiempo finito, este comportamiento lo podemos visualizar en la Figura \ref{SnyR=VIH(A)}. 
    
    Si el centro est\'a infectado $f_{0,R}(C)=1$ y los exteriores sanos $f_{0,R}(v)=0$, entonces en $t=1$ el centro pasa a muerto $f_{1,R}(C)=2$ y los v\'ertices exteriores se infectan $f_{1,R}(v)=1$. En $t=2$, los exteriores infectados mueren $f_{2,R}(v)=2$ y como $R=\lceil (n+1)/2\rceil$ entonces el centro vuelve a infectarse $f_{2,R}(C)=1$. En el siguiente paso, el centro muerto, mientras que los v\'ertices exteriores al solo tener un v\'ertice vecino que esta infectado el cual no es suficiente para que estos se vuelvan a infectar por lo tanto se curan $f_{3,R}(v)=0$. Finalmente, al no tener vecinos infectados, el centro se vuelve sano $f_{4,R}(C)=0$, y la infecci\'on se extingue completamente y todos los v\'ertices permanecen sanos, por lo que el sistema tambi\'en se cura r\'apidamente, como podemos observar en la Figura \ref{SnyR=VIH(B)}. 

    \begin{figure}[h]
    \centering
    \begin{subfigure}[h]{0.55\linewidth}
    \begin{tikzpicture}
    %tiempo 0
    % vertices
    \node[draw, circle, fill=white, label=above:{$C$}] (1) at (0,0) {};
    \node[draw, ellipse, minimum height=60pt,minimum width=20pt, fill=red, label={[label distance=1pt]267:$t=0$}] (2) at (1,0) {};
    \node[draw, ellipse, minimum height=60pt,minimum width=20pt, fill=red, label=above:{$v$}] (a) at (1,0) {};

    % Aristas
    \draw (1) -- (2);

    %tiempo 1
    % vertices
    \node[draw, circle, fill=red, label=above:{$C$}] (3) at (2,0) {};
    \node[draw, ellipse, minimum height=60pt,minimum width=20pt, fill=black, label={[label distance=1pt]267:$t=1$}] (4) at (3,0) {};
    \node[draw, ellipse, minimum height=60pt,minimum width=20pt, fill=black, label=above:{$v$}] (b) at (3,0) {};

    % Aristas
    \draw (3) -- (4);
    
    %tiempo 2
    % vertices
    \node[draw, circle, fill=black, label=above:{$C$}] (5) at (4,0) {};
    \node[draw, ellipse, minimum height=60pt,minimum width=20pt, fill=white, label={[label distance=1pt]267:$t=2$}] (6) at (5,0) {};
    \node[draw, ellipse, minimum height=60pt,minimum width=20pt, fill=white, label=above:{$v$}] (c) at (5,0) {};

    % Aristas
    \draw (5) -- (6);

    %tiempo 3
    % vertices
    \node[draw, circle, fill=white, label=above:{$C$}] (7) at (6,0) {};
    \node[draw, ellipse, minimum height=60pt,minimum width=20pt, fill=white, label={[label distance=1pt]267:$t=3$}] (8) at (7,0) {};
    \node[draw, ellipse, minimum height=60pt,minimum width=20pt, fill=white, label=above:{$v$}] (d) at (7,0) {};

    % Aristas
    \draw (7) -- (8);
    \end{tikzpicture}
    \caption{$f_{0,R}(C)=0$ y $f_{0,R}(v)=1$}
    \label{SnyR=VIH(A)}
    \end{subfigure}

    \begin{subfigure}[h]{0.7\linewidth}
    \begin{tikzpicture}
    %tiempo 0
    % vertices
    \node[draw, circle, fill=red, label=above:{$C$}] (1) at (0,0) {};
    \node[draw, ellipse, minimum height=60pt,minimum width=20pt, fill=white, label={[label distance=1pt]267:$t=0$}] (2) at (1,0) {};
    \node[draw, ellipse, minimum height=60pt,minimum width=20pt, fill=white, label=above:{$v$}] (a) at (1,0) {};

    % Aristas
    \draw (1) -- (2);

    %tiempo 1
    % vertices
    \node[draw, circle, fill=black, label=above:{$C$}] (3) at (2,0) {};
    \node[draw, ellipse, minimum height=60pt,minimum width=20pt, fill=red, label={[label distance=1pt]267:$t=1$}] (4) at (3,0) {};
    \node[draw, ellipse, minimum height=60pt,minimum width=20pt, fill=red, label=above:{$v$}] (b) at (3,0) {};

    % Aristas
    \draw (3) -- (4);
    
    %tiempo 2
    % vertices
    \node[draw, circle, fill=red, label=above:{$C$}] (5) at (4,0) {};
    \node[draw, ellipse, minimum height=60pt,minimum width=20pt, fill=black, label={[label distance=1pt]267:$t=2$}] (6) at (5,0) {};
    \node[draw, ellipse, minimum height=60pt,minimum width=20pt, fill=black, label=above:{$v$}] (c) at (5,0) {};

    % Aristas
    \draw (5) -- (6);

    %tiempo 3
    % vertices
    \node[draw, circle, fill=black, label=above:{$C$}] (7) at (6,0) {};
    \node[draw, ellipse, minimum height=60pt,minimum width=20pt, fill=white, label={[label distance=1pt]267:$t=3$}] (8) at (7,0) {};
    \node[draw, ellipse, minimum height=60pt,minimum width=20pt, fill=white, label=above:{$v$}] (d) at (7,0) {};

    % Aristas
    \draw (7) -- (8);

    %tiempo 4
    % vertices
    \node[draw, circle, fill=white, label=above:{$C$}] (9) at (8,0) {};
    \node[draw, ellipse, minimum height=60pt,minimum width=20pt, fill=white, label={[label distance=1pt]267:$t=4$}] (10) at (9,0) {};
    \node[draw, ellipse, minimum height=60pt,minimum width=20pt, fill=white, label=above:{$v$}] (e) at (9,0) {};

    % Aristas
    \draw (9) -- (10);
    \end{tikzpicture}
    \caption{$f_{0,R}(C)=1$ y $f_{0,R}(v)=0$}
    \label{SnyR=VIH(B)}
    \end{subfigure}
    \caption{Comportamiento del $VIH$ en $S_n$ con $n\geq3$ y $R=\lceil \frac{n+1}{2} \rceil$, primeros casos triviales.}
    \label{SnyR=vih}
    \end{figure}
    Ahora supongamos una configuraci\'on inicial en la que el centro $C$ est\'a sano $f_{0,R}(C)=0$, y de los $n$ v\'ertices exteriores, supongamos que hay un conjunto de estos v\'ertices que lo denotaremos como $B$, en donde, $|B|=n-\lceil n/2\rceil$ est\'an infectados y el otro conjunto $A$ de v\'ertices exteriores $|A|=\lceil n/2\rceil$ est\'an sanos. En el tiempo inicial $t=0$, el centro tiene $d_{0,1}(C)=|B|\ge 1$, por lo que en el siguiente paso se infecta $f_{1,R}(C)=1$, y los exteriores infectados mueren $f_{1,R}(B)=2$, mientras que los v\'ertices del conjunto $A$ permanecen sanos $f_{1,R}(A)=0$. En el tiempo $t=1$, los exteriores muertos tienen $d_{1,1}(B)=1<R$, por lo que pasan a sanos en el tiempo $t=2$ $f_{2,R}(B)=0$, el centro pasa a muerto $f_{2,R}(C)=2$ y como este estaba infectado en el tiempo anterior, entonces $f_{2,R}(A)=1$. 
    
    Para el siguiente tiempo $t=3$, primero vamos a ver ciertas condiciones que hay que cumplir, como las siguientes desigualdades $n > \lceil \frac{n+1}{2} \rceil \geq \lceil \frac{n}{2} \rceil$, adem\'as $\lceil \frac{n}{2} \rceil \geq n - \lceil \frac{n}{2} \rceil$, pero $\lceil \frac{n+1}{2} \rceil > n - \lceil \frac{n}{2} \rceil$. Por lo tanto vamos a tener dos casos que hay que considerar para pasar el tercer tiempo.

    En nuestro primer caso, como nuestro conjunto de v\'ertices externos $A$, esta infectado y adem\'as s\'i $|A|= \lceil \frac{n}{2} \rceil = \lceil \frac{n+1}{2} \rceil=R$. En este caso, para nuestro tiempo $t=3$, vamos a tener que $C$ est\'a infectado $f_{3,R}(C)=1$, mientras que los v\'ertices externos est\'an $f_{3,R}(A)=2$ y $f_{3,R}(B)=0$. De aqui es f\'acil ver que para el siguiente tiempo $t=4$, tendremos que $C$ estar\'a muerto $f_{4,R}(C)=2$, por lo tanto, nuestro conjunto $B$ que estaba sano en el paso anterior, entonces pasa a infectarse $f_{4,R}(B)=1$, mientras que el conjunto $A$ como solo tenia al centro como vecino infectado y este no es suficiente para que todos los v\'ertices de este conjunto se vuelvan a infectar, entonces se curan $f_{4,R}(A)=0$. Ahora recordemos que tenemos la condici\'on de que $\lceil \frac{n}{2} \rceil \geq n - \lceil \frac{n}{2} \rceil$, por ende, nuevamente tenemos dos situaciones que hay que considerar para el siguiente tiempo. 

    En el primer subcaso, si la cardinalidad de nuestro conjunto de v\'ertices externos $B$ es $|B|= n-\lceil \frac{n}{2} \rceil < \lceil \frac{n}{2} \rceil$, es decir $|B|<|A|$. En esta situaci\'on es muy sencillo, ya que para el tiempo $t=5$, tendremos que $C$ estar\'a sano $f_{5,R}(C)=0$, mientras que en los externos $f_{5,R}(A)=0$ y $f_{5,R}(B)=2$. Por lo tanto para el siguiente tiempo $t=6$, todos los v\'ertices de la gr\'afica de estrella estar\'an sanos, poni\'endole fin a la enfermedad, es comportamiento se puede observar en la Figura \ref{SnyR=VIH,|A|=R(1)}. El segundo subcaso, si  la cardinalidad de nuestro conjunto de v\'ertices externos $B$ es $|B|= n-\lceil \frac{n}{2} \rceil = \lceil \frac{n}{2} \rceil$, entonces, no existe ning\'un caso en el que sean iguales, ya que si $n$ es par, es decir $n=2k$ para alg\'un $k$ en los enteros, entonces, $\lceil \frac{n}{2} \rceil = \frac{n}{2} = k$ y $n - \lceil \frac{n}{2} \rceil = n - \frac{n}{2} = \frac{n}{2} = k$. Por lo tanto en $\lceil \frac{n+1}{2} \rceil = \lceil \frac{2k+1}{2} \rceil = \lceil \frac{2k}{2} \rceil + \lceil \frac{1}{2} \rceil = k+1$, es decir que $\lceil \frac{n+1}{2} \rceil = \frac{n}{2} + 1$.

    \begin{figure}[h]
    \centering
    \begin{subfigure}[h]{1\linewidth}
    \begin{tikzpicture}
    %tiempo 0
    % vertices
    \node[draw, circle, fill=white, label=above:{$C$}] (1) at (0,0) {};
    \node[draw, ellipse, minimum height=60pt,minimum width=20pt, fill=white, label=above:{$A$}] (2) at (-1,0) {};
    \node[draw, ellipse, minimum height=60pt,minimum width=20pt, fill=red, label={[label distance=1pt]253:$t=0$}] (3) at (1,0) {};
    \node[draw, ellipse, minimum height=60pt,minimum width=20pt, fill=red, label=above:{$B$}] (a) at (1,0) {};

    % Aristas
    \draw (1) -- (2);
    \draw (1) -- (3);

    %tiempo 1
    % vertices
    \node[draw, circle, fill=red, label=above:{$C$}] (4) at (4,0) {};
    \node[draw, ellipse, minimum height=60pt,minimum width=20pt, fill=white, label=above:{$A$}] (5) at (3,0) {};
    \node[draw, ellipse, minimum height=60pt,minimum width=20pt, fill=black, label={[label distance=1pt]253:$t=1$}] (6) at (5,0) {};
    \node[draw, ellipse, minimum height=60pt,minimum width=20pt, fill=black, label=above:{$B$}] (b) at (5,0) {};

    % Aristas
    \draw (4) -- (5);
    \draw (4) -- (6);
    
    %tiempo 2
    % vertices
    \node[draw, circle, fill=black, label=above:{$C$}] (7) at (8,0) {};
    \node[draw, ellipse, minimum height=60pt,minimum width=20pt, fill=red, label=above:{$A$}] (8) at (7,0) {};
    \node[draw, ellipse, minimum height=60pt,minimum width=20pt, fill=white, label={[label distance=1pt]253:$t=2$}] (9) at (9,0) {};
    \node[draw, ellipse, minimum height=60pt,minimum width=20pt, fill=white, label=above:{$B$}] (c) at (9,0) {};

    % Aristas
    \draw (7) -- (8);
    \draw (7) -- (9);

    %tiempo 3
    % vertices
    \node[draw, circle, fill=red, label=above:{$C$}] (10) at (12,0) {};
    \node[draw, ellipse, minimum height=60pt,minimum width=20pt, fill=black, label=above:{$A$}] (11) at (11,0) {};
    \node[draw, ellipse, minimum height=60pt,minimum width=20pt, fill=white, label={[label distance=1pt]253:$t=3$}] (12) at (13,0) {};
    \node[draw, ellipse, minimum height=60pt,minimum width=20pt, fill=white, label=above:{$B$}] (d) at (13,0) {};

    % Aristas
    \draw (10) -- (11);
    \draw (10) -- (12);
    \end{tikzpicture}
    \end{subfigure}
    
    \begin{subfigure}[h]{0.7\linewidth}
    \begin{tikzpicture}
    %tiempo 4
    % vertices
    \node[draw, circle, fill=black, label=above:{$C$}] (1) at (0,0) {};
    \node[draw, ellipse, minimum height=60pt,minimum width=20pt, fill=white, label=above:{$A$}] (2) at (-1,0) {};
    \node[draw, ellipse, minimum height=60pt,minimum width=20pt, fill=red, label={[label distance=1pt]253:$t=4$}] (3) at (1,0) {};
    \node[draw, ellipse, minimum height=60pt,minimum width=20pt, fill=red, label=above:{$B$}] (a) at (1,0) {};

    % Aristas
    \draw (1) -- (2);
    \draw (1) -- (3);
    
    %tiempo 5
    % vertices
    \node[draw, circle, fill=white, label=above:{$C$}] (4) at (4,0) {};
    \node[draw, ellipse, minimum height=60pt,minimum width=20pt, fill=white, label=above:{$A$}] (5) at (3,0) {};
    \node[draw, ellipse, minimum height=60pt,minimum width=20pt, fill=black, label={[label distance=1pt]253:$t=5$}] (6) at (5,0) {};
    \node[draw, ellipse, minimum height=60pt,minimum width=20pt, fill=black, label=above:{$B$}] (b) at (5,0) {};

    % Aristas
    \draw (4) -- (5);
    \draw (4) -- (6);

    %tiempo 6
    % vertices
    \node[draw, circle, fill=white, label=above:{$C$}] (7) at (8,0) {};
    \node[draw, ellipse, minimum height=60pt,minimum width=20pt, fill=white, label=above:{$A$}] (8) at (7,0) {};
    \node[draw, ellipse, minimum height=60pt,minimum width=20pt, fill=white, label={[label distance=1pt]253:$t=6$}] (9) at (9,0) {};
    \node[draw, ellipse, minimum height=60pt,minimum width=20pt, fill=white, label=above:{$B$}] (c) at (9,0) {};

    % Aristas
    \draw (7) -- (8);
    \draw (7) -- (9);
    \end{tikzpicture}
    \end{subfigure}
    \caption{Comportamiento del $VIH$ en $S_n$ con $n\geq3$ y $R=\lceil \frac{n+1}{2} \rceil$, s\'i $|A|=R$ y $|B|<|A|$, primer caso.}
    \label{SnyR=VIH,|A|=R(1)}
    \end{figure}
    Ahora, si la cardinalidad de nuestro conjunto de v\'ertices externos $A$ que estaba infectado en el tiempo $t=2$ es $|A|= \lceil \frac{n}{2} \rceil < \lceil \frac{n+1}{2} \rceil$, es decir, $|A|<R$. Si partimos de esto, entonces para el tiempo $t=3$, $C$ estar\'a sano $f_{3,R}(C)=0$, mientras que $f_{3,R}(A)=2$ y $f_{3,R}(B)=0$. Por ende de aqui es f\'acil ver que para el siguiente tiempo $t=4$ todos v\'ertices estar\'an sanos, poni\'endole fin a la enfermedad, como lo podemos observar en la Figura \ref{SnyR=VIH,|A|<R(1)}.
    
    \begin{figure}[h!]
    \centering
    \begin{subfigure}[h]{0.8\linewidth}
    \begin{tikzpicture}
    %tiempo 0
    % vertices
    \node[draw, circle, fill=white, label=above:{$C$}] (1) at (0,0) {};
    \node[draw, ellipse, minimum height=60pt,minimum width=20pt, fill=white, label=above:{$A$}] (2) at (-1,0) {};
    \node[draw, ellipse, minimum height=60pt,minimum width=20pt, fill=red, label={[label distance=1pt]253:$t=0$}] (3) at (1,0) {};
    \node[draw, ellipse, minimum height=60pt,minimum width=20pt, fill=red, label=above:{$B$}] (a) at (1,0) {};

    % Aristas
    \draw (1) -- (2);
    \draw (1) -- (3);

    %tiempo 1
    % vertices
    \node[draw, circle, fill=red, label=above:{$C$}] (4) at (4,0) {};
    \node[draw, ellipse, minimum height=60pt,minimum width=20pt, fill=white, label=above:{$A$}] (5) at (3,0) {};
    \node[draw, ellipse, minimum height=60pt,minimum width=20pt, fill=black, label={[label distance=1pt]253:$t=1$}] (6) at (5,0) {};
    \node[draw, ellipse, minimum height=60pt,minimum width=20pt, fill=black, label=above:{$B$}] (b) at (5,0) {};

    % Aristas
    \draw (4) -- (5);
    \draw (4) -- (6);
    
    %tiempo 2
    % vertices
    \node[draw, circle, fill=black, label=above:{$C$}] (7) at (8,0) {};
    \node[draw, ellipse, minimum height=60pt,minimum width=20pt, fill=red, label=above:{$A$}] (8) at (7,0) {};
    \node[draw, ellipse, minimum height=60pt,minimum width=20pt, fill=white, label={[label distance=1pt]253:$t=2$}] (9) at (9,0) {};
    \node[draw, ellipse, minimum height=60pt,minimum width=20pt, fill=white, label=above:{$B$}] (c) at (9,0) {};

    % Aristas
    \draw (7) -- (8);
    \draw (7) -- (9);
    \end{tikzpicture}
    \end{subfigure}
    
    \begin{subfigure}[h]{0.5\linewidth}
    \begin{tikzpicture}
    %tiempo 3
    % vertices
    \node[draw, circle, fill=white, label=above:{$C$}] (1) at (0,0) {};
    \node[draw, ellipse, minimum height=60pt,minimum width=20pt, fill=black, label=above:{$A$}] (2) at (-1,0) {};
    \node[draw, ellipse, minimum height=60pt,minimum width=20pt, fill=white, label={[label distance=1pt]253:$t=3$}] (3) at (1,0) {};
    \node[draw, ellipse, minimum height=60pt,minimum width=20pt, fill=white, label=above:{$B$}] (a) at (1,0) {};

    % Aristas
    \draw (1) -- (2);
    \draw (1) -- (3);
    
    %tiempo 4
    % vertices
    \node[draw, circle, fill=white, label=above:{$C$}] (4) at (4,0) {};
    \node[draw, ellipse, minimum height=60pt,minimum width=20pt, fill=white, label=above:{$A$}] (5) at (3,0) {};
    \node[draw, ellipse, minimum height=60pt,minimum width=20pt, fill=white, label={[label distance=1pt]253:$t=4$}] (6) at (5,0) {};
    \node[draw, ellipse, minimum height=60pt,minimum width=20pt, fill=white, label=above:{$B$}] (b) at (5,0) {};

    % Aristas
    \draw (4) -- (5);
    \draw (4) -- (6);
    \end{tikzpicture}
    \end{subfigure}
    \caption{Comportamiento del $VIH$ en $S_n$ con $n\geq3$ y $R=\lceil \frac{n+1}{2} \rceil$, s\'i $|A|<R$, primer caso.}
    \label{SnyR=VIH,|A|<R(1)}
    \end{figure}
    Finalmente vamos a considerar una \'ultima configuraci\'on inicial en la que el centro $C$ est\'a infectado $f_{0,R}(C)=1$, y de los $n$ v\'ertices exteriores, supongamos que hay un conjunto de estos v\'ertices que lo denotaremos como $B$, en donde, $|B|=n-\lceil n/2\rceil$ est\'an infectados y el otro conjunto $A$ de v\'ertices exteriores $|A|=\lceil n/2\rceil$ est\'an sanos. Con esto es f\'acil ver que para el tiempo $t=1$, $C$ va a estar muerto $f_{1,R}(C)=2$, mientras que $f_{1,R}(A)=1$ y $f_{1,R}(B)=2$. Ahora, para el siguiente tiempo $t=2$, vamos a considerar las condiciones $n > \lceil \frac{n+1}{2} \rceil \geq \lceil \frac{n}{2} \rceil$, adem\'as $\lceil \frac{n}{2} \rceil \geq n - \lceil \frac{n}{2} \rceil$, pero $\lceil \frac{n+1}{2} \rceil > n - \lceil \frac{n}{2} \rceil$. Por lo tanto para esta configuraci\'on inicial, nuevamente vamos a tener varios casos que hay que considerar para pasar al segundo tiempo.

    El primer caso es s\'i la cardinalidad del conjunto de v\'ertices externos $A$ es $|A|= \lceil \frac{n}{2} \rceil = \lceil \frac{n+1}{2} \rceil$, es decir, $|A|=R$. Como $R \geq \lceil \frac{n+1}{2} \rceil$, entonces para el tiempo $t=2$, vamos a tener que $C$ va a estar infectado $f_{2,R}(C)=1$, mientras que en los v\'ertices externos $f_{2,R}(A)=2$ y $f_{2,R}(B)=0$. Por lo tanto para el siguiente tiempo $t=3$, tendremos que $C$ est\'a muerto $f_{3,R}(C)=2$, $f_{3,R}(A)=0$ y $f_{3,R}(B)=1$. Recordando las condiciones iniciales antes mencionadas, sabemos que $\lceil \frac{n}{2} \rceil \geq n - \lceil \frac{n}{2} \rceil$, por ende, nuevamente tenemos dos situaciones que hay que considerar para el siguiente tiempo. 

    El primer subcaso, tenemos que s\'i la cardinalidad de nuestro conjunto de v\'ertices externos $B$ es $|B|= n-\lceil \frac{n}{2} \rceil < \lceil \frac{n}{2} \rceil$, en otras palabras, $|B|<|A|$. En este caso es f\'acil ver que para el tiempo $t=4$, $C$ estar\'a sano $f_{4,R}(C)=0$, mientras que en los v\'ertices externos $f_{4,R}(A)=0$ y $f_{4,R}(B)=2$. Finalmente para el siguiente tiempo $t=5$, todos los v\'ertices de la gr\'afica estar\'an sanos, llegando al fin de la enfermedad, este comportamiento se puede visualizar mejor en la Figura \ref{SnyR=VIH,|A|=R(2)}. En nuestro segundo subcaso, si nuestro conjunto de v\'ertices externos $B$ es $|B|= n-\lceil \frac{n}{2} \rceil = \lceil \frac{n}{2} \rceil=|A|$, este caso es an\'alogo al caso que vimos anteriormente, es decir, no existe ning\'un caso en el que sean iguales, ya que si $n$ es par, es decir $n=2k$ para alg\'un $k$ en los enteros, entonces, $\lceil \frac{n}{2} \rceil = \frac{n}{2} = k$ y $n - \lceil \frac{n}{2} \rceil = n - \frac{n}{2} = \frac{n}{2} = k$. Por lo tanto en $\lceil \frac{n+1}{2} \rceil = \lceil \frac{2k+1}{2} \rceil = \lceil \frac{2k}{2} \rceil + \lceil \frac{1}{2} \rceil = k+1$, en otros t\'erminos, $\lceil \frac{n+1}{2} \rceil = \frac{n}{2} + 1$, y con esto, contemplamos todas las posibles situaciones que generaba nuestro primer caso. 

    En nuestro segundo y \'ultlimo caso, tenemos que s\'i la cardinalidad de nuestro conjunto de v\'ertices externos $A$ es $|A|= \lceil \frac{n}{2} \rceil < \lceil \frac{n+1}{2} \rceil=R$, con esto es sencillo ver que para el tiempo $t=2$, vamos a tener que $C$ est\'a sano $f_{2,R}(C)=0$, mientras que los v\'ertices externos $f_{2,R}(A)=2$ y $f_{2,R}(B)=0$. De aqui es f\'acil notar que para nuestro siguiente tiempo $t=3$, todos los v\'ertices de nuestra gr\'afica estar\'an sanos y por ende la enfermedad se extingui\'o, esta din\'amica de propagaci\'on se puede observar en la Figura \ref{SnyR=VIH,|A|<R(2)}.
    \begin{figure}[h]
    \centering
    \begin{subfigure}[h]{0.8\linewidth}
    \begin{tikzpicture}
    %tiempo 0
    % vertices
    \node[draw, circle, fill=red, label=above:{$C$}] (1) at (0,0) {};
    \node[draw, ellipse, minimum height=60pt,minimum width=20pt, fill=white, label=above:{$A$}] (2) at (-1,0) {};
    \node[draw, ellipse, minimum height=60pt,minimum width=20pt, fill=red, label={[label distance=1pt]253:$t=0$}] (3) at (1,0) {};
    \node[draw, ellipse, minimum height=60pt,minimum width=20pt, fill=red, label=above:{$B$}] (a) at (1,0) {};

    % Aristas
    \draw (1) -- (2);
    \draw (1) -- (3);

    %tiempo 1
    % vertices
    \node[draw, circle, fill=black, label=above:{$C$}] (4) at (4,0) {};
    \node[draw, ellipse, minimum height=60pt,minimum width=20pt, fill=red, label=above:{$A$}] (5) at (3,0) {};
    \node[draw, ellipse, minimum height=60pt,minimum width=20pt, fill=black, label={[label distance=1pt]253:$t=1$}] (6) at (5,0) {};
    \node[draw, ellipse, minimum height=60pt,minimum width=20pt, fill=black, label=above:{$B$}] (b) at (5,0) {};

    % Aristas
    \draw (4) -- (5);
    \draw (4) -- (6);
    
    %tiempo 2
    % vertices
    \node[draw, circle, fill=red, label=above:{$C$}] (7) at (8,0) {};
    \node[draw, ellipse, minimum height=60pt,minimum width=20pt, fill=black, label=above:{$A$}] (8) at (7,0) {};
    \node[draw, ellipse, minimum height=60pt,minimum width=20pt, fill=white, label={[label distance=1pt]253:$t=2$}] (9) at (9,0) {};
    \node[draw, ellipse, minimum height=60pt,minimum width=20pt, fill=white, label=above:{$B$}] (c) at (9,0) {};

    % Aristas
    \draw (7) -- (8);
    \draw (7) -- (9);
    \end{tikzpicture}
    \end{subfigure}
    
    \begin{subfigure}[h]{0.8\linewidth}
    \begin{tikzpicture}
    %tiempo 3
    % vertices
    \node[draw, circle, fill=black, label=above:{$C$}] (1) at (0,0) {};
    \node[draw, ellipse, minimum height=60pt,minimum width=20pt, fill=white, label=above:{$A$}] (2) at (-1,0) {};
    \node[draw, ellipse, minimum height=60pt,minimum width=20pt, fill=red, label={[label distance=1pt]253:$t=3$}] (3) at (1,0) {};
    \node[draw, ellipse, minimum height=60pt,minimum width=20pt, fill=red, label=above:{$B$}] (a) at (1,0) {};

    % Aristas
    \draw (1) -- (2);
    \draw (1) -- (3);

    %tiempo 4
    % vertices
    \node[draw, circle, fill=white, label=above:{$C$}] (4) at (4,0) {};
    \node[draw, ellipse, minimum height=60pt,minimum width=20pt, fill=white, label=above:{$A$}] (5) at (3,0) {};
    \node[draw, ellipse, minimum height=60pt,minimum width=20pt, fill=black, label={[label distance=1pt]253:$t=4$}] (6) at (5,0) {};
    \node[draw, ellipse, minimum height=60pt,minimum width=20pt, fill=black, label=above:{$B$}] (b) at (5,0) {};

    % Aristas
    \draw (4) -- (5);
    \draw (4) -- (6);
    
    %tiempo 5
    % vertices
    \node[draw, circle, fill=white, label=above:{$C$}] (7) at (8,0) {};
    \node[draw, ellipse, minimum height=60pt,minimum width=20pt, fill=white, label=above:{$A$}] (8) at (7,0) {};
    \node[draw, ellipse, minimum height=60pt,minimum width=20pt, fill=white, label={[label distance=1pt]253:$t=5$}] (9) at (9,0) {};
    \node[draw, ellipse, minimum height=60pt,minimum width=20pt, fill=white, label=above:{$B$}] (c) at (9,0) {};

    % Aristas
    \draw (7) -- (8);
    \draw (7) -- (9);
    \end{tikzpicture}
    \end{subfigure}
    \caption{Comportamiento del $VIH$ en $S_n$ con $n\geq3$ y $R=\lceil \frac{n+1}{2} \rceil$, s\'i $|A|=R$ y $|B|<|A|$, segundo caso.}
    \label{SnyR=VIH,|A|=R(2)}
    \end{figure}

    \begin{figure}[h]
    \centering
    \begin{subfigure}[h]{1\linewidth}
    \begin{tikzpicture}
    %tiempo 0
    % vertices
    \node[draw, circle, fill=red, label=above:{$C$}] (1) at (0,0) {};
    \node[draw, ellipse, minimum height=60pt,minimum width=20pt, fill=white, label=above:{$A$}] (2) at (-1,0) {};
    \node[draw, ellipse, minimum height=60pt,minimum width=20pt, fill=red, label={[label distance=1pt]253:$t=0$}] (3) at (1,0) {};
    \node[draw, ellipse, minimum height=60pt,minimum width=20pt, fill=red, label=above:{$B$}] (a) at (1,0) {};

    % Aristas
    \draw (1) -- (2);
    \draw (1) -- (3);

    %tiempo 1
    % vertices
    \node[draw, circle, fill=black, label=above:{$C$}] (4) at (4,0) {};
    \node[draw, ellipse, minimum height=60pt,minimum width=20pt, fill=red, label=above:{$A$}] (5) at (3,0) {};
    \node[draw, ellipse, minimum height=60pt,minimum width=20pt, fill=black, label={[label distance=1pt]253:$t=1$}] (6) at (5,0) {};
    \node[draw, ellipse, minimum height=60pt,minimum width=20pt, fill=black, label=above:{$B$}] (b) at (5,0) {};

    % Aristas
    \draw (4) -- (5);
    \draw (4) -- (6);
    
    %tiempo 2
    % vertices
    \node[draw, circle, fill=white, label=above:{$C$}] (7) at (8,0) {};
    \node[draw, ellipse, minimum height=60pt,minimum width=20pt, fill=black, label=above:{$A$}] (8) at (7,0) {};
    \node[draw, ellipse, minimum height=60pt,minimum width=20pt, fill=white, label={[label distance=1pt]253:$t=2$}] (9) at (9,0) {};
    \node[draw, ellipse, minimum height=60pt,minimum width=20pt, fill=white, label=above:{$B$}] (c) at (9,0) {};

    % Aristas
    \draw (7) -- (8);
    \draw (7) -- (9);

    %tiempo 3
    % vertices
    \node[draw, circle, fill=white, label=above:{$C$}] (10) at (12,0) {};
    \node[draw, ellipse, minimum height=60pt,minimum width=20pt, fill=white, label=above:{$A$}] (11) at (11,0) {};
    \node[draw, ellipse, minimum height=60pt,minimum width=20pt, fill=white, label={[label distance=1pt]253:$t=3$}] (12) at (13,0) {};
    \node[draw, ellipse, minimum height=60pt,minimum width=20pt, fill=white, label=above:{$B$}] (d) at (13,0) {};

    % Aristas
    \draw (10) -- (11);
    \draw (10) -- (12);
    \end{tikzpicture}
    \end{subfigure}
    \caption{Comportamiento del $VIH$ en $S_n$ con $n\geq3$ y $R=\lceil \frac{n+1}{2} \rceil$, s\'i $|A|<R$, segundo caso.}
    \label{SnyR=VIH,|A|<R(2)}
    \end{figure}

    Finalmente, hemos demostrado que si $R<\lceil (n+1)/2\rceil$, la infecci\'on puede mantenerse de forma indefinida, y si $R\ge \lceil (n+1)/2\rceil$, el sistema siempre se cura completamente. 
\end{proof}

\subsection{VIH en trayectorias $P_n$}
La familia de las trayectorias $P_n$, tienen una estructura lineal que permite un an\'alisis m\'as directo del proceso de curaci\'on.

\begin{figure}[h!]
    \centering
    \begin{tabular}{|c|cccccccccc|}
        \hline
        \textbf{n} & \textbf{1} & \textbf{2} & \textbf{3} & \textbf{4} & \textbf{5} & \textbf{6} & \textbf{7} & \textbf{8} & \textbf{9} & \textbf{10}\\
        \hline
        \textbf{$VIH(P_n)$} & 1 & 2 & 2 & 2 & 2 & 2 & 2 & 2 & 2 & 2 \\
        \hline
    \end{tabular}
    \caption{Tabla $VIH(P_n)$}
    \label{tabla:VIH(P_n)}
\end{figure}
En este caso, los resultados computacionales vistos en la Figura \ref{tabla:VIH(P_n)} indicaron que valores relativamente peque\~nos del par\'ametro $R$ bastan para garantizar la recuperaci\'on total, lo que sugiere que la propagaci\'on de la infecci\'on se ve limitada por la baja conectividad de los v\'ertices terminales.
A partir de estas observaciones se formula el siguiente teorema:

\begin{theorem}\label{Teo:VIH(Pn)}
    $$
    VIH(P_n) = \left\lbrace
    \begin{array}{ll}
        1 \textup{ si } n = 1\\
        2 \textup{ e.o.c}
    \end{array}
    \right.
    $$
\end{theorem}

\begin{proof}
    Sea $P_n$ una trayectoria simple con v\'ertices $v_1,v_2,\dots,v_n$. Recordemos que en el modelo de Mukwembi cada v\'ertice $v$ puede encontrarse en tres posibles estados a cada instante $i$, los cuales se describen mediante la funci\'on $f_{i,R}:V(G)\rightarrow \{0,1,2\}$, donde $f_{i,R}(v)=0$ indica que el sitio se encuentra sano, $f_{i,R}(v)=1$ que est\'a infectado y $f_{i,R}(v)=2$ que est\'a muerto. Si $N(v)$ representa el conjunto de vecinos de $v$, definimos $N_{i,j}(v)=\{x\in N(v)\mid f_{i,R}(x)=j\}$ y $d_{i,j}(v)=|N_{i,j}(v)|$. Las reglas de transici\'on se expresan entonces de la siguiente forma: un v\'ertice sano pasa a infectado si $d_{i,1}(v)\ge 1$, un infectado pasa a muerto en el paso siguiente, y un v\'ertice muerto se reemplaza por infectado si y solo si $d_{i,1}(v)\ge R$, o por sano en caso contrario.

    Cuando $n=1$, la trayectoria est\'a compuesta por un \'unico v\'ertice $v_1$ sin vecinos, es decir, $d_{i,1}(v_1)=0$ para todo instante $i$. Si $R=1$, entonces las posibles configuraciones iniciales son $f_{0,1}(v_1)=0$ o $f_{0,1}(v_1)=1$. En el primer caso, el v\'ertice permanecer\'a sano indefinidamente, ya que no tiene vecinos infectados que puedan modificar su estado. En el segundo caso, si el v\'ertice inicia infectado, al paso siguiente se tiene $f_{1,1}(v_1)=2$, y dado que $d_{1,1}(v_1)=0<R$, se obtiene que $f_{2,1}(v_1)=0$. En ambos escenarios, el sistema alcanza un estado completamente sano en un n\'umero finito de pasos, por lo que la infecci\'on se extingue para $R=1$. De este modo se concluye que $VIH(P_1)=1$.

    Consideremos ahora el caso $n\ge 2$. Por el Teorema \ref{Teo:Delta+1}, se sabe que si $G$ contiene al menos una arista, la condici\'on $R=1$ no garantiza la extinci\'on de la infecci\'on. Para verlo intuitivamente; sea la configuraci\'on inicial tal que $f_{0,1}(v_1)=1$ y $f_{0,1}(x)=0$ para todo $x\neq v_1$. En el siguiente paso, el v\'ertice $v_1$ muere, pues estaba infectado, de modo que $f_{1,1}(v_1)=2$, mientras que su vecino inmediato $v_2$ se infecta, es decir, $f_{1,1}(v_2)=1$, dado que $d_{0,1}(v_2)=1$. En el instante posterior, $v_1$ vuelve a infectarse porque $d_{1,1}(v_1)=1\ge R$ y simult\'aneamente $v_2$ muere, de forma que $f_{2,1}(v_1)=1$ y $f_{2,1}(v_2)=2$. Este comportamiento se repite indefinidamente, generando un ciclo donde los estados de $v_1$ y $v_2$ alternan entre infectado y muerto sin alcanzar nunca el estado sano de manera simult\'anea. En consecuencia, $R=1$ no es suficiente para extinguir la infecci\'on en trayectorias con $n\ge 2$, por lo que se concluye que $VIH(P_n)\ge 2$, este comportamiento es parecido al ejemplo de la Figura \ref{S2yR=1}, ya que una estrella $S_1$ tiene la misma estructura que una trayectoria $P_2$.

    Por el Teorema \ref{Teo:Delta+1}, sabemos que $VIH(G)\leq \Delta(G)+1$. En una trayectoria, $\Delta(P_n)=2$, por lo que $VIH(P_n)\leq3$; sin embargo, mostraremos que basta con $R=2$.
    
    Procedamos a analizar el caso $R=2$. Supongamos una configuraci\'on inicial arbitraria $f_{0,2}$. Los v\'ertices extremos $v_1$ y $v_n$ de la trayectoria tienen grado 1, por lo que en cualquier instante $i$ se cumple que $d_{i,1}(v_1)\in\{0,1\}$ y de igual manera para $v_n$. Si en alg\'un instante $t$ se cumple que $f_{t,2}(v_1)=2$, entonces $d_{t,1}(v_1)\le 1 < R$, lo cual implica que $f_{t+1,2}(v_1)=0$.
    
    Esto significa que un extremo muerto siempre es reemplazado por un v\'ertice sano. Adem\'as, una vez que $v_1$ se encuentra en estado sano, no puede volver a infectarse en instantes posteriores, ya que para que un v\'ertice sano pase al estado infectado es necesario que tenga al menos un vecino infectado. Sin embargo, como $v_1$ tiene grado 1, su \'unico vecino debe encontrarse infectado en un instante determinado para que ocurra una reinfecci\'on. Si esto sucede, entonces $v_1$ pasa a infectado en el siguiente paso y necesariamente muere en el paso posterior. No obstante, al morir se cumple nuevamente que $d_{i,1}(v_1)\le 1<R$, por lo que el v\'ertice es reemplazado por uno sano. En consecuencia, cualquier posible reinfecci\'on del extremo es transitoria y no puede sostenerse en el tiempo. Por lo tanto, los v\'ertices extremos no pueden formar parte de un ciclo de reinfecci\'on y, una vez alcanzado el estado sano, permanecen definitivamente sanos.De esta manera, los extremos de la trayectoria act\'uan como barreras para la dina\'amica, ya que no pueden sostener infecci\'on persistente ni contribuir a la reinfecci\'on del resto de la gra\'afica.

    A partir de este comportamiento, la curaci\'on se propaga desde los extremos hacia el interior de la trayectoria. Los v\'ertices adyacentes a los extremos, al perder contacto con v\'ertices infectados, entran sucesivamente en las fases de infecci\'on, muerte y posterior sanaci\'on, hasta estabilizarse como sanos. Este proceso se repite progresivamente en direcci\'on al centro de la trayectoria, reduciendo paso a paso la regi\'on de v\'ertices infectados o muertos. Dado que en una trayectoria no existen ciclos que permitan mantener una fuente de infecci\'on constante, el proceso termina necesariamente en un estado completamente sano tras un n\'umero finito de iteraciones, como se muestra en la Figura \ref{PnyR=2}.
    \begin{figure}[h]
        \centering
        \begin{tikzpicture}
            %tiempo 0
            % vertices
            \node[draw, circle, fill=white, label=left:{$t=0$}] (1) at (0,0) {};
            \node[draw, circle, fill=red, label=above:{$ $}] (2) at (1,0) {};
            \node[draw, circle, fill=white, label=above:{$ $}] (3) at (2,0) {};
            \node[draw, circle, fill=red, label=above:{$ $}] (4) at (3,0) {};
            \node[draw, circle, fill=white, label=above:{$ $}] (5) at (4,0) {};
            
            % Aristas
            \draw (1) -- (2);
            \draw (2) -- (3);
            \draw (3) -- (4);
            \draw (4) -- (5);

            %tiempo 1
            % vertices
            \node[draw, circle, fill=red, label=left:{$t=1$}] (1) at (0,-1) {};
            \node[draw, circle, fill=black, label=above:{$ $}] (2) at (1,-1) {};
            \node[draw, circle, fill=red, label=above:{$ $}] (3) at (2,-1) {};
            \node[draw, circle, fill=black, label=above:{$ $}] (4) at (3,-1) {};
            \node[draw, circle, fill=red, label=above:{$ $}] (5) at (4,-1) {};
            
            % Aristas
            \draw (1) -- (2);
            \draw (2) -- (3);
            \draw (3) -- (4);
            \draw (4) -- (5);

            %tiempo 2
            % vertices
            \node[draw, circle, fill=black, label=left:{$t=2$}] (1) at (0,-2) {};
            \node[draw, circle, fill=red, label=above:{$ $}] (2) at (1,-2) {};
            \node[draw, circle, fill=black, label=above:{$ $}] (3) at (2,-2) {};
            \node[draw, circle, fill=red, label=above:{$ $}] (4) at (3,-2) {};
            \node[draw, circle, fill=black, label=above:{$ $}] (5) at (4,-2) {};
            
            % Aristas
            \draw (1) -- (2);
            \draw (2) -- (3);
            \draw (3) -- (4);
            \draw (4) -- (5);

            %tiempo 3
            % vertices
            \node[draw, circle, fill=white, label=left:{$t=3$}] (1) at (6,0) {};
            \node[draw, circle, fill=black, label=above:{$ $}] (2) at (7,0) {};
            \node[draw, circle, fill=red, label=above:{$ $}] (3) at (8,0) {};
            \node[draw, circle, fill=black, label=above:{$ $}] (4) at (9,0) {};
            \node[draw, circle, fill=white, label=above:{$ $}] (5) at (10,0) {};
            
            % Aristas
            \draw (1) -- (2);
            \draw (2) -- (3);
            \draw (3) -- (4);
            \draw (4) -- (5);

            %tiempo 4
            % vertices
            \node[draw, circle, fill=white, label=left:{$t=4$}] (1) at (6,-1) {};
            \node[draw, circle, fill=white, label=above:{$ $}] (2) at (7,-1) {};
            \node[draw, circle, fill=black, label=above:{$ $}] (3) at (8,-1) {};
            \node[draw, circle, fill=white, label=above:{$ $}] (4) at (9,-1) {};
            \node[draw, circle, fill=white, label=above:{$ $}] (5) at (10,-1) {};
            
            % Aristas
            \draw (1) -- (2);
            \draw (2) -- (3);
            \draw (3) -- (4);
            \draw (4) -- (5);

            %tiempo 5
            % vertices
            \node[draw, circle, fill=white, label=left:{$t=5$}] (1) at (6,-2) {};
            \node[draw, circle, fill=white, label=above:{$ $}] (2) at (7,-2) {};
            \node[draw, circle, fill=white, label=above:{$ $}] (3) at (8,-2) {};
            \node[draw, circle, fill=white, label=above:{$ $}] (4) at (9,-2) {};
            \node[draw, circle, fill=white, label=above:{$ $}] (5) at (10,-2) {};
            
            % Aristas
            \draw (1) -- (2);
            \draw (2) -- (3);
            \draw (3) -- (4);
            \draw (4) -- (5);
        \end{tikzpicture}
        \caption{Comportamiento de $P_5$ con $R=2$.}
        \label{PnyR=2}
    \end{figure}
    De esta forma, para $R=2$, la infecci\'on siempre se extingue en $P_n$, independientemente de la configuraci\'on inicial. Dado que se ha mostrado que $R=1$ no es suficiente y que $R=2$ garantiza la extinci\'on, concluimos que $VIH(P_n)=2$ para toda trayectoria con $n\ge 2$.
\end{proof}
\subsection{VIH en gr\'aficas completas $K_n$}
En el caso de las gr\'aficas completas $K_n$, donde cada v\'ertice se conecta con todos los dem\'as, gracias a las resultados que se muestran en la Figura \ref{tabla:VIH(K_n)}, se observ\'o el valor m\'as alto de $VIH(G)$ entre todas las familias estudiadas. Este resultado es consistente con la interpretaci\'on del par\'ametro $R$: en una gr\'afica completamente conectada, cualquier v\'ertice recuperado tiene el m\'aximo n\'umero posible de vecinos infectados, por lo que requiere una resistencia m\'as elevada para evitar la re infecci\'on.

\begin{figure}[h!]
    \centering
    \begin{tabular}{|c|cccccccccc|}
        \hline
        \textbf{n} & \textbf{1} & \textbf{2} & \textbf{3} & \textbf{4} & \textbf{5} & \textbf{6} & \textbf{7} & \textbf{8} & \textbf{9} & \textbf{10}\\
        \hline
        \textbf{$VIH(K_n)$} & 1 & 2 & 2 & 3 & 3 & 4 & 4 & 5 & 5 & 6 \\
        \hline
    \end{tabular}
    \caption{Tabla $VIH(K_n)$}
    \label{tabla:VIH(K_n)}
\end{figure}
De este modo, el n\'umero $VIH(K_n)$ puede interpretarse como una funci\'on creciente del orden $n$, que refleja la densidad m\'axima de conexiones de la gr\'afica.
\begin{theorem}\label{Teo:VIH(Kn)}
    $$
    VIH(K_n) = \lceil \frac{n+1}{2} \rceil
    $$
\end{theorem}

\begin{proof}

    Comencemos analizando el caso trivial. Cuando $n=1$, la gr\'afica completa $K_1$ consta de un solo v\'ertice $v$ sin vecinos, es decir, $d_{i,1}(v)=0$ para todo instante $i$. Si tomamos $R=1$, entonces las posibles configuraciones iniciales son $f_{0,1}(v)=0$ o $f_{0,1}(v)=1$. 
    
    En el primer caso, el v\'ertice permanecer\'a sano indefinidamente, ya que al no tener vecinos infectados no existen condiciones que modifiquen su estado, por lo que $f_{i,1}(v)=0$ para todo $i\geq 0$. En el segundo caso, si el v\'ertice inicia infectado, se tiene que $f_{1,1}(v)=2$ (por la regla de que un v\'ertice infectado muere en el siguiente instante). Adem\'as, dado que $d_{1,1}(v)=0<R$, se cumple que $f_{2,1}(v)=0$. Por lo tanto, en ambos escenarios el sistema alcanza un estado completamente sano en un n\'umero finito de pasos, lo que implica que la infecci\'on se extingue para $R=1$. De este modo se concluye que $VIH(K_1)=1$.

    Consideremos ahora el caso $n=2$ y $R=1$. La gr\'afica completa $K_2$ est\'a formada por dos v\'ertices $v_1$ y $v_2$ que son adyacentes entre s\'i, por lo que es equivalente a una trayectoria $P_2$ o tambi\'en a una estrella $S_1$. De acuerdo con los resultados previos para estas estructuras, si $R=1$ la infecci\'on no se extingue, pues los v\'ertices alternan sus estados entre infectado y muerto de forma indefinida. 
    
    En efecto, si se parte de una configuraci\'on $f_{0,1}$ donde $f_{0,1}(v_1)=1$ y $f_{0,1}(v_2)=0$, al paso siguiente se tiene $f_{1,1}(v_1)=2$ y $f_{1,1}(v_2)=1$. En el instante $2$, ocurre que $f_{2,1}(v_1)=1$ y $f_{2,1}(v_2)=2$, reproduci\'endose as\'i un ciclo de infecci\'on sin convergencia a un estado sano, tal como se observa en la Figura \ref{S1yR=1}. 
    
    En cambio, si se toma $R=2$, los v\'ertices muertos no alcanzan el umbral de reinfecci\'on, ya que $d_{i,1}(v)<R$ en todo instante $i$, y por lo tanto el sistema converge a un estado completamente sano en un n\'umero finito de pasos, como se muestra en la Figura \ref{S1yR=2}. De este modo, $R=2$ es suficiente para garantizar la desaparici\'on de la infecci\'on en $K_2$.

    Prosigamos con el caso general donde $n\geq 3$. Sea $R=\lceil \tfrac{n+1}{2}\rceil-1$. Particionemos el conjunto de v\'ertices $V(K_n)$ en dos subconjuntos disjuntos $A$ y $B$, tales que $|A|=\lceil \tfrac{n}{2}\rceil$ y $|B|=n-\lceil \tfrac{n}{2}\rceil$. Si $n$ es impar, se cumple que $|A|>|B|$, de modo que $|A|=\lceil \tfrac{n+1}{2}\rceil$, mientras que si $n$ es par se tiene $|A|=|B|=\lceil \tfrac{n+1}{2}\rceil-1$.
    
    Consideremos una configuraci\'on inicial $f_{0,R}$ tal que los v\'ertices de $A$ est\'an sanos, es decir, $f_{0,R}(A)=0$, y los v\'ertices de $B$ est\'an infectados, de modo que $f_{0,R}(B)=1$. Como $K_n$ es completa, cada v\'ertice de $A$ est\'a conectado con todos los de $B$ y viceversa, por lo que en el siguiente paso se obtiene $f_{1,R}(A)=1$ y $f_{1,R}(B)=2$.

    Ahora bien, cada v\'ertice muerto en $B$ tiene exactamente $d_{1,1}(B)=|A|$ vecinos infectados, y dado que $|A|\ge R$, estos v\'ertices se re-infectan en el siguiente instante, es decir, $f_{2,R}(B)=1$. De forma an\'aloga, los v\'ertices infectados de $A$ mueren, de modo que $f_{2,R}(A)=2$. Este proceso se repite indefinidamente, alternando los estados de infecci\'on y muerte entre los conjuntos $A$ y $B$ como se muestra en la Figura \ref{KnyR<VIH}, lo que impide que la infecci\'on se extinga. 
    
    Por lo tanto, cuando $R=\lceil \tfrac{n+1}{2}\rceil-1$, existe una configuraci\'on inicial que mantiene la infecci\'on activa para todo tiempo $i$, concluy\'endose que en este caso $R$ no es suficiente para curar el sistema.

    \begin{figure}[h!]
    \centering
    \begin{tikzpicture}
    %tiempo 0
    % vertices
    \node[draw, ellipse, minimum height=60pt,minimum width=20pt, fill=white, label=above:{$A$}] (1) at (-1,0) {};
    \node[draw, ellipse, minimum height=60pt,minimum width=20pt, fill=red, label={[label distance=1pt]258:$t=0$}] (2) at (0.5,0) {};
    \node[draw, ellipse, minimum height=60pt,minimum width=20pt, fill=red, label=above:{$B$}] (a) at (0.5,0) {};

    % Aristas
    \draw (1) -- (2);

    %tiempo 1
    % vertices
    \node[draw, ellipse, minimum height=60pt,minimum width=20pt, fill=red, label=above:{$A$}] (3) at (2,0) {};
    \node[draw, ellipse, minimum height=60pt,minimum width=20pt, fill=black, label={[label distance=1pt]258:$t=1$}] (4) at (3.5,0) {};
    \node[draw, ellipse, minimum height=60pt,minimum width=20pt, fill=black, label=above:{$B$}] (b) at (3.5,0) {};

    % Aristas
    \draw (3) -- (4);
    
    %tiempo 2
    % vertices
    \node[draw, ellipse, minimum height=60pt,minimum width=20pt, fill=black, label=above:{$A$}] (5) at (5,0) {};
    \node[draw, ellipse, minimum height=60pt,minimum width=20pt, fill=red, label={[label distance=1pt]258:$t=2$}] (6) at (6.5,0) {};
    \node[draw, ellipse, minimum height=60pt,minimum width=20pt, fill=red, label=above:{$B$}] (c) at (6.5,0) {};

    % Aristas
    \draw (5) -- (6);

    %tiempo 3
    % vertices
    \node[draw, ellipse, minimum height=60pt,minimum width=20pt, fill=red, label=above:{$A$}] (7) at (8,0) {};
    \node[draw, ellipse, minimum height=60pt,minimum width=20pt, fill=black, label={[label distance=1pt]258:$t=3$}] (8) at (9.5,0) {};
    \node[draw, ellipse, minimum height=60pt,minimum width=20pt, fill=black, label=above:{$B$}] (d) at (9.5,0) {};

    % Aristas
    \draw (7) -- (8);
    \end{tikzpicture}
    \caption{Comportamiento del $VIH$ en $K_n$ con $n\geq3$ y $R=\lceil \frac{n+1}{2} \rceil-1$}
    \label{KnyR<VIH}
    \end{figure}
    Por otro lado, si se incrementa el par\'ametro a $R=\lceil \tfrac{n+1}{2}\rceil$, el n\'umero de vecinos infectados de cualquier v\'ertice muerto satisface $d_{i,1}(v)<R$ para todo instante $i$, lo cual impide que un v\'ertice en estado muerto pueda re-infectarse, como se observa en la Figura \ref{KnyR=VIH}. 
    
    En consecuencia, cualquier configuraci\'on inicial $f_{0,R}$ conduce en un n\'umero finito de pasos a un estado en el que todos los v\'ertices son sanos, es decir, $f_{t,R}(v)=0$ para todo $v\in V(K_n)$ y para todo $t$ suficientemente grande. Por lo tanto, el valor cr\'itico del par\'ametro de reemplazo para el cual la infecci\'on deja de propagarse y el sistema se cura completamente es $R=\lceil \tfrac{n+1}{2}\rceil$.

    \begin{figure}[h]
    \centering
    \begin{subfigure}[h]{\linewidth}
    \begin{tikzpicture}
    %tiempo 0
    % vertices
    \node[draw, ellipse, minimum height=60pt,minimum width=20pt, fill=white, label=above:{$A$}] (1) at (-1,0) {};
    \node[draw, ellipse, minimum height=60pt,minimum width=20pt, fill=red, label={[label distance=1pt]258:$t=0$}] (2) at (0.5,0) {};
    \node[draw, ellipse, minimum height=60pt,minimum width=20pt, fill=red, label=above:{$B$}] (a) at (0.5,0) {};

    % Aristas
    \draw (1) -- (2);

    %tiempo 1
    % vertices
    \node[draw, ellipse, minimum height=60pt,minimum width=20pt, fill=red, label=above:{$A$}] (3) at (2,0) {};
    \node[draw, ellipse, minimum height=60pt,minimum width=20pt, fill=black, label={[label distance=1pt]258:$t=1$}] (4) at (3.5,0) {};
    \node[draw, ellipse, minimum height=60pt,minimum width=20pt, fill=black, label=above:{$B$}] (b) at (3.5,0) {};

    % Aristas
    \draw (3) -- (4);
    
    %tiempo 2
    % vertices
    \node[draw, ellipse, minimum height=60pt,minimum width=20pt, fill=black, label=above:{$A$}] (5) at (5,0) {};
    \node[draw, ellipse, minimum height=60pt,minimum width=20pt, fill=red, label={[label distance=1pt]258:$t=2$}] (6) at (6.5,0) {};
    \node[draw, ellipse, minimum height=60pt,minimum width=20pt, fill=red, label=above:{$B$}] (c) at (6.5,0) {};

    % Aristas
    \draw (5) -- (6);

    %tiempo 3
    % vertices
    \node[draw, ellipse, minimum height=60pt,minimum width=20pt, fill=white, label=above:{$A$}] (7) at (8,0) {};
    \node[draw, ellipse, minimum height=60pt,minimum width=20pt, fill=black, label={[label distance=1pt]258:$t=3$}] (8) at (9.5,0) {};
    \node[draw, ellipse, minimum height=60pt,minimum width=20pt, fill=black, label=above:{$B$}] (d) at (9.5,0) {};

    % Aristas
    \draw (7) -- (8);

    %tiempo 4
    % vertices
    \node[draw, ellipse, minimum height=60pt,minimum width=20pt, fill=white, label=above:{$A$}] (9) at (11,0) {};
    \node[draw, ellipse, minimum height=60pt,minimum width=20pt, fill=white, label={[label distance=1pt]258:$t=4$}] (10) at (12.5,0) {};
    \node[draw, ellipse, minimum height=60pt,minimum width=20pt, fill=white, label=above:{$B$}] (e) at (12.5,0) {};

    % Aristas
    \draw (11) -- (10);
    \end{tikzpicture}
    \caption{Si $n$ es impar}
    \label{KnyR=VIH,nimpar}
    \end{subfigure}
    \begin{subfigure}[h]{0.78\linewidth}
    \begin{tikzpicture}
    %tiempo 0
    % vertices
    \node[draw, ellipse, minimum height=60pt,minimum width=20pt, fill=white, label=above:{$A$}] (1) at (-1,0) {};
    \node[draw, ellipse, minimum height=60pt,minimum width=20pt, fill=red, label={[label distance=1pt]258:$t=0$}] (2) at (0.5,0) {};
    \node[draw, ellipse, minimum height=60pt,minimum width=20pt, fill=red, label=above:{$B$}] (a) at (0.5,0) {};

    % Aristas
    \draw (1) -- (2);

    %tiempo 1
    % vertices
    \node[draw, ellipse, minimum height=60pt,minimum width=20pt, fill=red, label=above:{$A$}] (3) at (2,0) {};
    \node[draw, ellipse, minimum height=60pt,minimum width=20pt, fill=black, label={[label distance=1pt]258:$t=1$}] (4) at (3.5,0) {};
    \node[draw, ellipse, minimum height=60pt,minimum width=20pt, fill=black, label=above:{$B$}] (b) at (3.5,0) {};

    % Aristas
    \draw (3) -- (4);
    
    %tiempo 2
    % vertices
    \node[draw, ellipse, minimum height=60pt,minimum width=20pt, fill=black, label=above:{$A$}] (5) at (5,0) {};
    \node[draw, ellipse, minimum height=60pt,minimum width=20pt, fill=white, label={[label distance=1pt]258:$t=2$}] (6) at (6.5,0) {};
    \node[draw, ellipse, minimum height=60pt,minimum width=20pt, fill=white, label=above:{$B$}] (c) at (6.5,0) {};

    % Aristas
    \draw (5) -- (6);

    %tiempo 3
    % vertices
    \node[draw, ellipse, minimum height=60pt,minimum width=20pt, fill=white, label=above:{$A$}] (7) at (8,0) {};
    \node[draw, ellipse, minimum height=60pt,minimum width=20pt, fill=white, label={[label distance=1pt]258:$t=3$}] (8) at (9.5,0) {};
    \node[draw, ellipse, minimum height=60pt,minimum width=20pt, fill=white, label=above:{$B$}] (d) at (9.5,0) {};

    % Aristas
    \draw (7) -- (8);
    \end{tikzpicture}
    \caption{Si $n$ es par}
    \label{KnyR=VIH,npar}
    \end{subfigure}
    \caption{Comportamiento del $VIH$ en $K_n$ con $n\geq3$ y $R=\lceil \frac{n+1}{2} \rceil$}
    \label{KnyR=VIH}
    \end{figure}
\end{proof}

\subsection{VIH en ruedas $W_n$}
De manera similar, las gr\'aficas ruedas $W_n$, que combinan un ciclo con un v\'ertice central conectado a todos los dem\'as, presentan un comportamiento intermedio entre los ciclos y las completas. El v\'ertice central mantiene un alto grado de conectividad, pero la presencia del ciclo perif\'erico introduce una estructura que distribuye parcialmente la infecci\'on.

\begin{figure}[h!]
    \centering
    \begin{tabular}{|c|cccccccccc|}
        \hline
        \textbf{n} & \textbf{1} & \textbf{2} & \textbf{3} & \textbf{4} & \textbf{5} & \textbf{6} & \textbf{7} & \textbf{8} & \textbf{9} & \textbf{10}\\
        \hline
        \textbf{$VIH(W_n)$} & 1 & 2 & 2 & 3 & 4 & 5 & 6 & 7 & 8 & 9 \\
        \hline
    \end{tabular}
    \caption{Tabla $VIH(W_n)$}
    \label{tabla:VIH(W_n)}
\end{figure}

Los valores de $R$ obtenidos en la Figura \ref{tabla:VIH(W_n)}, muestran un equilibrio entre la densidad y la dispersi\'on de la gr\'afica.
\begin{theorem}\label{Con:VIH(Wn)}
    $$
    VIH(W_n) = \left\lbrace
    \begin{array}{ll}
        n \textup{ si } n = 1 \textup{ \'o } n=2\\
        n-1 \textup{ si } n\geq3
    \end{array}
    \right.
    $$
\end{theorem}
\begin{proof}
    Sea $W_n$ la rueda con $n$ v\'ertices, es decir, la gr\'afica formada por un ciclo $C_{n-1}$ y un v\'ertice central $C$ adyacente a todos los v\'ertices del ciclo. Empecemos con el caso trivial, cuando $n=1$, la gr\'afica $W_1$ consiste en un solo v\'ertice, y como lo hemos esta viendo anteriormente en los distintos tipos de familias de gr\'aficas, es evidente que $VIH(W_1)=1$, pues cualquier condici\'on inicial alcanza el estado sano en un paso y no existe posibilidad de reinfecci\'on al no haber vecinos. Si $n=2$, la gr\'afica $W_2$ es isomorfa a $S_1$, cuyos dos v\'ertices est\'an conectados por una sola arista. Por los resultados previamente demostrados para gr\'aficas de orden dos, o bien verificando directamente todas las configuraciones posibles, se obtiene que $VIH(W_2)=2$, ya que con $R=2$ la infecci\'on siempre se extingue en tiempo finito, como lo muestra en la Figura \ref{S1}.

    Consideremos ahora el caso general $n\ge 3$. Para el caso en el que $n=3$ observamos que $W_3$ es isomorfa a $K_3$, de donde se deduce que $VIH(K_3) = 2$, por lo tanto $VIH(W_3)=2$. 
    
    Para $n\ge4$ debemos establecer la igualdad $VIH(W_n)=n-1$ se probar\'an dos casos complementarios: primero que $VIH(W_n)\le n-1$ y luego que $VIH(W_n)\ge n-1$. Para la primera desigualdad, por el Teorema \ref{Teo:Mukwembi}, sabemos que Mukwembi demostr\'o que para cualquier gr\'afica $G$ de orden $n$, si se elige $R=n-1$, entonces la infecci\'on se extingue en tiempo finito sin importar la condici\'on inicial. En otras palabras, el valor $R=n-1$ es siempre un par\'ametro de reemplazo que termina en un estado de curaci\'on para toda gr\'afica conexa de orden $n$. Aplicando este resultado al caso particular $G=W_n$, se concluye que $R=n-1$ garantiza la extinci\'on del virus en todos los escenarios, y por tanto $VIH(W_n)\le n-1$.

    Para probar la desigualdad inversa, se construye una condici\'on inicial en la que la infecci\'on no se extingue cuando $R=n-2$. Sea $f_{0,n-2}$ definida de la siguiente manera para $W_4$: el v\'ertice central $C$ y exactamente un v\'ertice del ciclo, digamos $v_1$, comienzan infectados, es decir, $f_{0,2}(C)=f_{0,2}(v_1)=1$, mientras que todos los dem\'as v\'ertices del ciclo $A = V(W_4) \setminus \{C,v_1\}$ est\'an sanos.
    
    Al aplicar las reglas del modelo se obtiene la siguiente evoluci\'on. En el siguiente instante, $t=1$, $C$ y $v_1$ se convierten en muertos, por lo que $f_{1,2}(C)=f_{1,2}(v_1)=2$. Al mismo tiempo, todos los v\'ertices sanos adyacentes a un infectado en $t=0$ se vuelven infectados, de modo que todos los v\'ertices del ciclo $A$, al ser vecinos del centro, se infectan. En consecuencia, en $t=1$ el centro $C$ est\'a muerto y todos los v\'ertices del ciclo $A$ est\'an infectados. 
    
    En el instante siguiente, $t=2$, los v\'ertices infectados en $A$ pasan a muertos. Como el centro $C$ estaba muerto en $t=1$, debemos determinar si vuelve a infectarse. En efecto, $C$ tiene $n-2=2$ vecinos, todos infectados en $t=1$, por lo que $d_{1,1}(C)=2=R$, y entonces $f_{2,2}(C)=1$, esto mismo pasa con el v\'ertice $v_1$ ya que $d_{1,1}(v_1)=2=R$ . As\'i, en $t=2$ el centro $C$ y el v\'ertice $v_1$ vuelven a infectarse mientras todos los v\'ertices del ciclo $A$ permanecen muertos. 
    
    En el tiempo $t=3$, el centro y $v_1$ pasan a muertos nuevamente, y como los v\'ertices del ciclo $A$ eran muertos en $t=2$ pero tienen $d_{2,1}(A)=2=R$, entonces todos ellos vuelven a infectarse, de manera que la configuraci\'on en $t=3$ es esencialmente la misma que la de $t=1$: $C$ y $v_1$ muertos, y $A$ infectado. De esta forma, la din\'amica se vuelve peri\'odica, alternando entre las configuraciones $(C$ y $v_1$ muertos, $A$ infectado$)$ y $(C$ y $v_1$ infectado, $A$ muerto$)$, este comportamiento lo podemos observar en la Figura \ref{W4yR=2}.

    \begin{figure}[h]
    \centering
    \begin{tikzpicture}
        % t=0
        % vertices W_4
        \node[draw, circle, fill=red, label=right:{$ $}] (1) at (0,0) {};
        \node[draw, circle, fill=red, label=above:{$ $}] (2) at (0,1) {};
        \node[draw, circle, fill=white, label={[label distance=6pt]355:$t=0$}] (3) at (-1,-1) {};
        \node[draw, circle, fill=white, label=below:{$ $}] (4) at (1,-1) {};

        % Aristas W_4
        \draw (1) -- (2);
        \draw (1) -- (3);
        \draw (1) -- (4);
        \draw (2) -- (3);
        \draw (2) -- (4);
        \draw (3) -- (4);

        % t=1
        % vertices W_4
        \node[draw, circle, fill=black, label=right:{$ $}] (5) at (3,0) {};
        \node[draw, circle, fill=black, label=above:{$ $}] (6) at (3,1) {};
        \node[draw, circle, fill=red, label={[label distance=6pt]355:$t=1$}] (7) at (2,-1) {};
        \node[draw, circle, fill=red, label=below:{$ $}] (8) at (4,-1) {};

        % Aristas W_4
        \draw (5) -- (6);
        \draw (5) -- (7);
        \draw (5) -- (8);
        \draw (6) -- (7);
        \draw (6) -- (8);
        \draw (7) -- (8);

        % t=2
        % vertices W_4
        \node[draw, circle, fill=red, label=right:{$ $}] (9) at (6,0) {};
        \node[draw, circle, fill=red, label=above:{$ $}] (10) at (6,1) {};
        \node[draw, circle, fill=black, label={[label distance=6pt]355:$t=2$}] (11) at (5,-1) {};
        \node[draw, circle, fill=black, label=below:{$ $}] (12) at (7,-1) {};

        % Aristas W_4
        \draw (9) -- (10);
        \draw (9) -- (11);
        \draw (9) -- (12);
        \draw (10) -- (11);
        \draw (10) -- (12);
        \draw (11) -- (12);

        % t=3
        % vertices W_4
        \node[draw, circle, fill=black, label=right:{$ $}] (13) at (9,0) {};
        \node[draw, circle, fill=black, label=above:{$ $}] (14) at (9,1) {};
        \node[draw, circle, fill=red, label={[label distance=6pt]355:$t=3$}] (15) at (8,-1) {};
        \node[draw, circle, fill=red, label=below:{$ $}] (16) at (10,-1) {};

        % Aristas W_4
        \draw (13) -- (14);
        \draw (13) -- (15);
        \draw (13) -- (16);
        \draw (14) -- (15);
        \draw (14) -- (16);
        \draw (15) -- (16);

        % t=4
        % vertices W_4
        \node[draw, circle, fill=red, label=right:{$ $}] (17) at (12,0) {};
        \node[draw, circle, fill=red, label=above:{$ $}] (18) at (12,1) {};
        \node[draw, circle, fill=black, label={[label distance=6pt]355:$t=4$}] (19) at (11,-1) {};
        \node[draw, circle, fill=black, label=below:{$ $}] (20) at (13,-1) {};

        % Aristas W_4
        \draw (17) -- (18);
        \draw (17) -- (19);
        \draw (17) -- (20);
        \draw (18) -- (19);
        \draw (18) -- (20);
        \draw (19) -- (20); 
    \end{tikzpicture}
    \caption{Comportamiento del VIH en $W_4$ con $R=n-2=2$.}
    \label{W4yR=2}
    \end{figure}

    Para fortalecer el argumento de que $R=n-2$ no es un suficiente para que la enfermedad se cure en $W_n$, es conveniente analizar expl\'icitamente los primeros casos no triviales, pues muestran con claridad el mecanismo que impide la extinci\'on. En la construcci\'on utilizada previamente, el centro $C$ y uno de los v\'ertices del ciclo, que denotamos por $v_1$, comienzan infectados y todos los dem\'as v\'ertices del ciclo est\'an sanos. Para $W_4$ observamos que la evoluci\'on entra en un ciclo peri\'odico en el cual alternan las configuraciones donde el centro est\'a muerto y el ciclo est'a infectado, con aquellas donde el centro est\'a infectado y el ciclo est\'a muerto. Sin embargo, es natural preguntarse si este comportamiento persiste cuando aumenta el tama\~no de la rueda.

    En la rueda $W_5$, el ciclo tiene cuatro v\'ertices, que denotaremos por $v_1,v_2,v_3,v_4$ en orden circular. El centro se denotar\'a por $C$. El par\'ametro de reemplazo es $R=n-2=3$, en la configuraci\'on inicial $f_{0,3}$ coloca al centro y a $v_1$ infectados, y a los restantes v\'ertices del ciclo sanos, la din\'amica de propagaci\'on de este ejemplo se podr\'a observar en la Figura \ref{W5yR=3}. 

    En el instante $t=1$, tanto $C$ como $v_1$ pasan del estado infectado al estado muerto. Al mismo tiempo, cada uno de los v\'ertices sanos del ciclo observa que $d_{0,1}(v)=1$ debido al centro infectado en $t=0$, por lo que todos ellos pasan al estado infectado. De este modo en el instante $t=1$ se tiene que $f_{1,3}(C)=2$, $f_{1,3}(v_1)=2$, y $f_{1,3}(v_2)=f_{1,3}(v_3)=f_{1,3}(v_4)=1$.

    En el instante $t=2$, los v\'ertices infectados del ciclo, es decir, $v_2,v_3,v_4$, pasan a estado muerto. El centro estaba muerto en $t=1$, y como tiene $d_{1,1}(C)=3$ vecinos infectados, y este n\'umero satisface $d_{1,1}(C)\geq R$, entonces se vuelve a infectar, $f_{2,3}(C)=1$. El v\'ertice $v_1$, que estaba muerto desde $t=1$, no cumple $d_{1,1}(v_1)\geq R$, por lo que pasa a sano, es decir, $f_{2,3}(v_1)=0$. En $t=2$ tenemos entonces al centro infectado y a todos los v\'ertices del ciclo muertos excepto $v_1$, que ha vuelto al estado sano.

    En el instante $t=3$, el centro infectado pasa a muerto nuevamente, pues $f_{2,3}(C)=1$ implica $f_{3,3}(C)=2$. Todos los v\'ertices del ciclo que estaban muertos en $t=2$ como cada uno de ellos tienen solo un vecino infectado, se cumple $d_{2,1}(v)\leq R$, y por tanto vuelven al estado sano. El v\'ertice $v_1$, que estaba sano en $t=2$, pasa a infectado porque $d_{2,1}(v_1)\geq 1$. En consecuencia, en $t=3$ todos los v\'ertices del ciclo est\'an sanos excepto $v_1$ que esta infectado y el centro est\'a muerto.

    En el instante $t=4$ el v\'ertice $v_1$ que estaba infectado vuelve a muerto, y como el centro estaba muerto en $t=3$ pero tiene $d_{3,1}(C)=1\leq 3$, vuelve al estado sano. Mientras que en los v\'ertices ciclo que estaban sanos en $t=3$, como $v_2$ y $v_4$ son adyacentes a $v_1$, entonces $d_{3,1}(v) \geq 1$ por lo tanto se infectan, sin embargo el v\'ertice $v_3$ en el tiempo $t=3$ no cuenta con vecinos infectados que alteren su comportamiento, por ende permanecer\'a sano para el tiempo $t=4$.

    En $t=5$, el centro y el v\'ertice $v_3$ pasan a infectarse ya que cumplen con el criterio de $d_{4,1}(v)\geq 1$, en cambio los v\'ertices $v_2$ y $v_4$ pasan a muertos, mientras que $v_1$ como en el tiempo anterior estaba muerto y $d_{4,1}(v_1) < R$, entonces sanara en $t=5$. En el siguiente tiempo $t=6$, el centro y el v\'ertice $v_3$ pasan a muertos, los v\'ertices $v_2$ y $v_4$ pasan a sanos ya que cumplen con la condici\'on de $d_{5,1}(v) < R$, mientras que $v_1$ como en el tiempo anterior estaba susceptible y $d_{5,1}(v_1) \geq 1$, entonces se infectara en $t=6$.
    
    Finalmente en el instante $t=7$ los v\'ertices $C$ y $v_3$ que estaban muertos pasan a sanos ya que $d_{6,1}(v) < R$, los v\'ertices $v_2$ y $v_4$ se infectan por $d_{6,1}(v)\geq 1$, mientras que $v_1$ muere, notamos que el comportamiento de este tiempo est\'a replicando la configuraci\'on de $t=4$. Esta repetici\'on exacta marca el inicio de un ciclo peri\'odico, decimos que la din\'amica entra en un ciclo de periodo tres, pues $t=4$ reproduce la configuraci\'on de $t=7$ y esta a su vez coincidir\'a con la de $t=10$. De este modo la infecci\'on nunca se extingue en $W_5$ cuando $R=n-2=3$. 
    \begin{figure}[h]
    \centering
    \begin{tikzpicture}
        %t=0
        % vertices W_5
        \node[draw, circle, fill=red, label=right:{$ $}] (1) at (1,0) {};
        \node[draw, circle, fill=white, label=above:{$ $}] (2) at (2,1) {};
        \node[draw, circle, fill=red, label=above:{$ $}] (3) at (0,1) {};
        \node[draw, circle, fill=white, label={[label distance=6pt]355:$t=0$}] (4) at (0,-1) {};
        \node[draw, circle, fill=white, label=below:{$ $}] (5) at (2,-1) {};

        % Aristas W_5
        \draw (1) -- (2);
        \draw (1) -- (3);
        \draw (1) -- (4);
        \draw (1) -- (5);
        \draw (2) -- (3);
        \draw (3) -- (4);
        \draw (4) -- (5);
        \draw (2) -- (5);
        
        %t=1
        % vertices W_5
        \node[draw, circle, fill=black, label=right:{$ $}] (1) at (4,0) {};
        \node[draw, circle, fill=red, label=above:{$ $}] (2) at (5,1) {};
        \node[draw, circle, fill=black, label=above:{$ $}] (3) at (3,1) {};
        \node[draw, circle, fill=red, label={[label distance=6pt]355:$t=1$}] (4) at (3,-1) {};
        \node[draw, circle, fill=red, label=below:{$ $}] (5) at (5,-1) {};

        % Aristas W_5
        \draw (1) -- (2);
        \draw (1) -- (3);
        \draw (1) -- (4);
        \draw (1) -- (5);
        \draw (2) -- (3);
        \draw (3) -- (4);
        \draw (4) -- (5);
        \draw (2) -- (5);

        %t=2
        % vertices W_5
        \node[draw, circle, fill=red, label=right:{$ $}] (1) at (7,0) {};
        \node[draw, circle, fill=black, label=above:{$ $}] (2) at (8,1) {};
        \node[draw, circle, fill=white, label=above:{$ $}] (3) at (6,1) {};
        \node[draw, circle, fill=black, label={[label distance=6pt]355:$t=2$}] (4) at (6,-1) {};
        \node[draw, circle, fill=black, label=below:{$ $}] (5) at (8,-1) {};

        % Aristas W_5
        \draw (1) -- (2);
        \draw (1) -- (3);
        \draw (1) -- (4);
        \draw (1) -- (5);
        \draw (2) -- (3);
        \draw (3) -- (4);
        \draw (4) -- (5);
        \draw (2) -- (5);

        %t=3
        % vertices W_5
        \node[draw, circle, fill=black, label=right:{$ $}] (1) at (10,0) {};
        \node[draw, circle, fill=white, label=above:{$ $}] (2) at (11,1) {};
        \node[draw, circle, fill=red, label=above:{$ $}] (3) at (9,1) {};
        \node[draw, circle, fill=white, label={[label distance=6pt]355:$t=3$}] (4) at (9,-1) {};
        \node[draw, circle, fill=white, label=below:{$ $}] (5) at (11,-1) {};

        % Aristas W_5
        \draw (1) -- (2);
        \draw (1) -- (3);
        \draw (1) -- (4);
        \draw (1) -- (5);
        \draw (2) -- (3);
        \draw (3) -- (4);
        \draw (4) -- (5);
        \draw (2) -- (5);

        %t=4
        % vertices W_5
        \node[draw, circle, fill=white, label=right:{$ $}] (1) at (1,-3) {};
        \node[draw, circle, fill=red, label=above:{$ $}] (2) at (2,-2) {};
        \node[draw, circle, fill=black, label=above:{$ $}] (3) at (0,-2) {};
        \node[draw, circle, fill=red, label={[label distance=6pt]355:$t=4$}] (4) at (0,-4) {};
        \node[draw, circle, fill=white, label=below:{$ $}] (5) at (2,-4) {};

        % Aristas W_5
        \draw (1) -- (2);
        \draw (1) -- (3);
        \draw (1) -- (4);
        \draw (1) -- (5);
        \draw (2) -- (3);
        \draw (3) -- (4);
        \draw (4) -- (5);
        \draw (2) -- (5);

        %t=5
        % vertices W_5
        \node[draw, circle, fill=red, label=right:{$ $}] (1) at (4,-3) {};
        \node[draw, circle, fill=black, label=above:{$ $}] (2) at (5,-2) {};
        \node[draw, circle, fill=white, label=above:{$ $}] (3) at (3,-2) {};
        \node[draw, circle, fill=black, label={[label distance=6pt]355:$t=5$}] (4) at (3,-4) {};
        \node[draw, circle, fill=red, label=below:{$ $}] (5) at (5,-4) {};

        % Aristas W_5
        \draw (1) -- (2);
        \draw (1) -- (3);
        \draw (1) -- (4);
        \draw (1) -- (5);
        \draw (2) -- (3);
        \draw (3) -- (4);
        \draw (4) -- (5);
        \draw (2) -- (5);

        %t=6
        % vertices W_5
        \node[draw, circle, fill=black, label=right:{$ $}] (1) at (7,-3) {};
        \node[draw, circle, fill=white, label=above:{$ $}] (2) at (8,-2) {};
        \node[draw, circle, fill=red, label=above:{$ $}] (3) at (6,-2) {};
        \node[draw, circle, fill=white, label={[label distance=6pt]355:$t=6$}] (4) at (6,-4) {};
        \node[draw, circle, fill=black, label=below:{$ $}] (5) at (8,-4) {};

        % Aristas W_5
        \draw (1) -- (2);
        \draw (1) -- (3);
        \draw (1) -- (4);
        \draw (1) -- (5);
        \draw (2) -- (3);
        \draw (3) -- (4);
        \draw (4) -- (5);
        \draw (2) -- (5);

        %t=7
        % vertices W_5
        \node[draw, circle, fill=white, label=right:{$ $}] (1) at (10,-3) {};
        \node[draw, circle, fill=red, label=above:{$ $}] (2) at (11,-2) {};
        \node[draw, circle, fill=black, label=above:{$ $}] (3) at (9,-2) {};
        \node[draw, circle, fill=red, label={[label distance=6pt]355:$t=7$}] (4) at (9,-4) {};
        \node[draw, circle, fill=white, label=below:{$ $}] (5) at (11,-4) {};

        % Aristas W_5
        \draw (1) -- (2);
        \draw (1) -- (3);
        \draw (1) -- (4);
        \draw (1) -- (5);
        \draw (2) -- (3);
        \draw (3) -- (4);
        \draw (4) -- (5);
        \draw (2) -- (5);
    \end{tikzpicture}   
    \caption{Comportamiento del VIH en $W_5$ con $R=n-2=3$.}
    \label{W5yR=3}
    \end{figure}

    En el caso de $W_6$ el ciclo tiene cinco v\'ertices, que se nombran $v_1,v_2,v_3,v_4,v_5$ en orden circular; el centro se denota nuevamente por $C$. El valor del par\'ametro es $R=n-2=4$, la configuraci\'on inicial fija $f_{0,4}(C)=1$ y $f_{0,4}(v_1)=1$, mientras que el resto de los v\'ertices del ciclo est\'an sanos, la din\'amica de propgaci\'on de este ejemplo se explicara a continuaci\'on y se podr\'a visualizar en la Figura \ref{W6yR=4}.

    En el instante $t=1$, tanto el centro como $v_1$ pasan al estado muerto. Como cada uno de los v\'ertices $v_2,v_3,v_4,v_5$ tiene $d_{0,1}(v)\geq1$ debido al centro infectado en el instante previo, todos ellos pasan a estado infectado. Por tanto, $f_{1,4}(C)=2$, $f_{1,4}(v_1)=2$ y $f_{1,4}(v_2)=f_{1,4}(v_3)=f_{1,4}(v_4)=f_{1,4}(v_5)=1$.

    En el instante $t=2$, todos los v\'ertices del ciclo que estaban infectados pasan a muertos. Mientras, el centro como en $t=1$ tenia $d_{1,1}(C)=4\ge R$, el centro vuelve al estado infectado en el instante $t=2$. El v\'ertice $v_1$, que estaba muerto en $t=1$, no cumple $d_{1,1}(v_1)\ge R$, y por ello vuelve a sano.

    En $t=3$, el centro infectado en $t=2$ pasa al estado muerto, mientras que todos los v\'ertices del ciclo, que estaban muertos en $t=2$ pero ten\'ian al centro como vecino infectado, vuelven al estado sano porque $d_{2,1}(v)< R$, en cambio como el v\'ertice $v_1$ estaba susceptible en el tiempo anterior y tenia un vecino que era el centro que estaba infectado, por lo tanto en el tiempo $t=3$ se volvera a infectar.

    En el instante $t=4$ el v\'ertice $v_1$ que estaba infectado vuelve a muerto, y como el centro estaba muerto en $t=3$ pero tiene $d_{3,1}(C)=1\leq 3$, vuelve al estado sano. Mientras que en los v\'ertices ciclo que estaban sanos en $t=3$, como $v_2$ y $v_5$ son adyacentes a $v_1$, entonces $d_{3,1}(v) \geq 1$ por lo tanto se infectan, sin embargo el v\'ertice $v_3$ y $v_4$ en el tiempo $t=3$ no cuentan con vecinos infectados que alteren su comportamiento, por ende permanecer\'an sano para el tiempo $t=4$.

    En $t=5$, el centro y los v\'ertices $v_3$ y $v_4$ pasan a infectarse ya que cumplen con el criterio de $d_{4,1}(v)\geq 1$, en cambio los v\'ertices $v_2$ y $v_5$ pasan a muertos, mientras que $v_1$ como en el tiempo anterior estaba muerto y $d_{4,1}(v_1) < R$, entonces sanara en $t=5$. En el siguiente tiempo $t=6$, nuevamente el centro y los v\'ertices $v_3$ y $v_4$ pasan a muertos, los v\'ertices $v_2$ y $v_5$ pasan a sanos ya que cumplen con la condici\'on de $d_{5,1}(v) < R$, mientras que $v_1$ como en el tiempo anterior estaba susceptible y $d_{5,1}(v_1) \geq 1$, entonces se infectara en $t=6$.
    
    Finalmente en el instante $t=7$ los v\'ertices $C$, $v_3$ y $v_4$ que estaban muertos pasan a sanos ya que $d_{6,1}(v) < R$, los v\'ertices $v_2$ y $v_5$ se infectan por $d_{6,1}(v)\geq 1$, mientras que $v_1$ muere, notamos que el comportamiento de este tiempo est\'a replicando la configuraci\'on de $t=4$. Esto marca el inicio de un ciclo peri\'odico, decimos que la din\'amica entra en un ciclo de periodo tres, pues $t=4$ reproduce la configuraci\'on de $t=7$ y esta a su vez coincidir\'a con la de $t=10$. De este modo la infecci\'on nunca se extingue en $W_6$ cuando $R=n-2=4$.

    \begin{figure}[h]
    \centering
    \begin{tikzpicture}
        %t=0
        % vertices W_6
        \node[draw, circle, fill=red, label={[label distance=28pt]270:$t=0$}] (1) at (1,0) {};
        \node[draw, circle, fill=red, label=above:{$ $}] (2) at (1,1) {};
        \node[draw, circle, fill=white, label=above:{$ $}] (3) at (1.9,0.2) {};
        \node[draw, circle, fill=white, label=below:{$ $}] (4) at (1.5,-1) {};
        \node[draw, circle, fill=white, label=below:{$ $}] (5) at (0.5,-1) {};
        \node[draw, circle, fill=white, label=below:{$ $}] (6) at (0.1,0.2) {};
        
        % Aristas W_6
        \draw (1) -- (2);
        \draw (1) -- (3);
        \draw (1) -- (4);
        \draw (1) -- (5);
        \draw (1) -- (6);
        \draw (2) -- (3);
        \draw (3) -- (4);
        \draw (4) -- (5);
        \draw (5) -- (6);
        \draw (6) -- (2);

        %t=1
        % vertices W_6
        \node[draw, circle, fill=black, label={[label distance=28pt]270:$t=1$}] (1) at (3.8,0) {};
        \node[draw, circle, fill=black, label=above:{$ $}] (2) at (3.8,1) {};
        \node[draw, circle, fill=red, label=above:{$ $}] (3) at (4.7,0.2) {};
        \node[draw, circle, fill=red, label=below:{$ $}] (4) at (4.3,-1) {};
        \node[draw, circle, fill=red, label=below:{$ $}] (5) at (3.3,-1) {};
        \node[draw, circle, fill=red, label=below:{$ $}] (6) at (2.9,0.2) {};
        
        % Aristas W_6
        \draw (1) -- (2);
        \draw (1) -- (3);
        \draw (1) -- (4);
        \draw (1) -- (5);
        \draw (1) -- (6);
        \draw (2) -- (3);
        \draw (3) -- (4);
        \draw (4) -- (5);
        \draw (5) -- (6);
        \draw (6) -- (2);

        %t=2
        % vertices W_6
        \node[draw, circle, fill=red, label={[label distance=28pt]270:$t=2$}] (1) at (6.6,0) {};
        \node[draw, circle, fill=white, label=above:{$ $}] (2) at (6.6,1) {};
        \node[draw, circle, fill=black, label=above:{$ $}] (3) at (7.5,0.2) {};
        \node[draw, circle, fill=black, label=below:{$ $}] (4) at (7.1,-1) {};
        \node[draw, circle, fill=black, label=below:{$ $}] (5) at (6.1,-1) {};
        \node[draw, circle, fill=black, label=below:{$ $}] (6) at (5.7,0.2) {};
        
        % Aristas W_6
        \draw (1) -- (2);
        \draw (1) -- (3);
        \draw (1) -- (4);
        \draw (1) -- (5);
        \draw (1) -- (6);
        \draw (2) -- (3);
        \draw (3) -- (4);
        \draw (4) -- (5);
        \draw (5) -- (6);
        \draw (6) -- (2);

        %t=3
        % vertices W_6
        \node[draw, circle, fill=black, label={[label distance=28pt]270:$t=3$}] (1) at (9.4,0) {};
        \node[draw, circle, fill=red, label=above:{$ $}] (2) at (9.4,1) {};
        \node[draw, circle, fill=white, label=above:{$ $}] (3) at (10.3,0.2) {};
        \node[draw, circle, fill=white, label=below:{$ $}] (4) at (9.9,-1) {};
        \node[draw, circle, fill=white, label=below:{$ $}] (5) at (8.9,-1) {};
        \node[draw, circle, fill=white, label=below:{$ $}] (6) at (8.5,0.2) {};
        
        % Aristas W_6
        \draw (1) -- (2);
        \draw (1) -- (3);
        \draw (1) -- (4);
        \draw (1) -- (5);
        \draw (1) -- (6);
        \draw (2) -- (3);
        \draw (3) -- (4);
        \draw (4) -- (5);
        \draw (5) -- (6);
        \draw (6) -- (2);

        %t=4 
        % vertices W_6
        \node[draw, circle, fill=white, label={[label distance=28pt]270:$t=4$}] (1) at (1,-3) {};
        \node[draw, circle, fill=black, label=above:{$ $}] (2) at (1,-2) {};
        \node[draw, circle, fill=red, label=above:{$ $}] (3) at (1.9,-2.8) {};
        \node[draw, circle, fill=white, label=below:{$ $}] (4) at (1.5,-4) {};
        \node[draw, circle, fill=white, label=below:{$ $}] (5) at (0.5,-4) {};
        \node[draw, circle, fill=red, label=below:{$ $}] (6) at (0.1,-2.8) {};
        
        % Aristas W_6
        \draw (1) -- (2);
        \draw (1) -- (3);
        \draw (1) -- (4);
        \draw (1) -- (5);
        \draw (1) -- (6);
        \draw (2) -- (3);
        \draw (3) -- (4);
        \draw (4) -- (5);
        \draw (5) -- (6);
        \draw (6) -- (2);

        %t=5 
        % vertices W_6
        \node[draw, circle, fill=red, label={[label distance=28pt]270:$t=5$}] (1) at (3.8,-3) {};
        \node[draw, circle, fill=white, label=above:{$ $}] (2) at (3.8,-2) {};
        \node[draw, circle, fill=black, label=above:{$ $}] (3) at (4.7,-2.8) {};
        \node[draw, circle, fill=red, label=below:{$ $}] (4) at (4.3,-4) {};
        \node[draw, circle, fill=red, label=below:{$ $}] (5) at (3.3,-4) {};
        \node[draw, circle, fill=black, label=below:{$ $}] (6) at (2.9,-2.8) {};
        
        % Aristas W_6
        \draw (1) -- (2);
        \draw (1) -- (3);
        \draw (1) -- (4);
        \draw (1) -- (5);
        \draw (1) -- (6);
        \draw (2) -- (3);
        \draw (3) -- (4);
        \draw (4) -- (5);
        \draw (5) -- (6);
        \draw (6) -- (2);

        %t=6 
        % vertices W_6
        \node[draw, circle, fill=black, label={[label distance=28pt]270:$t=6$}] (1) at (6.6,-3) {};
        \node[draw, circle, fill=red, label=above:{$ $}] (2) at (6.6,-2) {};
        \node[draw, circle, fill=white, label=above:{$ $}] (3) at (7.5,-2.8) {};
        \node[draw, circle, fill=black, label=below:{$ $}] (4) at (7.1,-4) {};
        \node[draw, circle, fill=black, label=below:{$ $}] (5) at (6.1,-4) {};
        \node[draw, circle, fill=white, label=below:{$ $}] (6) at (5.7,-2.8) {};
        
        % Aristas W_6
        \draw (1) -- (2);
        \draw (1) -- (3);
        \draw (1) -- (4);
        \draw (1) -- (5);
        \draw (1) -- (6);
        \draw (2) -- (3);
        \draw (3) -- (4);
        \draw (4) -- (5);
        \draw (5) -- (6);
        \draw (6) -- (2);

        %t=7 
        % vertices W_6
        \node[draw, circle, fill=white, label={[label distance=28pt]270:$t=7$}] (1) at (9.4,-3) {};
        \node[draw, circle, fill=black, label=above:{$ $}] (2) at (9.4,-2) {};
        \node[draw, circle, fill=red, label=above:{$ $}] (3) at (10.3,-2.8) {};
        \node[draw, circle, fill=white, label=below:{$ $}] (4) at (9.9,-4) {};
        \node[draw, circle, fill=white, label=below:{$ $}] (5) at (8.9,-4) {};
        \node[draw, circle, fill=red, label=below:{$ $}] (6) at (8.5,-2.8) {};
        
        % Aristas W_6
        \draw (1) -- (2);
        \draw (1) -- (3);
        \draw (1) -- (4);
        \draw (1) -- (5);
        \draw (1) -- (6);
        \draw (2) -- (3);
        \draw (3) -- (4);
        \draw (4) -- (5);
        \draw (5) -- (6);
        \draw (6) -- (2);
    \end{tikzpicture}   
    \caption{Comportamiento del VIH en $W_6$ con $R=n-2=4$.}
    \label{W6yR=4}
    \end{figure}
    
    Para la gr\'afica $W_7$, el ciclo contiene seis v\'ertices, denotados por $v_1,v_2,v_3,v_4,v_5,v_6$ en orden circular, y el centro se denota por $C$. En la configuraci\'on inicial $f_{0,5}$, tanto el centro como $v_1$ se encuentran infectados, mientras que los dem\'as v\'ertices del ciclo est\'an sanos, en la Figura \ref{W7yR=5} se podr\'a observar la din\'amica que se explicara a continuaci\'on.

    En el instante $t=1$, los v\'ertices infectados, es decir, $C$ y $v_1$, pasan inmediatamente al estado muerto. Mientras que todos los v\'ertices sanos del ciclo, como $d_{0,1}(v)\geq1$ debido al centro infectado en el instante inicial, por lo que se infectan. En consecuencia, $f_{1,5}(C)=2$, $f_{1,5}(v_1)=2$ y $f_{1,5}(v_j)=1$ para cada $j\ge 2$.

    En $t=2$, todos los v\'ertices del ciclo que estaban infectados pasan al estado muerto, es decir, $f_{2,5}(v_j)=2$ para cada $j\ge 2$. El centro estaba muerto en $t=1$, pero como cuenta con $d_{1,1}(C)=5\geq R$ v\'ertices infectados en su vecindad, la regla lo obliga a volver al estado infectado, es decir, $f_{2,5}(C)=1$. El v\'ertice $v_1$, que estaba muerto, vuelve al estado sano porque no satisface $d_{1,1}(v_1)\ge R$. De esta manera, en $t=2$ aparece el centro infectado, $v_1$ sano y el resto del ciclo muerto.

    Ahora, en $t=3$, el centro infectado pasa nuevamente al estado muerto, mientras que todos los v\'ertices del ciclo, que estaban muertos en $t=2$ pasan a sanos porque $d_{2,1}(v)= 1<R$ para cada uno de ellos, por otro lado, el v\'ertice $v_1$, como estaba susceptible en el tiempo anterior, se infectara en este tiempo ya que $d_{2,1}(v)\geq1$.

    Pasemos al instante $t=4$, la din\'amica contin\'ua el v\'ertice $v_1$ que estaba infectado vuelve a muerto, y como el centro estaba muerto en $t=3$ pero tiene $d_{3,1}(C)=1\leq R$, vuelve al estado sano. Mientras que en los v\'ertices ciclo que estaban sanos en $t=3$, como $v_2$ y $v_6$ son adyacentes a $v_1$, entonces $d_{3,1}(v) \geq 1$ por lo tanto se infectan, sin embargo el v\'ertice $v_3$ $v_4$ y $v_5$ en el tiempo $t=3$ no cuentan con vecinos infectados que alteren su comportamiento, por ende permanecer\'an sano para el tiempo $t=4$.

    En $t=5$, el centro y los v\'ertices $v_3$ y $v_5$ pasan a infectarse ya que cumplen con el criterio de $d_{4,1}(v)\geq 1$, en cambio los v\'ertices $v_2$ y $v_6$ pasan a muertos, mientras que $v_1$ como en el tiempo anterior estaba muerto y $d_{4,1}(v_1) < R$, entonces sanara en $t=5$, por otro lado el v\'ertice $v_4$ permanecer\'a sano ya que no contaba con vecinos infectados que alteran su comportamiento en este tiempo. 
    
    En el siguiente tiempo $t=6$, nuevamente el centro y los v\'ertices $v_3$ y $v_5$ pasan a muertos, los v\'ertices $v_2$ y $v_6$ pasan a sanos ya que cumplen con la condici\'on de $d_{5,1}(v) < R$, mientras que $v_1$ y $v_4$ como en el tiempo anterior estaban susceptibles y $d_{5,1}(v_1) \geq 1$, entonces se infectaran en $t=6$. Ahora, en $t=7$ los v\'ertices $C$, $v_3$ y $v_5$ que estaban muertos pasan a sanos ya que $d_{6,1}(v) < R$, los v\'ertices $v_2$ y $v_6$ se infectan por $d_{6,1}(v)\geq 1$, mientras que $v_1$ y $v_4$ mueren.

    Finalmente, en el instante $t=8$ la gr\'afica adopta exactamente la misma configuraci\'on del instante $t=5$, ya que el $C$, $v_3$ y $v_5$ se vuelven a infectar por cumplir con $d_{7,1}(v)\geq 1$, los v\'ertices $v_2$ y $v_6$ mueren, mientras que $v_1$ y $v_4$ sanan, ya que en el tiempo anterior estaban muertos y no cumpl\'ian con $d_{7,1}(v) \geq R$. Entonces, decimos que la din\'amica entra en un ciclo de periodo tres, pues $t=5$ reproduce la configuraci\'on de $t=8$ y esta a su vez coincidir\'a con la de $t=11$. De este modo la infecci\'on nunca se extingue en $W_7$ cuando $R=n-2=5$.

    Concluimos que para toda rueda $W_n$ con $n\ge 7$ el mismo mecanismo se mantiene sin variaciones esenciales. En efecto, la condici\'on inicial considerada provoca que en $t=1$ todos los v\'ertices a excepci\'on de uno del ciclo se encuentren infectados, de modo que el centro tiene $n-2$ vecinos infectados. Como $R=n-2$, el criterio de $R$ se cumple, independientemente de cu\'an grande sea la rueda. Esto implica que el centro siempre se re-infecta en $t=2$ y que toda la estructura regresa a la configuraci\'on c\'iclica descrita anteriormente, produciendo un ciclo de periodo tres que impide la extinci\'on. A partir de $W_7$ no aparece ning\'un fen\'omeno nuevo, puesto que disponer de $n-2$ v\'ertices infectados en el ciclo es suficiente para cumplir la condici\'on de reinfecci\'on del centro en todos los casos. El modelo, por tanto, exhibe un comportamiento estable e invariable para todas las ruedas grandes.

    \begin{figure}[h]
    \centering
    \begin{tikzpicture}
        %t=0
        % vertices W_6
        \node[draw, circle, fill=red, label={[label distance=28pt]270:$t=0$}] (1) at (1,0) {};
        \node[draw, circle, fill=white, label=above:{$ $}] (2) at (1.5,1) {};
        \node[draw, circle, fill=white, label=below:{$ $}] (3) at (2,0) {};
        \node[draw, circle, fill=white, label=below:{$ $}] (4) at (1.5,-1) {};
        \node[draw, circle, fill=white, label=below:{$ $}] (5) at (0.5,-1) {};
        \node[draw, circle, fill=white, label=below:{$ $}] (6) at (0,0) {};
        \node[draw, circle, fill=red, label=above:{$ $}] (7) at (0.5,1) {};
        
        % Aristas W_6
        \draw (1) -- (2);
        \draw (1) -- (3);
        \draw (1) -- (4);
        \draw (1) -- (5);
        \draw (1) -- (6);
        \draw (1) -- (7);
        \draw (2) -- (3);
        \draw (3) -- (4);
        \draw (4) -- (5);
        \draw (5) -- (6);
        \draw (6) -- (7);
        \draw (7) -- (2);

        %t=1
        % vertices W_6
        \node[draw, circle, fill=black, label={[label distance=28pt]270:$t=1$}] (1) at (4,0) {};
        \node[draw, circle, fill=red, label=above:{$ $}] (2) at (4.5,1) {};
        \node[draw, circle, fill=red, label=below:{$ $}] (3) at (5,0) {};
        \node[draw, circle, fill=red, label=below:{$ $}] (4) at (4.5,-1) {};
        \node[draw, circle, fill=red, label=below:{$ $}] (5) at (3.5,-1) {};
        \node[draw, circle, fill=red, label=below:{$ $}] (6) at (3,0) {};
        \node[draw, circle, fill=black, label=above:{$ $}] (7) at (3.5,1) {};
        
        % Aristas W_6
        \draw (1) -- (2);
        \draw (1) -- (3);
        \draw (1) -- (4);
        \draw (1) -- (5);
        \draw (1) -- (6);
        \draw (1) -- (7);
        \draw (2) -- (3);
        \draw (3) -- (4);
        \draw (4) -- (5);
        \draw (5) -- (6);
        \draw (6) -- (7);
        \draw (7) -- (2);

        %t=2
        % vertices W_6
        \node[draw, circle, fill=red, label={[label distance=28pt]270:$t=2$}] (1) at (7,0) {};
        \node[draw, circle, fill=black, label=above:{$ $}] (2) at (7.5,1) {};
        \node[draw, circle, fill=black, label=below:{$ $}] (3) at (8,0) {};
        \node[draw, circle, fill=black, label=below:{$ $}] (4) at (7.5,-1) {};
        \node[draw, circle, fill=black, label=below:{$ $}] (5) at (6.5,-1) {};
        \node[draw, circle, fill=black, label=below:{$ $}] (6) at (6,0) {};
        \node[draw, circle, fill=white, label=above:{$ $}] (7) at (6.5,1) {};
        
        % Aristas W_6
        \draw (1) -- (2);
        \draw (1) -- (3);
        \draw (1) -- (4);
        \draw (1) -- (5);
        \draw (1) -- (6);
        \draw (1) -- (7);
        \draw (2) -- (3);
        \draw (3) -- (4);
        \draw (4) -- (5);
        \draw (5) -- (6);
        \draw (6) -- (7);
        \draw (7) -- (2);

        %t=3
        % vertices W_6
        \node[draw, circle, fill=black, label={[label distance=28pt]270:$t=3$}] (1) at (10,0) {};
        \node[draw, circle, fill=white, label=above:{$ $}] (2) at (10.5,1) {};
        \node[draw, circle, fill=white, label=below:{$ $}] (3) at (11,0) {};
        \node[draw, circle, fill=white, label=below:{$ $}] (4) at (10.5,-1) {};
        \node[draw, circle, fill=white, label=below:{$ $}] (5) at (9.5,-1) {};
        \node[draw, circle, fill=white, label=below:{$ $}] (6) at (9,0) {};
        \node[draw, circle, fill=red, label=above:{$ $}] (7) at (9.5,1) {};
        
        % Aristas W_6
        \draw (1) -- (2);
        \draw (1) -- (3);
        \draw (1) -- (4);
        \draw (1) -- (5);
        \draw (1) -- (6);
        \draw (1) -- (7);
        \draw (2) -- (3);
        \draw (3) -- (4);
        \draw (4) -- (5);
        \draw (5) -- (6);
        \draw (6) -- (7);
        \draw (7) -- (2);

        %t=4
        % vertices W_6
        \node[draw, circle, fill=white, label={[label distance=28pt]270:$t=4$}] (1) at (13,0) {};
        \node[draw, circle, fill=red, label=above:{$ $}] (2) at (13.5,1) {};
        \node[draw, circle, fill=white, label=below:{$ $}] (3) at (14,0) {};
        \node[draw, circle, fill=white, label=below:{$ $}] (4) at (13.5,-1) {};
        \node[draw, circle, fill=white, label=below:{$ $}] (5) at (12.5,-1) {};
        \node[draw, circle, fill=red, label=below:{$ $}] (6) at (12,0) {};
        \node[draw, circle, fill=black, label=above:{$ $}] (7) at (12.5,1) {};
        
        % Aristas W_6
        \draw (1) -- (2);
        \draw (1) -- (3);
        \draw (1) -- (4);
        \draw (1) -- (5);
        \draw (1) -- (6);
        \draw (1) -- (7);
        \draw (2) -- (3);
        \draw (3) -- (4);
        \draw (4) -- (5);
        \draw (5) -- (6);
        \draw (6) -- (7);
        \draw (7) -- (2);

        %t=5
        % vertices W_6
        \node[draw, circle, fill=red, label={[label distance=28pt]270:$t=5$}] (1) at (2.5,-3) {};
        \node[draw, circle, fill=black, label=above:{$ $}] (2) at (3,-2) {};
        \node[draw, circle, fill=red, label=below:{$ $}] (3) at (3.5,-3) {};
        \node[draw, circle, fill=white, label=below:{$ $}] (4) at (3,-4) {};
        \node[draw, circle, fill=red, label=below:{$ $}] (5) at (2,-4) {};
        \node[draw, circle, fill=black, label=below:{$ $}] (6) at (1.5,-3) {};
        \node[draw, circle, fill=white, label=above:{$ $}] (7) at (2,-2) {};
        
        % Aristas W_6
        \draw (1) -- (2);
        \draw (1) -- (3);
        \draw (1) -- (4);
        \draw (1) -- (5);
        \draw (1) -- (6);
        \draw (1) -- (7);
        \draw (2) -- (3);
        \draw (3) -- (4);
        \draw (4) -- (5);
        \draw (5) -- (6);
        \draw (6) -- (7);
        \draw (7) -- (2);

        %t=6
        % vertices W_6
        \node[draw, circle, fill=black, label={[label distance=28pt]270:$t=6$}] (1) at (5.5,-3) {};
        \node[draw, circle, fill=white, label=above:{$ $}] (2) at (6,-2) {};
        \node[draw, circle, fill=black, label=below:{$ $}] (3) at (6.5,-3) {};
        \node[draw, circle, fill=red, label=below:{$ $}] (4) at (6,-4) {};
        \node[draw, circle, fill=black, label=below:{$ $}] (5) at (5,-4) {};
        \node[draw, circle, fill=white, label=below:{$ $}] (6) at (4.5,-3) {};
        \node[draw, circle, fill=red, label=above:{$ $}] (7) at (5,-2) {};
        
        % Aristas W_6
        \draw (1) -- (2);
        \draw (1) -- (3);
        \draw (1) -- (4);
        \draw (1) -- (5);
        \draw (1) -- (6);
        \draw (1) -- (7);
        \draw (2) -- (3);
        \draw (3) -- (4);
        \draw (4) -- (5);
        \draw (5) -- (6);
        \draw (6) -- (7);
        \draw (7) -- (2);

        %t=7
        % vertices W_6
        \node[draw, circle, fill=white, label={[label distance=28pt]270:$t=7$}] (1) at (8.5,-3) {};
        \node[draw, circle, fill=red, label=above:{$ $}] (2) at (9,-2) {};
        \node[draw, circle, fill=white, label=below:{$ $}] (3) at (9.5,-3) {};
        \node[draw, circle, fill=black, label=below:{$ $}] (4) at (9,-4) {};
        \node[draw, circle, fill=white, label=below:{$ $}] (5) at (8,-4) {};
        \node[draw, circle, fill=red, label=below:{$ $}] (6) at (7.5,-3) {};
        \node[draw, circle, fill=black, label=above:{$ $}] (7) at (8,-2) {};
        
        % Aristas W_6
        \draw (1) -- (2);
        \draw (1) -- (3);
        \draw (1) -- (4);
        \draw (1) -- (5);
        \draw (1) -- (6);
        \draw (1) -- (7);
        \draw (2) -- (3);
        \draw (3) -- (4);
        \draw (4) -- (5);
        \draw (5) -- (6);
        \draw (6) -- (7);
        \draw (7) -- (2);

        %t=8
        % vertices W_6
        \node[draw, circle, fill=red, label={[label distance=28pt]270:$t=8$}] (1) at (11.5,-3) {};
        \node[draw, circle, fill=black, label=above:{$ $}] (2) at (12,-2) {};
        \node[draw, circle, fill=red, label=below:{$ $}] (3) at (12.5,-3) {};
        \node[draw, circle, fill=white, label=below:{$ $}] (4) at (12,-4) {};
        \node[draw, circle, fill=red, label=below:{$ $}] (5) at (11,-4) {};
        \node[draw, circle, fill=black, label=below:{$ $}] (6) at (10.5,-3) {};
        \node[draw, circle, fill=white, label=above:{$ $}] (7) at (11,-2) {};
        
        % Aristas W_6
        \draw (1) -- (2);
        \draw (1) -- (3);
        \draw (1) -- (4);
        \draw (1) -- (5);
        \draw (1) -- (6);
        \draw (1) -- (7);
        \draw (2) -- (3);
        \draw (3) -- (4);
        \draw (4) -- (5);
        \draw (5) -- (6);
        \draw (6) -- (7);
        \draw (7) -- (2);
    \end{tikzpicture}   
    \caption{Comportamiento del VIH en $W_7$ con $R=n-2=5$.}
    \label{W7yR=5}
    \end{figure}
    As\'i, los casos $W_5$, $W_6$ y $W_7$ ilustran de manera clara c\'omo el centro cumple en cada paso la condici\'on de reinfecci\'on necesaria para evitar la extinci\'on, y justifican que, desde $W_7$ en adelante, el patr\'on observado se mantiene sin modificaciones estructurales. Esto completa la demostraci\'on de que para $R=n-2$ la infecci\'on no se extingue en $W_n$ para ning\'un $n\ge 4$, y por tanto establece rigurosamente que $VIH(W_n)\ge n-1$.
\end{proof}

\subsection{VIH en bipartitas completas $K_{n,m}$}
Sabemos que una gr\'afica $G$ se dice \textit{bipartita} si su conjunto de v\'ertices puede particionarse en dos subconjuntos disjuntos $A$ y $B$ tales que ninguna arista une v\'ertices del mismo subconjunto. Si cada v\'ertice de $A$ est\'a conectado con todos los v\'ertices de $V_2$, se obtiene la \textit{gr\'afica bipartita completa} $K_{n,m}$, donde $|A| = n$ y $|B| = m$. Por otra parte, las gr\'aficas bipartitas completas $K_{n,m}$ presentan un comportamiento m\'as complejo, dado que su estructura se divide en dos subconjuntos de v\'ertices donde las conexiones ocurren \'unicamente entre particiones distintas.

Este tipo de topolog\'ia genera una din\'amica de propagaci\'on asim\'etrica, en la que la infecci\'on puede alternar entre los dos conjuntos sin extenderse completamente a todos los v\'ertices al mismo tiempo. Sin embargo, debido a la simetr\'ia estructural de las gr\'aficas bipartitas completas, basta considerar el caso en el que $n \geq m$, ya que intercambiar las particiones no altera el comportamiento del modelo ni el valor del n\'umero $VIH(G)$.

\begin{figure}[h]
    \centering
    \begin{tabular}{|c|cccccccccc|}
        \hline
        \textbf{n/m} & \textbf{1} & \textbf{2} & \textbf{3} & \textbf{4} & \textbf{5} & \textbf{6} & \textbf{7} & \textbf{8} & \textbf{9} & \textbf{10} \\
        \hline
        \textbf{1} & 2 & . & . & . & . & . & . & . & . & . \\
        \textbf{2} & 2 & 3 & . & . & . & . & . & . & . & . \\
        \textbf{3} & 2 & 3 & 4 & . & . & . & . & . & . & . \\
        \textbf{4} & 3 & 3 & 4 & 5 & . & . & . & . & . & . \\
        \textbf{5} & 3 & 3 & 4 & 5 & 6 & . & . & . & . & . \\
        \textbf{6} & 4 & 4 & 4 & 5 & 6 & 7 & . & . & . & . \\
        \textbf{7} & 4 & 4 & 4 & 5 & 6 & 7 & 8 & . & . & . \\
        \textbf{8} & 5 & 5 & 5 & 5 & 6 & 7 & 8 & 9 & . & . \\
        \textbf{9} & 5 & 5 & 5 & 5 & 6 & 7 & 8 & 9 & 10 & . \\
        \textbf{10} & 6 & 6 & 6 & 6 & 6 & 7 & 8 & 9 & 10 & 11 \\
        \hline
    \end{tabular}
    \caption{Tabla $VIH(K_{n,m})$ con $n\geq m$}
    \label{tabla:VIH(K_n,m)n>=m}
\end{figure}
Los resultados experimentales mostrados en la Figura \ref{tabla:VIH(K_n,m)n>=m}, indican que el n\'umero $VIH(K_{n,m})$ depende tanto de la relaci\'on entre $n$ y $m$ como del tama\~no total de la gr\'afica, lo que sugiere la existencia de una relaci\'on no lineal entre ambos par\'ametros. Gracias esto, llegamos al siguiente teorema de $VIH$ para gr\'aficas bipartitas completas $K_{n,m}$ bajo la c\'ondicion $n\geq m$: 

\begin{theorem}\label{Teo:VIH(Kn,m)}
    $$
    VIH(K_{n,m}) = \left\lbrace
    \begin{array}{ll}
        m+1 \textup{ si } n \leq 2m\\
        \lceil \frac{n+1}{2} \rceil \textup{ si } n \geq 2m
    \end{array}
    \right.
    $$
\end{theorem}
\begin{proof}
    Sea $G = K_{n,m}$ una gr\'afica bipartita completa con conjuntos de v\'ertices $A$ y $B$, donde $|A| = n$ y $|B| = m$. Recordemos que una gr\'afica se dice bipartita si su conjunto de v\'ertices puede particionarse en dos subconjuntos disjuntos de manera que ninguna arista une v\'ertices del mismo subconjunto. 

    Comencemos con el caso en que $n \leq 2m$. En particular, cuando $n = m$ y el par\'ametro de reemplazo es $R = \Delta(K_{n,m}) = n = m$, la infecci\'on no se extingue. En efecto, si en el estado inicial todos los v\'ertices de $A$ est\'an infectados y los de $B$ sanos, es decir, $f_{0,R}(A)=1$ y $f_{0,R}(B)=0$, entonces en el siguiente instante cada v\'ertice de $B$ tiene $d_{0,1}(B)=n$ vecinos infectados, por lo que $f_{1,R}(A)=2$ y $f_{1,R}(B)=1$. En el tiempo $t=1$, cada v\'ertice en $A$ tiene $d_{1,1}(A)=m=R$, de manera que en el siguiente paso $f_{2,R}(A)=1$ y $f_{2,R}(B)=2$. Este proceso se repite indefinidamente, alternando los estados de infecci\'on y muerte entre ambos subconjuntos, como se muestra en la Figura \ref{KnmyR=m,n<=2m}. Sin embargo, si $R = m+1$, un v\'ertice muerto necesitar\'ia de $m+1$ vecinos infectados para re-infectarse, lo cual es imposible pues su grado es $m$. Por tanto, con $R=m+1$ la infecci\'on desaparece en un n\'umero finito de pasos, lo que justifica que $VIH(K_{n,m})=m+1$ para $n=m$.

    \begin{figure}[h]
    \centering
    \begin{tikzpicture}
    %tiempo 0
    % vertices
    \node[draw, ellipse, minimum height=60pt,minimum width=20pt, fill=red, label=above:{$A$}] (1) at (-1,0) {};
    \node[draw, ellipse, minimum height=60pt,minimum width=20pt, fill=white, label={[label distance=1pt]258:$t=0$}] (2) at (0.5,0) {};
    \node[draw, ellipse, minimum height=60pt,minimum width=20pt, fill=white, label=above:{$B$}] (a) at (0.5,0) {};

    % Aristas
    \draw (1) -- (2);

    %tiempo 1
    % vertices
    \node[draw, ellipse, minimum height=60pt,minimum width=20pt, fill=black, label=above:{$A$}] (3) at (2,0) {};
    \node[draw, ellipse, minimum height=60pt,minimum width=20pt, fill=red, label={[label distance=1pt]258:$t=1$}] (4) at (3.5,0) {};
    \node[draw, ellipse, minimum height=60pt,minimum width=20pt, fill=red, label=above:{$B$}] (b) at (3.5,0) {};

    % Aristas
    \draw (3) -- (4);
    
    %tiempo 2
    % vertices
    \node[draw, ellipse, minimum height=60pt,minimum width=20pt, fill=red, label=above:{$A$}] (5) at (5,0) {};
    \node[draw, ellipse, minimum height=60pt,minimum width=20pt, fill=black, label={[label distance=1pt]258:$t=2$}] (6) at (6.5,0) {};
    \node[draw, ellipse, minimum height=60pt,minimum width=20pt, fill=black, label=above:{$B$}] (c) at (6.5,0) {};

    % Aristas
    \draw (5) -- (6);

    %tiempo 3
    % vertices
    \node[draw, ellipse, minimum height=60pt,minimum width=20pt, fill=black, label=above:{$A$}] (7) at (8,0) {};
    \node[draw, ellipse, minimum height=60pt,minimum width=20pt, fill=red, label={[label distance=1pt]258:$t=3$}] (8) at (9.5,0) {};
    \node[draw, ellipse, minimum height=60pt,minimum width=20pt, fill=red, label=above:{$B$}] (d) at (9.5,0) {};

    % Aristas
    \draw (7) -- (8);
    \end{tikzpicture}
    \caption{Comportamiento del $VIH$ en $K_{n,m}$ con $n\leq2m$ y $R=m$}
    \label{KnmyR=m,n<=2m}
    \end{figure}

    Ahora consideremos el caso en que $n>m$ pero $n\leq 2m$. En este escenario, el comportamiento es an\'alogo. Si tomamos $R=\delta(K_{n,m})=m$ y la configuraci\'on inicial $f_{0,R}(A)=1$ y $f_{0,R}(B)=0$, los v\'ertices de $B$ tienen inicialmente $d_{0,1}(B)=n$, por lo que en el siguiente paso se obtiene $f_{1,R}(A)=2$ y $f_{1,R}(B)=1$. En el instante $t=1$, los v\'ertices de $A$ tienen $d_{1,1}(A)=m=R$, as\'i que $f_{2,R}(A)=1$ y $f_{2,R}(B)=2$. Este patr\'on de alternancia se mantiene indefinidamente. Si aumentamos el par\'ametro a $R=m+1$, los v\'ertices muertos del conjunto $A$ requerir\'ian $m+1$ vecinos infectados para reinfectarse, pero dado que su grado es $m$, la infecci\'on no puede sostenerse y el sistema se cura completamente tras un n\'umero finito de pasos. Por lo tanto, para todo $n\leq 2m$, el valor cr\'itico del par\'ametro es $R=m+1$, y en consecuencia $VIH(K_{n,m})=m+1$.

    Analicemos ahora el caso en que $n \geq 2m$. Si tomamos $n=2m$ y $R=m$, la din\'amica se comporta igual que en los casos anteriores: el sistema oscila sin curarse, alternando entre estados donde $A$ est\'a infectado y $B$ muerto, y viceversa. En cambio, si $R=m+1$, la enfermedad se extingue, pues los v\'ertices muertos carecen de suficientes vecinos infectados para reinfectarse. En este caso se cumple que $m+1 = \lceil \tfrac{n+1}{2}\rceil,$ ya que como $n=2m$, entonces: $$\lceil \frac{n+1}{2} \rceil=\lceil \frac{2m+1}{2} \rceil = \lceil \frac{2(m+\frac{1}{2})}{2} \rceil=\lceil m+\frac{1}{2} \rceil= m+\lceil \frac{1}{2} \rceil=m+1.$$

    Cuando $n > 2m$, el razonamiento requiere un an\'alisis m\'as detallado. Si $R = \lceil \tfrac{n+1}{2}\rceil - 1$, existir\'a siempre una configuraci\'on inicial capaz de mantener la infecci\'on indefinidamente. Para justificarlo, particionemos $A$ en dos subconjuntos $A_1$ y $A_2$, y $B$ en $B_1$ y $B_2$, tales que $|A_1| = \lceil \tfrac{n}{2}\rceil$, $|A_2| = n - \lceil \tfrac{n}{2}\rceil$, $|B_1| = \lceil \tfrac{m}{2}\rceil$ y $|B_2| = m - \lceil \tfrac{m}{2}\rceil$. Entonces se cumple que $|A_1|\geq |A_2|$ y $|B_1|\geq |B_2|$. 

    Si $n$ es impar, se tiene que $|A_1| > |A_2|$ y por tanto $|A_1| = \lceil \tfrac{n+1}{2}\rceil$. En cambio, si $n$ es par, se cumple $|A_1| = |A_2| = \lceil \tfrac{n+1}{2}\rceil - 1$. En ambos casos, si se toma una configuraci\'on inicial en la que los v\'ertices de $A_1$ est\'an infectados y los dem\'as sanos, se genera un ciclo de infecci\'on que alterna de la misma forma que en los casos anteriores, sin que la infecci\'on se extinga. Esto ocurre porque con $R = \lceil \tfrac{n+1}{2}\rceil - 1$, los v\'ertices muertos de $A$ tienen un n\'umero suficiente de vecinos infectados en $B$ para reinfectarse en el siguiente paso, manteniendo la din\'amica indefinidamente, esto se puede observar mejor en la Figura \ref{KnmyR=VIH-1,n>2m}. 

    \begin{figure}[h]
    \centering
    \begin{subfigure}[h]{0.85\linewidth}
    \begin{tikzpicture}
    %tiempo 0
    % vertices
    \node[draw, ellipse, minimum height=60pt,minimum width=1pt, fill=red] (1) at (-1,0) {$A_1$};
    \node[draw, ellipse, minimum height=60pt,minimum width=1pt, fill=white] (2) at (0.5,0) {$B_1$};
    \node[draw, ellipse, minimum height=60pt,minimum width=1pt, fill=white] (3) at (-1,-2.2) {$A_2$};
    \node[draw, ellipse, minimum height=60pt,minimum width=1pt, fill=white, label={[label distance=1pt]258:$t=0$}] (4) at (0.5,-2.2) {$B_2$};

    % Aristas
    \draw (1) -- (2);
    \draw (1) -- (4);
    \draw (3) -- (2);
    \draw (3) -- (4);

    %tiempo 1
    % vertices
    \node[draw, ellipse, minimum height=60pt,minimum width=1pt, fill=black] (1) at (2,0) {\textcolor{white}{$A_1$}};
    \node[draw, ellipse, minimum height=60pt,minimum width=1pt, fill=red] (2) at (3.5,0) {$B_1$};
    \node[draw, ellipse, minimum height=60pt,minimum width=1pt, fill=white] (3) at (2,-2.2) {$A_2$};
    \node[draw, ellipse, minimum height=60pt,minimum width=1pt, fill=red, label={[label distance=1pt]258:$t=1$}] (4) at (3.5,-2.2) {$B_2$};
    
    % Aristas
    \draw (1) -- (2);
    \draw (1) -- (4);
    \draw (3) -- (2);
    \draw (3) -- (4);

    %tiempo 2
    % vertices
    \node[draw, ellipse, minimum height=60pt,minimum width=1pt, fill=white] (1) at (5,0) {$A_1$};
    \node[draw, ellipse, minimum height=60pt,minimum width=1pt, fill=black] (2) at (6.5,0) {\textcolor{white}{$B_1$}};
    \node[draw, ellipse, minimum height=60pt,minimum width=1pt, fill=red] (3) at (5,-2.2) {$A_2$};
    \node[draw, ellipse, minimum height=60pt,minimum width=1pt, fill=black, label={[label distance=1pt]258:$t=2$}] (4) at (6.5,-2.2) {\textcolor{white}{$B_2$}};
    
    % Aristas
    \draw (1) -- (2);
    \draw (1) -- (4);
    \draw (3) -- (2);
    \draw (3) -- (4);

    %tiempo 3
    % vertices
    \node[draw, ellipse, minimum height=60pt,minimum width=1pt, fill=white] (1) at (8,0) {$A_1$};
    \node[draw, ellipse, minimum height=60pt,minimum width=1pt, fill=red] (2) at (9.5,0) {$B_1$};
    \node[draw, ellipse, minimum height=60pt,minimum width=1pt, fill=black] (3) at (8,-2.2) {\textcolor{white}{$A_2$}};
    \node[draw, ellipse, minimum height=60pt,minimum width=1pt, fill=red, label={[label distance=1pt]258:$t=3$}] (4) at (9.5,-2.2) {$B_2$};
    
    % Aristas
    \draw (1) -- (2);
    \draw (1) -- (4);
    \draw (3) -- (2);
    \draw (3) -- (4);
    \end{tikzpicture}
    \end{subfigure}
    \begin{subfigure}[h]{0.41\linewidth}
    \begin{tikzpicture}
    %tiempo 4
    % vertices
    \node[draw, ellipse, minimum height=60pt,minimum width=1pt, fill=red] (1) at (-1,0) {$A_1$};
    \node[draw, ellipse, minimum height=60pt,minimum width=1pt, fill=black] (2) at (0.5,0) {\textcolor{white}{$B_1$}};
    \node[draw, ellipse, minimum height=60pt,minimum width=1pt, fill=white] (3) at (-1,-2.2) {$A_2$};
    \node[draw, ellipse, minimum height=60pt,minimum width=1pt, fill=black, label={[label distance=1pt]258:$t=4$}] (4) at (0.5,-2.2) {\textcolor{white}{$B_2$}};

    % Aristas
    \draw (1) -- (2);
    \draw (1) -- (4);
    \draw (3) -- (2);
    \draw (3) -- (4);

    %tiempo 5
    % vertices
    \node[draw, ellipse, minimum height=60pt,minimum width=1pt, fill=black] (1) at (2,0) {\textcolor{white}{$A_1$}};
    \node[draw, ellipse, minimum height=60pt,minimum width=1pt, fill=red] (2) at (3.5,0) {$B_1$};
    \node[draw, ellipse, minimum height=60pt,minimum width=1pt, fill=white] (3) at (2,-2.2) {$A_2$};
    \node[draw, ellipse, minimum height=60pt,minimum width=1pt, fill=red, label={[label distance=1pt]258:$t=5$}] (4) at (3.5,-2.2) {$B_2$};
    
    % Aristas
    \draw (1) -- (2);
    \draw (1) -- (4);
    \draw (3) -- (2);
    \draw (3) -- (4);
    \end{tikzpicture}
    \end{subfigure}
    \caption{Comportamiento del $VIH$ en $K_{n,m}$ con $n>2m$ y $R=\lceil \frac{n+1}{2}\rceil-1$}
    \label{KnmyR=VIH-1,n>2m}
    \end{figure}

    Por lo tanto, el valor cr\'itico para el cual la infecci\'on deja de propagarse es $R = \lceil \tfrac{n+1}{2}\rceil$, ya que un incremento de una unidad en $R$ impide que los v\'ertices muertos puedan reinfectarse, garantizando que el sistema se cure completamente en un n\'umero finito de pasos.

    De esta forma, concluimos que para las gr\'aficas bipartitas completas $K_{n,m}$, el n\'umero VIH est\'a dado por: $$ VIH(K_{n,m}) = \left\lbrace \begin{array}{ll} m+1 & \textup{si } n \leq 2m,\\ \lceil \tfrac{n+1}{2} \rceil & \textup{si } n \geq 2m, \end{array} \right.$$ lo cual coincide con los resultados obtenidos mediante simulaci\'on computacional y con la din\'amica del modelo propuesta por Mukwembi.
\end{proof}
\subsection{VIH en ciclos $C_n$}
El siguiente caso corresponde a los ciclos $C_n$, cuya estructura cerrada introduce una simetr\'ia que modifica el comportamiento de la propagaci\'on. A diferencia de las trayectorias, en el ciclo no existen v\'ertices terminales, por lo que el proceso de curaci\'on depende de la interacci\'on circular entre los v\'ertices.

\begin{figure}[h!]
    \centering
    \begin{tabular}{|c|cccccccccc|}
        \hline
        \textbf{n} & \textbf{1} & \textbf{2} & \textbf{3} & \textbf{4} & \textbf{5} & \textbf{6} & \textbf{7} & \textbf{8} & \textbf{9} & \textbf{10}\\
        \hline
        \textbf{$VIH(C_n)$} & 1 & 2 & 2 & 3 & 2 & 3 & 2 & 3 & 2 & 3 \\
        \hline
    \end{tabular}
    \caption{Tabla $VIH(C_n)$}
    \label{tabla:VIH(C_n)}
\end{figure}
Los resultados experimentales los cuales se observan en la Figura \ref{tabla:VIH(C_n)}, mostraron que el n\'umero $VIH(C_n)$ requiere un valor ligeramente mayor que en el caso lineal, lo que refleja la influencia de la conectividad c\'iclica en la persistencia de la infecci\'on.
\begin{theorem}\label{Teo:VIH(Cn)}
    $$
    VIH(C_n) = \left\lbrace
    \begin{array}{ll}
        \lceil\frac{n+1}{2} \rceil \textup{ si } n < 5\\
        2 \textup{ si } n\geq5 \textup{ y } n \textup{ es impar }\\
        3 \textup{ si } n\geq5 \textup{ y } n \textup{ es par }
    \end{array}
    \right.
    $$
\end{theorem}
\begin{proof}
    Sea $C_n$ el ciclo de longitud $n\geq 1$. Recordemos que, de acuerdo con el Teorema \ref{Teo:Delta+1}, si $G$ es una gr\'afica que contiene al menos una arista, entonces se cumple que $2\leq VIH(G)\leq \Delta(G)+1$. Como en el caso de los ciclos se tiene $\Delta(C_n)=2$, obtenemos inmediatamente que $2\leq VIH(C_n)\leq 3$. Por tanto, basta con analizar cu\'ando el par\'ametro $R=2$ es suficiente para extinguir la infecci\'on y cu\'ando no lo es, de modo que determinemos el valor exacto de $VIH(C_n)$. Antes de ello, consideremos los casos triviales: si $n=1$, el ciclo $C_1$ consiste de un solo v\'ertice, por lo que se sigue de los resultados previos que $VIH(C_1)=1$. Para $n=2$, la gr\'afica $C_2$ tiene la misma estructura que $P_2$, $S_1$, $K_2$ y $W_2$, las cuales ya se demostr\'o que satisfacen $VIH(G)=2$. En el caso $n=3$, observamos que $C_3$ es isomorfa a $K_3$, de donde se deduce que $VIH(C_3)=2$. Finalmente, $C_4$ tiene la misma estructura que $K_{2,2}$, para la cual se prob\'o que $VIH(K_{2,2})=3$, y por tanto $VIH(C_4)=3$. Con esto, los casos $n<5$ quedan establecidos y coinciden con el teorema propuesta.

    Supongamos ahora que $n\geq5$. Como ya sabemos que $2\leq VIH(C_n)\leq 3$, debemos mostrar que cuando $n$ es impar, el valor $R=2$ es suficiente para que la infecci\'on desaparezca para cualquier condici\'on inicial, mientras que cuando $n$ es par, el valor $R=2$ no es suficiente y, por tanto, se requiere $R=3$.

    Para el caso en que $n$ es par, mostraremos que $R=2$ no es suficiente para que la enfermedad se cure. Sea $V(C_n)={v_1,v_2,\dots,v_n}$ con las aristas $(v_i,v_{i+1})$ y $(v_n,v_1)$, y consideremos la coloraci\'on inicial alternante dada por $f_{0,2}(v_i)=1$ si $i$ es impar y $f_{0,2}(v_i)=0$ si $i$ es par. Esta configuraci\'on es propia en el sentido de que ning\'un par de v\'ertices adyacentes est\'a simult\'aneamente infectado.

    Analicemos la evoluci\'on del proceso con $R=2$. En el instante $t=1$, cada v\'ertice inicialmente infectado pasa al estado muerto por la regla que establece que si $f_{t,R}(v)=1$, entonces $f_{t+1,R}(v)=2$. Por otro lado, cada v\'ertice inicialmente sano tiene al menos un vecino infectado (de hecho, ambos), y por tanto se infecta en el instante siguiente. As\'i, la configuraci\'on $f_{1,2}$ es exactamente la complementaria de $f_{0,2}$, es decir, los v\'ertices infectados en el instante $1$ son aquellos que estaban sanos al inicio. 
    
    En el instante $t=2$, ocurre nuevamente la misma transici\'on: los infectados del tiempo anterior pasan a muertos y los muertos cuyos dos vecinos estaban infectados se vuelven a infectar, lo que reconstituye la configuraci\'on inicial. Por tanto, el sistema entra en un ciclo de periodo $2$ en el que las clases alternan indefinidamente, manteniendo siempre v\'ertices infectados, como se puede observar en la Figura \ref{C6yR=2}. En consecuencia, la infecci\'on nunca desaparece para esta condici\'on inicial, lo que implica que $R=2$ no es suficiente en ciclos pares. Dado que la cota superior garantiza que $VIH(C_n)\leq3$, concluimos que para todo $n\geq5$ par se cumple $VIH(C_n)=3$.

    \begin{figure}[h]
    \centering
    \begin{tikzpicture}
        %t=0
        % vertices W_6 
        \node[label={[label distance=29pt]270:$t=0$}] (1) at (1,0) {};
        \node[draw, circle, fill=white, label=above:{$ $}] (2) at (1.5,1) {};
        \node[draw, circle, fill=red, label=below:{$ $}] (3) at (2,0) {};
        \node[draw, circle, fill=white, label=below:{$ $}] (4) at (1.5,-1) {};
        \node[draw, circle, fill=red, label=below:{$ $}] (5) at (0.5,-1) {};
        \node[draw, circle, fill=white, label=below:{$ $}] (6) at (0,0) {};
        \node[draw, circle, fill=red, label=above:{$ $}] (7) at (0.5,1) {};
        
        % Aristas W_6
        \draw (2) -- (3);
        \draw (3) -- (4);
        \draw (4) -- (5);
        \draw (5) -- (6);
        \draw (6) -- (7);
        \draw (7) -- (2);

        %t=1
        % vertices W_6
        \node[label={[label distance=29pt]270:$t=1$}] (1) at (4,0) {};
        \node[draw, circle, fill=red, label=above:{$ $}] (2) at (4.5,1) {};
        \node[draw, circle, fill=black, label=below:{$ $}] (3) at (5,0) {};
        \node[draw, circle, fill=red, label=below:{$ $}] (4) at (4.5,-1) {};
        \node[draw, circle, fill=black, label=below:{$ $}] (5) at (3.5,-1) {};
        \node[draw, circle, fill=red, label=below:{$ $}] (6) at (3,0) {};
        \node[draw, circle, fill=black, label=above:{$ $}] (7) at (3.5,1) {};
        
        % Aristas W_6
        \draw (2) -- (3);
        \draw (3) -- (4);
        \draw (4) -- (5);
        \draw (5) -- (6);
        \draw (6) -- (7);
        \draw (7) -- (2);

        %t=2
        % vertices W_6
        \node[label={[label distance=29pt]270:$t=2$}] (1) at (7,0) {};
        \node[draw, circle, fill=black, label=above:{$ $}] (2) at (7.5,1) {};
        \node[draw, circle, fill=red, label=below:{$ $}] (3) at (8,0) {};
        \node[draw, circle, fill=black, label=below:{$ $}] (4) at (7.5,-1) {};
        \node[draw, circle, fill=red, label=below:{$ $}] (5) at (6.5,-1) {};
        \node[draw, circle, fill=black, label=below:{$ $}] (6) at (6,0) {};
        \node[draw, circle, fill=red, label=above:{$ $}] (7) at (6.5,1) {};
        
        % Aristas W_6
        \draw (2) -- (3);
        \draw (3) -- (4);
        \draw (4) -- (5);
        \draw (5) -- (6);
        \draw (6) -- (7);
        \draw (7) -- (2);

        %t=3
        % vertices W_6
        \node[label={[label distance=29pt]270:$t=3$}] (1) at (10,0) {};
        \node[draw, circle, fill=red, label=above:{$ $}] (2) at (10.5,1) {};
        \node[draw, circle, fill=black, label=below:{$ $}] (3) at (11,0) {};
        \node[draw, circle, fill=red, label=below:{$ $}] (4) at (10.5,-1) {};
        \node[draw, circle, fill=black, label=below:{$ $}] (5) at (9.5,-1) {};
        \node[draw, circle, fill=red, label=below:{$ $}] (6) at (9,0) {};
        \node[draw, circle, fill=black, label=above:{$ $}] (7) at (9.5,1) {};
        
        % Aristas W_6
        \draw (2) -- (3);
        \draw (3) -- (4);
        \draw (4) -- (5);
        \draw (5) -- (6);
        \draw (6) -- (7);
        \draw (7) -- (2);

        %t=4
        % vertices W_6
        \node[label={[label distance=29pt]270:$t=4$}] (1) at (13,0) {};
        \node[draw, circle, fill=black, label=above:{$ $}] (2) at (13.5,1) {};
        \node[draw, circle, fill=red, label=below:{$ $}] (3) at (14,0) {};
        \node[draw, circle, fill=black, label=below:{$ $}] (4) at (13.5,-1) {};
        \node[draw, circle, fill=red, label=below:{$ $}] (5) at (12.5,-1) {};
        \node[draw, circle, fill=black, label=below:{$ $}] (6) at (12,0) {};
        \node[draw, circle, fill=red, label=above:{$ $}] (7) at (12.5,1) {};
        
        % Aristas W_6
        \draw (2) -- (3);
        \draw (3) -- (4);
        \draw (4) -- (5);
        \draw (5) -- (6);
        \draw (6) -- (7);
        \draw (7) -- (2);

    \end{tikzpicture}   
    \caption{Comportamiento del VIH en $C_6$ con $R=2$.}
    \label{C6yR=2}
    \end{figure}

    Pasemos ahora al caso en que $n$ es impar. Queremos mostrar que $R=2$ es suficiente, es decir, que para cualquier condici\'on inicial $f_{0,2}$, el proceso termina en un estado libre de infectados tras un n\'umero finito de pasos. 
    
    Supongamos, con el fin de llegar a una contradicci\'on, que existe una condici\'on inicial en la que la infecci\'on no se extingue. Como el conjunto de configuraciones posibles es finito y la din\'amica del modelo es determinista, debe existir un ciclo peri\'odico de configuraciones que se repite indefinidamente. Si la infecci\'on no desaparece, entonces en cada paso de dicho ciclo hay al menos un v\'ertice infectado. Sin embargo, por la regla de evoluci\'on, todo v\'ertice infectado en un instante se convierte en muerto en el siguiente, de modo que para que la infecci\'on persista, debe haber muertos que sean re-infectados. Con $R=2$, esto s\'olo ocurre si ambos vecinos de un muerto estaban infectados en el paso anterior. Esto implica que la infecci\'on se propaga alternadamente entre dos subconjuntos disjuntos de v\'ertices, de manera que en cada paso los infectados ocupan una de las dos clases y en el siguiente la otra. 
    
    Este comportamiento peri\'odico de alternancia s\'olo es posible si el ciclo admite una 2-coloraci'on propia, es decir, si se puede dividir $V(C_n)$ en dos subconjuntos $A$ y $B$ tales que cada v\'ertice tenga a sus dos vecinos en el subconjunto opuesto. No obstante, tal partici\'on s\'olo es posible cuando $n$ es par, y por tanto cuando $n$ es impar esta estructura alternante no puede existir.

    Por lo tanto, para $n$ impar y $R=2$, la infecci\'on necesariamente desaparece tras un n\'umero finito de pasos, sin importar la condici\'on inicial. Esto demuestra que $R=2$ es suficiente para todo $n\geq5$ impar, y por tanto $VIH(C_n)=2$ en este caso.
\end{proof}

\subsection{VIH en rejillas $R_{n,m}$}
Finalmente, se analiz\'o el comportamiento del modelo sobre las gr\'aficas tipo rejilla $R_{n,m}$. En este caso, la estructura bidimensional y regular de la gr\'afica produce un patr\'on de propagaci\'on m\'as controlado y localizado, donde los efectos de frontera influyen significativamente en la din\'amica del sistema.

Debido a la simetr\'ia de la estructura de la rejilla, basta considerar el caso en que $n \geq m$, pues intercambiar las dimensiones no altera la topolog\'ia de la gr\'afica ni modifica el resultado de la simulaci\'on.

%Acompletar tabla Rn,m con n>=m
\begin{figure}[h]
    \centering
    \begin{tabular}{|c|cccccccccc|}
        \hline
        \textbf{n/m} & \textbf{1} & \textbf{2} & \textbf{3} & \textbf{4} & \textbf{5} & \textbf{6} & \textbf{7} & \textbf{8} & \textbf{9} & \textbf{10} \\
        \hline
        \textbf{1} & 1 & . & . & . & . & . & . & . & . & . \\
        \textbf{2} & 2 & 3 & . & . & . & . & . & . & . & . \\
        \textbf{3} & 2 & 3 & 4 & . & . & . & . & . & . & . \\
        \textbf{4} & 2 & 3 & 4 & 4 & . & . & . & . & . & . \\
        \textbf{5} & 2 & 3 & 4 & 4 & 4 & . & . & . & . & . \\
        \textbf{6} & 2 & 3 & 4 & . & . & . & . & . & . & . \\
        \textbf{7} & 2 & 3 & 4 & . & . & . & . & . & . & . \\
        \textbf{8} & 2 & 3 & . & . & . & . & . & . & . & . \\
        \textbf{9} & 2 & 3 & . & . & . & . & . & . & . & . \\
        \textbf{10} & 2 & 3 & . & . & . & . & . & . & . & . \\
        \hline
    \end{tabular}
    \caption{Tabla $VIH(R_{n,m})$ con $n\geq m$}
    \label{tabla:VIH(R_n,m)n>=m}
\end{figure}
Las simulaciones vistas en la Figura \ref{tabla:VIH(R_n,m)n>=m} mostraron que el valor de $R$ necesario para alcanzar la curaci\'on total crece lentamente con el tama\~no de la rejilla, lo que refleja la estabilidad de la estructura y su capacidad para contener la propagaci\'on de la infecci\'on incluso en configuraciones de mayor tama\~no. Gracias esto, llegamos al siguiente teorema de $VIH$ para rejillas $R_{n,m}$ bajo la c\'ondicion $n\geq m$:

\begin{theorem}\label{Con:VIH(Rn,m)}
    $$
    VIH(R_{n,m}) = \left\lbrace
    \begin{array}{ll}
        1 \textup{ si } n=m=1\\
        2 \textup{ si } m=1 \textup{ y } n > m\\
        3 \textup{ si } m=2 \textup{ y } n \geq m\\
        4 \textup{ si } m\geq 3 \textup{ y } n \geq m
    \end{array}
    \right.
    $$
\end{theorem}

\begin{proof}
    A continuaci\'on demostraremos el valor del par\'ametro $VIH(R_{n,m})$ para toda rejilla finita $R_{n,m}$ con $n\geq m$.

    Si $m=1$ la rejilla se identifica con la misma estructura que una trayectoria $P_n$, como esta se demostr\'o anteriormente en el Teorema \ref{Teo:VIH(Pn)}, entonces se obtiene que $VIH(P_1)=1$ y $VIH(P_n)=2$ para $n\geq 2$, en otras palabras $VIH(R_{1,1})=1$ y $VIH(R_{n,1})=2$ para $n\geq 2>m$.

    Ahora, si $m=2$ estas rejillas son tipo escalera, adem\'as siempre contienen un ciclo $C_4$ como subgr\'afica y por el Teorema \ref{Teo:VIH(Cn)}, sabemos que $VIH(C_4)=3$, por lo que $VIH(R_{n,2})\geq VIH(C_4)=3$, esto gracias a los lemas \ref{Lem:BorrarVertices} y \ref{Lem:BorrarAristas} antes mencionados de monoton\'ia. Adem\'as, $\Delta(R_{n,2})=3$, y por el Teorema \ref{Teo:Delta+1} se tiene $VIH(R_{n,2})\leq 4$, por ende al final tenemos que $3\leq VIH(R_{n,2})\leq 4$. Basta mostrar que $R=3$ es suficiente para extinguir la infecci\'on.

    Fijamos el par\'ametro de reemplazo $R=3$. La rejilla $R_{n,2}$ tiene v\'ertices se dividen en; los cuatro v\'ertices de las esquinas con $\deg(v)=2$, a este conjunto de v\'ertices lo denotaremos como $V_a$, y los v\'ertices del borde con $\deg(v)=3$, los cuales denotaremos como $V_b$. La regla crucial de Mukwembi establece que una c\'elula muerta $f_{i,3}(v)=2$ solo se re-infecta si tiene $d_{i,1}(v) \ge R$ vecinos infectados.

    Consideremos los v\'ertices de grado $\deg(v)=2$, que son las esquinas. Para estos, el n\'umero m\'aximo de vecinos infectados es $d_{i,1}(v_a) \le 2$, donde $v_a \in V_a$. Dado que $2 < R=3$, la condici\'on de re-infecci\'on $d_{i,1}(v) \ge 3$ nunca se satisface. En consecuencia, cualquier c\'elula muerta en una esquina ser\'a reemplazada inevitablemente por una c\'elula sana $f_{i+1,3}(v_a)=0$. Esto establece el concepto que llamaremos como frontera curativa: la infecci\'on no puede ser sostenida en la periferia, y la transici\'on completa de una esquina es $f_{i-1,3}(v_a)=1$, $f_{i,3}(v_a)=2$ y $f_{i+1,3}(v_a)=0$.

    Los \'unicos v\'ertices capaces de sostener el ciclo de re-infecci\'on son los v\'ertices del borde $V_b$ con $\deg(v_b)=3$. Para que un v\'ertice $v_b \in V_b$ se re-infecte, es imprescindible que sus tres vecinos est\'en infectados de forma simult\'anea $d_{i,1}(v_b)=3$. Sin embargo, la existencia de la frontera curativa, definida por los v\'ertices de las esquinas $V_a$, asegura que la infecci\'on en la periferia es transitoria e insostenible. Esta insostenibilidad se debe a que las c\'elulas muertas en $V_a$ se curan obligatoriamente, eliminando permanentemente los focos infecciosos en el exterior. De este modo, la re-infecci\'on de los v\'ertices $V_b$, aun cuando es posible al coincidir $R=3$, no logra generar nuevos infectados a una tasa suficiente para contrarrestar la p\'erdida neta de c\'elulas infecciosas viables causada por el ciclo de curaci\'on definitivo en $V_a$. En esencia, las c\'elulas $V_b$ no pueden mantener la condici\'on $d_{i,1}(v_b)=3$ el tiempo suficiente porque sus vecinos en la esquina est\'an programados para curarse. Por lo tanto, el sistema es incapaz de entrar en un ciclo persistente, y la infecci\'on debe extinguirse en un n\'umero finito de pasos. Se concluye que $R=3$ es suficiente para que la enfermedad se extinga, probando as\'i que $VIH(R_{n,2})=3$.

    Para el caso en donde $m\ge 3$, mostraremos primero que $VIH(R_{3,3})=4$ y posteriormente generalizaremos el resultado para toda rejilla $R_{n,m}$ con $n\ge m\ge 3$ usando los lemas de monoton\'ia. Primero analicemos la rejilla $R_{3,3}$, esta gr\'afica contiene un ciclo $C_4$ como subgr\'afica y, por el Teorema \ref{Teo:VIH(Cn)}, sabemos que $VIH(C_4)=3$, por lo que $VIH(R_{3,3})\ge VIH(C_4)=3$. Adem\'as, $\Delta(R_{3,3})=4$, y por la cota general $VIH(G)\le \Delta(G)+1$ se tiene que $VIH(R_{3,3})\le 5$. De esta forma obtenemos $3\le VIH(R_{3,3})\le 5$.

    Para concluir que $VIH(R_{3,3})=4$, basta probar dos afirmaciones; la primera es probar que con $R=3$ la infecci\'on no necesariamente se extingue, es decir, existe una condici\'on inicial para la cual la enfermedad no se cura; y la segunda es probar que con $R=4$ la infecci\'on siempre se extingue.

    Veamos el caso cuando $R=3$ para una rejilla $R_{3,3}$ en donde denotaremos a cada v\'ertice como $v_{n,m}$, es decir, el v\'ertice de la fila $n$ y columna $m$, para una mejor visualizaci\'on de la din\'amica de propagaci\'on que explicaremos a continuaci\'on vea la Figura \ref{R3,3yR=3}.
    
    En la condici\'on inicial, $t=0$, la configuraci\'on consta de tres v\'ertices infectados $f_{0,3}(v_{1,1})=f_{0,3}(v_{1,2})=f_{0,3}(v_{2,2})=1$, mientras los dem\'as v\'ertices est\'an sanos. En la primera transici\'on $t=1$, por la regla dos, los tres v\'ertices infectados en $t=0$ transitan inmediatamente a muertos en $t=1$, es decir, $f_{1,3}(v_{1,1})=f_{1,3}(v_{1,2})=f_{1,3}(v_{2,2})=2$ . Simult\'aneamente, los v\'ertices sanos en $t=0$ que son adyacentes a al menos un infectado $d_{0,1}(v)\geq 1$ se infectaran, en caso contrario permanecer\'an sanos. Por lo tanto, en $t=1$ la configuraci\'on de los dem\'as v\'ertices es $f_{1,3}(v_{1,3})=f_{1,3}(v_{2,1})=f_{1,3}(v_{2,3})=f_{1,3}(v_{3,2})=1$ y $f_{1,3}(v_{3,1})=f_{1,3}(v_{3,3})=0$

    En $t=2$, los v\'ertices que estaban infectados en $t=1$ pasaran a muertos $f_{2,3}(v_{1,3})=f_{2,3}(v_{2,1})=f_{2,3}(v_{2,3})=f_{2,3}(v_{3,2})=2$, los que estaban susceptibles como cada uno tenia $d_{1,1}(v)\geq 1$ entonces se infectar\'an $f_{2,3}(v_{3,1})=f_{2,3}(v_{3,3})=1$, mientras que los que estaban muertos el v\'ertice $v_{2,2}$ es el \'unico que tiene $d_{1,1}(v_{2,2})\geq R=3$ por ende se volver\'a a infectar en el siguiente tiempo $f_{2,3}(v_{2,2})=1$, el resto de v\'ertices muertos pasaran a sanos $f_{2,3}(v_{1,1})=f_{2,3}(v_{1,2})=0$. 

    Pasemos al tiempo $t=3$, nuevamente los que estaban infectados en el tiempo anterior pasaran a muertos en este tiempo $f_{3,3}(v_{2,2})=f_{3,3}(v_{3,1})=f_{3,3}(v_{3,3})=2$, los que estaban muertos el v\'ertice $v_{3,2}$ es el \'unico que tiene $d_{2,1}(v_{3,2})\geq R=3$ por ende se volver\'a a infectar en el siguiente tiempo $f_{3,3}(v_{3,2})=1$, el resto de v\'ertices muertos pasaran a sanos $f_{3,3}(v_{1,3})=f_{3,3}(v_{2,1})=f_{3,3}(v_{2,3})=0$, mientras que los que estaban susceptibles el v\'ertice $v_{1,2}$ tiene $d_{2,1}(v_{1,2})\geq 1$ entonces se infectar\'a $f_{3,3}(v_{1,2})=1$ el el otro permanecer\'a sano $f_{3,3}(v_{1,1})=0$.

    En el instante $t=4$, los v\'ertices que estaban muertos, ninguno tiene $d_{3,1}(v)\geq R=3$, por lo tanto sanaran $f_{4,3}(v_{2,2})=f_{4,3}(v_{3,1})=f_{4,3}(v_{3,3})=0$. Los v\'ertices que estaban susceptibles los \'unicos que tienen $d_{3,1}(v)\geq 1$ son $v_{1,1}$ y $v_{1,3}$, por ende se infectaran $f_{4,3}(v_{1,1})=f_{4,3}(v_{1,3})=1$ mientras que los otros permanecer\'an sanos $f_{4,3}(v_{2,1})=f_{4,3}(v_{2,3})=0$, finalmente los que estaban infectados pasaran a muertos $f_{4,3}(v_{1,2})=f_{4,3}(v_{3,2})=2$.

    Veamos ahora el instante $t=5$, los infectados pasaran a muertos $f_{5,3}(v_{1,1})=f_{5,3}(v_{1,3})=2$, los que estaban muertos ninguno tiene $d_{4,1}(v)\geq R=3$, por lo tanto sanaran $f_{5,3}(v_{1,2})=f_{5,3}(v_{3,2})=0$. Por otro lado, los susceptibles los \'unicos que tienen $d_{4,1}(v)\geq 1$ son $v_{2,1}$ y $v_{2,3}$, por ende se infectaran $f_{5,3}(v_{2,1})=f_{5,3}(v_{2,3})=1$, el resto permanecer\'an sanos $f_{5,3}(v_{2,2})=f_{5,3}(v_{3,1})=f_{5,3}(v_{3,3})=0$.

    En $t=6$, los que estaban muertos como no tienen $d_{5,1}(v)\geq R=3$, entonces sanaran $f_{6,3}(v_{1,1})=f_{6,3}(v_{1,3})=0$, los que estaban infectados forzosamente morir\'an $f_{6,3}(v_{2,1})=f_{6,3}(v_{2,3})=2$. Mientras que,en los susceptibles los v\'ertices $v_{2,2}$, $v_{3,1}$ y $v_{3,3}$ tienen $d_{5,1}(v)\geq 1$ por lo tato se infectaran $f_{6,3}(v_{2,2})=f_{6,3}(v_{3,1})=f_{6,3}(v_{3,3})=1$, y los restantes permanecer\'an sanos $f_{6,3}(v_{1,2})=f_{6,3}(v_{3,2})=0$.

    Finalmente en $t=7$, los infectados pasaran a muertos $f_{7,3}(v_{2,2})=f_{7,3}(v_{3,1})=f_{7,3}(v_{3,3})=2$, en el conjunto de v\'ertices que estaban muertos ninguno cumple con la condici\'on $d_{6,1}(v)\geq R=3$ ppr lo tanto nacer\'an $f_{7,3}(v_{2,1})=f_{7,3}(v_{2,3})=0$. Por \'ultimo en los susceptibles los v\'ertices que tienen  $d_{6,1}(v)\geq 1$ se infectaran, en caso contrario se quedaran como sanos, $f_{7,3}(v_{1,2})=f_{7,3}(v_{3,2})=1$ y $f_{7,3}(v_{1,1})=f_{7,3}(v_{1,3})=0$. Podemos observar que en este tiempo la configuraci\'on es la misma que en el tiempo $t=4$, por lo tanto podemos deducir que  en le siguiente tiempo $t=8$, va a ser igual que el tiempo $t=5$ y as\'i sucesivamente, formando as\'i un ciclo repetitivo de periodo tres. De este modo la infecci\'on nunca se extingue en $R_{3,3}$ cuando $R = 3$.

    \begin{figure}[h]
    \centering
    \begin{tikzpicture}
        %tiempo 0 n=fila , m =columna
        % vertices
        \node[draw, circle, fill=red, label=above:{$ $}] (a) at (0,0) {};
        \node[draw, circle, fill=red, label=below:{$t=0$}] (b) at (1,0) {};
        \node[draw, circle, fill=white, label=above:{$ $}] (c) at (2,0) {};
        \node[draw, circle, fill=white, label=above:{$ $}] (d) at (0,1) {};
        \node[draw, circle, fill=red, label=above:{$ $}] (e) at (1,1) {};
        \node[draw, circle, fill=white, label=above:{$ $}] (f) at (2,1) {};
        \node[draw, circle, fill=white, label=above:{$ $}] (g) at (0,2) {};
        \node[draw, circle, fill=white, label=above:{$ $}] (h) at (1,2) {};
        \node[draw, circle, fill=white, label=above:{$ $}] (i) at (2,2) {};

        % Aristas
        \draw (a) -- (b);
        \draw (b) -- (c);
        \draw (d) -- (e);
        \draw (e) -- (f);
        \draw (g) -- (h);
        \draw (h) -- (i);
        \draw (a) -- (d);
        \draw (d) -- (g);
        \draw (b) -- (e);
        \draw (e) -- (h);
        \draw (c) -- (f);
        \draw (f) -- (i);
        
        %tiempo 1 n=fila , m =columna
        % vertices
        \node[draw, circle, fill=black, label=above:{$ $}] (a) at (3,0) {};
        \node[draw, circle, fill=black, label=below:{$t=1$}] (b) at (4,0) {};
        \node[draw, circle, fill=red, label=above:{$ $}] (c) at (5,0) {};
        \node[draw, circle, fill=red, label=above:{$ $}] (d) at (3,1) {};
        \node[draw, circle, fill=black, label=above:{$ $}] (e) at (4,1) {};
        \node[draw, circle, fill=red, label=above:{$ $}] (f) at (5,1) {};
        \node[draw, circle, fill=white, label=above:{$ $}] (g) at (3,2) {};
        \node[draw, circle, fill=red, label=above:{$ $}] (h) at (4,2) {};
        \node[draw, circle, fill=white, label=above:{$ $}] (i) at (5,2) {};

        % Aristas
        \draw (a) -- (b);
        \draw (b) -- (c);
        \draw (d) -- (e);
        \draw (e) -- (f);
        \draw (g) -- (h);
        \draw (h) -- (i);
        \draw (a) -- (d);
        \draw (d) -- (g);
        \draw (b) -- (e);
        \draw (e) -- (h);
        \draw (c) -- (f);
        \draw (f) -- (i);

        %tiempo 2 n=fila , m =columna
        % vertices
        \node[draw, circle, fill=white, label=above:{$ $}] (a) at (6,0) {};
        \node[draw, circle, fill=white, label=below:{$t=2$}] (b) at (7,0) {};
        \node[draw, circle, fill=black, label=above:{$ $}] (c) at (8,0) {};
        \node[draw, circle, fill=black, label=above:{$ $}] (d) at (6,1) {};
        \node[draw, circle, fill=red, label=above:{$ $}] (e) at (7,1) {};
        \node[draw, circle, fill=black, label=above:{$ $}] (f) at (8,1) {};
        \node[draw, circle, fill=red, label=above:{$ $}] (g) at (6,2) {};
        \node[draw, circle, fill=black, label=above:{$ $}] (h) at (7,2) {};
        \node[draw, circle, fill=red, label=above:{$ $}] (i) at (8,2) {};

        % Aristas
        \draw (a) -- (b);
        \draw (b) -- (c);
        \draw (d) -- (e);
        \draw (e) -- (f);
        \draw (g) -- (h);
        \draw (h) -- (i);
        \draw (a) -- (d);
        \draw (d) -- (g);
        \draw (b) -- (e);
        \draw (e) -- (h);
        \draw (c) -- (f);
        \draw (f) -- (i);

        %tiempo 3 n=fila , m =columna
        % vertices
        \node[draw, circle, fill=white, label=above:{$ $}] (a) at (9,0) {};
        \node[draw, circle, fill=red, label=below:{$t=3$}] (b) at (10,0) {};
        \node[draw, circle, fill=white, label=above:{$ $}] (c) at (11,0) {};
        \node[draw, circle, fill=white, label=above:{$ $}] (d) at (9,1) {};
        \node[draw, circle, fill=black, label=above:{$ $}] (e) at (10,1) {};
        \node[draw, circle, fill=white, label=above:{$ $}] (f) at (11,1) {};
        \node[draw, circle, fill=black, label=above:{$ $}] (g) at (9,2) {};
        \node[draw, circle, fill=red, label=above:{$ $}] (h) at (10,2) {};
        \node[draw, circle, fill=black, label=above:{$ $}] (i) at (11,2) {};

        % Aristas
        \draw (a) -- (b);
        \draw (b) -- (c);
        \draw (d) -- (e);
        \draw (e) -- (f);
        \draw (g) -- (h);
        \draw (h) -- (i);
        \draw (a) -- (d);
        \draw (d) -- (g);
        \draw (b) -- (e);
        \draw (e) -- (h);
        \draw (c) -- (f);
        \draw (f) -- (i);

        %tiempo 4 n=fila , m =columna
        % vertices
        \node[draw, circle, fill=red, label=above:{$ $}] (a) at (0,-3) {};
        \node[draw, circle, fill=black, label=below:{$t=4$}] (b) at (1,-3) {};
        \node[draw, circle, fill=red, label=above:{$ $}] (c) at (2,-3) {};
        \node[draw, circle, fill=white, label=above:{$ $}] (d) at (0,-2) {};
        \node[draw, circle, fill=white, label=above:{$ $}] (e) at (1,-2) {};
        \node[draw, circle, fill=white, label=above:{$ $}] (f) at (2,-2) {};
        \node[draw, circle, fill=white, label=above:{$ $}] (g) at (0,-1) {};
        \node[draw, circle, fill=black, label=above:{$ $}] (h) at (1,-1) {};
        \node[draw, circle, fill=white, label=above:{$ $}] (i) at (2,-1) {};

        % Aristas
        \draw (a) -- (b);
        \draw (b) -- (c);
        \draw (d) -- (e);
        \draw (e) -- (f);
        \draw (g) -- (h);
        \draw (h) -- (i);
        \draw (a) -- (d);
        \draw (d) -- (g);
        \draw (b) -- (e);
        \draw (e) -- (h);
        \draw (c) -- (f);
        \draw (f) -- (i);
        
        %tiempo 5 n=fila , m =columna
        % vertices
        \node[draw, circle, fill=black, label=above:{$ $}] (a) at (3,-3) {};
        \node[draw, circle, fill=white, label=below:{$t=5$}] (b) at (4,-3) {};
        \node[draw, circle, fill=black, label=above:{$ $}] (c) at (5,-3) {};
        \node[draw, circle, fill=red, label=above:{$ $}] (d) at (3,-2) {};
        \node[draw, circle, fill=white, label=above:{$ $}] (e) at (4,-2) {};
        \node[draw, circle, fill=red, label=above:{$ $}] (f) at (5,-2) {};
        \node[draw, circle, fill=white, label=above:{$ $}] (g) at (3,-1) {};
        \node[draw, circle, fill=white, label=above:{$ $}] (h) at (4,-1) {};
        \node[draw, circle, fill=white, label=above:{$ $}] (i) at (5,-1) {};

        % Aristas
        \draw (a) -- (b);
        \draw (b) -- (c);
        \draw (d) -- (e);
        \draw (e) -- (f);
        \draw (g) -- (h);
        \draw (h) -- (i);
        \draw (a) -- (d);
        \draw (d) -- (g);
        \draw (b) -- (e);
        \draw (e) -- (h);
        \draw (c) -- (f);
        \draw (f) -- (i);

        %tiempo 6 n=fila , m =columna
        % vertices
        \node[draw, circle, fill=white, label=above:{$ $}] (a) at (6,-3) {};
        \node[draw, circle, fill=white, label=below:{$t=6$}] (b) at (7,-3) {};
        \node[draw, circle, fill=white, label=above:{$ $}] (c) at (8,-3) {};
        \node[draw, circle, fill=black, label=above:{$ $}] (d) at (6,-2) {};
        \node[draw, circle, fill=red, label=above:{$ $}] (e) at (7,-2) {};
        \node[draw, circle, fill=black, label=above:{$ $}] (f) at (8,-2) {};
        \node[draw, circle, fill=red, label=above:{$ $}] (g) at (6,-1) {};
        \node[draw, circle, fill=white, label=above:{$ $}] (h) at (7,-1) {};
        \node[draw, circle, fill=red, label=above:{$ $}] (i) at (8,-1) {};

        % Aristas
        \draw (a) -- (b);
        \draw (b) -- (c);
        \draw (d) -- (e);
        \draw (e) -- (f);
        \draw (g) -- (h);
        \draw (h) -- (i);
        \draw (a) -- (d);
        \draw (d) -- (g);
        \draw (b) -- (e);
        \draw (e) -- (h);
        \draw (c) -- (f);
        \draw (f) -- (i);

        %tiempo 7 n=fila , m =columna
        % vertices
        \node[draw, circle, fill=white, label=above:{$ $}] (a) at (9,-3) {};
        \node[draw, circle, fill=red, label=below:{$t=7$}] (b) at (10,-3) {};
        \node[draw, circle, fill=white, label=above:{$ $}] (c) at (11,-3) {};
        \node[draw, circle, fill=white, label=above:{$ $}] (d) at (9,-2) {};
        \node[draw, circle, fill=black, label=above:{$ $}] (e) at (10,-2) {};
        \node[draw, circle, fill=white, label=above:{$ $}] (f) at (11,-2) {};
        \node[draw, circle, fill=black, label=above:{$ $}] (g) at (9,-1) {};
        \node[draw, circle, fill=red, label=above:{$ $}] (h) at (10,-1) {};
        \node[draw, circle, fill=black, label=above:{$ $}] (i) at (11,-1) {};

        % Aristas
        \draw (a) -- (b);
        \draw (b) -- (c);
        \draw (d) -- (e);
        \draw (e) -- (f);
        \draw (g) -- (h);
        \draw (h) -- (i);
        \draw (a) -- (d);
        \draw (d) -- (g);
        \draw (b) -- (e);
        \draw (e) -- (h);
        \draw (c) -- (f);
        \draw (f) -- (i);
    \end{tikzpicture}
    \caption{Comportamiento del VIH en $R_{3,3}$ con $R=3$.}
    \label{R3,3yR=3}
    \end{figure}

    Ahora, fijemos el par\'ametro $R=4$, primero analicemos la clasificaci\'on de los nueve v\'ertices; la gr\'afica $R_{3,3}$ consta de cuatro esquinas que nuevamente denotaremos como $V_a$ con $deg(v_a)=2$ donde $v_a \in V_a$, cuartro v\'ertices de borde $V_b$ con $\deg(v_b)=3$ donde $v_b \in V_b$. Para cualquiera de estos ochos v\'ertices, el n\'umero m\'aximo de vecinos infectados es $d_{i,1}(v) \le 3$. Dado que $3 < R=4$, la condici\'on de re-infecci\'on $d_{i,1}(v) \ge 4$ es imposible de satisfacer. De esta forma, cualquier c\'elula muerta en la periferia siempre\ se\ cura.

    Sin embargo la gr\'afica $R_{3,3}$ consta de un v\'ertice central que denotaremos como $v_c$ y adem\'as $\deg(v_c)=4$. Este es el \'unico v\'ertice que, si est\'a muerto, tiene la capacidad de re-infectarse, requiriendo que sus cuatro vecinos est\'en infectados, es decir, $d_{i,1}(v_c)=4$.

    Entonces, la demostraci\'on de la extinci\'on se basa en el aislamiento de $v_c$. La periferia de la rejilla, que incluye los v\'ertices $V_a$ y $V_b$ con $\deg(v) \le 3$, funcionan como una frontera curativa permanente ya que el par\'ametro $R=4$ excede su grado. Para que la infecci\'on se mantenga en el tiempo, el v\'ertice central $v_c$ debe participar en el ciclo re-infecci\'on. Sin embargo, este ciclo es insostenible, la re-infecci\'on de $v_c$ exige la m\'axima saturaci\'on infecciosa de los cuatro vecinos infectados, y todos sus vecinos, al ser v\'ertices de borde $V_b$, est\'an sujetos a la curaci\'on obligatoria de la frontera. La carga infecciosa que muere no puede ser compensada por las re-infecciones, ya que el costo de regenerar el \'unico v\'ertice posible $v_c$ es demasiado alto y sus fuentes infecciosas son transitorias. Consecuentemente, en un n\'umero finito de pasos, la infecci\'on se extingue en $R_{3,3}$, probando $VIH(R_{3,3}) = 4$.

    Finalmente, la generalizaci\'on de $R_{n,m}$ con $n \ge m \ge 3$ se sigue por los lemas \ref{Lem:BorrarVertices} y \ref{Lem:BorrarAristas}. La existencia de $R_{3,3}$ como subgr\'afica implica $VIH(R_{n,m}) \ge VIH(R_{3,3}) = 4$. Puesto que $\Delta(R_{n,m})=4$, la cota superior general $VIH(R_{n,m}) \le 5$ acota el resultado, confirmando que $VIH(R_{n,m})=4$ para este rango de rejillas.
    
\end{proof}

A partir de los teoremas anteriores se deduce que el n\'umero $VIH(G)$ presenta un comportamiento dependiente de la estructura de la gr\'afica. En particular, las familias con mayor grado m\'aximo tienden a requerir valores de $R$ m\'as altos para garantizar la curaci\'on total. Esto se debe a que una mayor conectividad entre v\'ertices aumenta la probabilidad de re-infecci\'on, lo que exige un umbral de resistencia m\'as elevado.

Por el contrario, en gr\'aficas con estructura lineal o menos densa, como las trayectorias $P_n$ o los ciclos $C_n$, se observa que valores m\'as peque\~nos de $R$ son suficientes para lograr la curaci\'on completa. En cambio, en las familias de gr\'aficas completas y ruedas, donde todos los v\'ertices est\'an altamente conectados, los valores de $VIH(G)$ alcanzan sus niveles m\'as altos, lo que confirma que la densidad estructural incrementa la persistencia de la infecci\'on antes de que el sistema pueda curarse. Por otra parte, en las gr\'aficas bipartitas y en las rejillas se observan valores intermedios, reflejando un equilibrio entre la regularidad local y la dispersi\'on global de la estructura.

\chapter{Conclusiones y trabajo futuro}
%En este cap\'itulo se incluir\'an las conclusiones y posibles l\'ineas de trabajo que podr\'ian surgir a partir de este proyecto

En este trabajo se estudi\'o un modelo matem\'atico para la propagaci\'on del Virus de Inmunodeficiencia Humana en los ganglios linf\'aticos, utilizando herramientas de la teor\'ia de gr\'aficas. Dado que la geometr\'ia real de los ganglios linf\'aticos no se conoce con precisi\'on, el enfoque adoptado permite modelar la situaci\'on de manera general, considerando distintas posibles estructuras y analizando c\'omo estas influyen en la din\'amica de la infecci\'on. A trav\'es del modelo de Mukwembi se logr\'o describir, mediante reglas sencillas, la evoluci\'on temporal de c\'elulas sanas, infectadas y muertas dentro de una gr\'afica.

Uno de los resultados centrales del trabajo es el n\'umero de $VIH(G)$ de una gr\'afica, que representa el valor m\'inimo del par\'ametro de reemplazo $R$ para el cual la infecci\'on desaparece independientemente de la condici\'on inicial. Los resultados obtenidos muestran que la estructura de la gr\'afica tiene un papel fundamental en la persistencia o eliminaci\'on del virus. En particular, se observa que cuando el reemplazo de c\'elulas muertas favorece completamente a las c\'elulas sanas, es decir, cuando la probabilidad de que una c\'elula muerta se reemplace por una c\'elula sana es igual a uno, el modelo predice que la enfermedad termina por desaparecer en un n\'umero finito de pasos.

El an\'alisis detallado de distintas familias de gr\'aficas permiti\'o identificar comportamientos espec\'ificos y establecer resultados exactos en varios casos. Las simulaciones computacionales fueron una herramienta clave para explorar el modelo, formular teoremas y apoyar las demostraciones te\'oricas. En conjunto, el trabajo muestra que incluso con un modelo relativamente simple es posible obtener informaci\'on relevante sobre la din\'amica de la infecci\'on y su relaci\'on con la estructura del entorno en el que se propaga.

A partir de este proyecto se abren diversas l\'ineas de trabajo futuro. Una posibilidad natural es estudiar el modelo en gr\'aficas aleatorias, con el objetivo de analizar el comportamiento promedio del n\'umero VIH y comparar estos resultados con los obtenidos para familias deterministas. Otra l\'inea de inter\'es es profundizar en el caso de los \'arboles, donde la ausencia de ciclos podr\'ia influir de manera significativa en la propagaci\'on y eliminaci\'on de la infecci\'on.

Asimismo, resulta relevante estudiar no solo si la infecci\'on desaparece, sino tambi\'en cu\'anto tiempo tarda en hacerlo. En este sentido, ser\'ia interesante analizar el tiempo m\'inimo $t$ a partir del cual todos los v\'ertices se encuentran sanos, y c\'omo este tiempo depende tanto de la gr\'afica como del valor del par\'ametro $R$. Este tipo de estudio permitir\'ia una comprensi\'on m\'as fina de la din\'amica del modelo y complementar\'ia los resultados obtenidos en esta tesis.

En conclusi\'on, el trabajo presentado ofrece una base s\'olida para el estudio del modelo de Mukwembi desde un punto de vista matem\'atico y computacional, y deja abiertas varias direcciones de investigaci\'on que pueden ser desarrolladas en trabajos posteriores.



%%%%%%%%%%%%%%%%%%%%%%%%%%%%%%%%%%%%%%%%%%%%%%%%%%%%%%%
% La bibliograf\'ia debe aparecer de manera homologada
%%%%%%%%%%%%%%%%%%%%%%%%%%%%%%%%%%%%%%%%%%%%%%%%%%%%%%%

\begin{thebibliography}{0}
	%\bibitem{}
	\bibitem{Mukwembi} Mukwembi, S. (2008). A note on the effects of replenishment of depleted cells on HIV infection dynamics: A graph-theoretic approach. Physica A: Statistical Mechanics and its Applications, 387(5-6), 1200-1204.
    \bibitem{Harary} Harary, F. (1969). Graph Theory. Reading, Mass. : Perseus Books.
    \bibitem{Matthes} Matthes, E. (2024). Curso intensivo de Python: Introducci\'on pr\'actica a la programaci\'on basada en proyectos (3.a ed.). Madrid: Anaya Multimedia.
\end{thebibliography}

\end{document}


